% Non-infectious diseases — Biology Upper Sixth — Cours signé OnBuch+
% Official MINESEC syllabus. Compiler avec : tectonic BIO06-non-infectious-diseases.tex
\newcommand{\DOCMATIERE}{Biology}
\newcommand{\DOCNIVEAU}{Upper Sixth}
\newcommand{\DOCMODULE}{Upper Sixth — Biology}
\newcommand{\DOCLECON}{6}
\newcommand{\DOCTITRE}{Non-infectious diseases}
% OnBuch+ — préambule « cours signés » ENRICHI (compilé avec tectonic / XeLaTeX)
% Système visuel « Pop » d'OnBuch+ : fond crème (jamais blanc), bordures encre
% épaisses, ombres « dures » décalées, titres Archivo Black en capitales.
% Version MATHÉMATIQUES : graphes (pgfplots/tikz), boîtes pédagogiques
% (définition, propriété, méthode, attention, le savais-tu, exercice résolu,
% exercice, TP, bilan), page de garde et signature OnBuch+ ancrée en bas.
%
% Macros attendues dans main.tex AVANT \input{preamble} :
%   \DOCMATIERE  \DOCNIVEAU  \DOCMODULE  \DOCLECON (numéro)  \DOCTITRE
\documentclass[11pt]{article}

\usepackage{fontspec}
\usepackage[english]{babel}
\usepackage[a4paper,margin=2cm,top=3.4cm,bottom=2.8cm,headheight=1.4cm,footskip=1.3cm]{geometry}
\usepackage{amsmath,amssymb,amsfonts}
\usepackage{mathtools} % \xmapsto, \coloneqq... souvent utilisés par les modèles
\usepackage{cancel} % simplifications barrées (bilans de forces)
% \abs{z} (module, valeur absolue) et \norm{u} (norme), utilisés par les modèles
\providecommand{\abs}{}\let\abs\relax\DeclarePairedDelimiter{\abs}{\lvert}{\rvert}
\providecommand{\norm}{}\let\norm\relax\DeclarePairedDelimiter{\norm}{\lVert}{\rVert}
\usepackage[table]{xcolor} % [table] : fournit aussi \rowcolors (alternance de lignes)
\usepackage{pagecolor}
\usepackage{tikz}
\usetikzlibrary{arrows.meta,positioning,calc,shapes.geometric,shapes.misc,shapes.arrows,decorations.pathmorphing,decorations.pathreplacing,decorations.markings,patterns,shadows,mindmap,trees,fit,backgrounds,matrix,intersections,angles,quotes}
\usepackage{pgfplots}
\usepackage[version=4]{mhchem} % équations nucléaires \ce{^{235}_{92}U}
\usepackage[european,siunitx]{circuitikz} % schémas électriques
\pgfplotsset{compat=1.18}
\usepgfplotslibrary{fillbetween,groupplots} % aires entre courbes (calcul intégral)
\usepackage{enumitem}
\usepackage{fancyhdr}
% [most] charge toutes les bibliothèques tcolorbox (listings, minted, raster,
% documentation...) et gonfle inutilement la mémoire TeX sur un document long
% (~300 boîtes) : on ne charge que ce qui est réellement utilisé.
\usepackage[skins,breakable]{tcolorbox}
\usepackage{graphicx}
\usepackage{colortbl}
\usepackage{booktabs}
\usepackage{tabularx}
\usepackage{diagbox} % case d'en-tête diagonale des tableaux à double entrée
% Colonnes extensibles centrée / gauche / droite (C, L, R), souvent utilisées par les modèles
\newcolumntype{C}{>{\centering\arraybackslash}X}
\newcolumntype{L}{>{\raggedright\arraybackslash}X}
\newcolumntype{R}{>{\raggedleft\arraybackslash}X}
\usepackage{array}
\usepackage{multicol}
\usepackage{siunitx}
% parse-numbers=false : accepte aussi \SI{2,5 \times 10^{-3}}{...}
\sisetup{locale=UK,output-decimal-marker={.},group-minimum-digits=4,parse-numbers=false}
\DeclareSIUnit\ohm{\ensuremath{\symup{\Omega}}} % U+2126 absent de la police
\DeclareSIUnit\division{div} % réglages d'oscilloscope (ms/div, V/div)
\DeclareSIUnit\tr{tr} % tours (vitesse de rotation, ex. \SI{1500}{\tr\per\minute})
\usepackage{qrcode}
% \degree : commande fréquemment utilisée par les modèles pour un angle ou une
% température en dehors de \SI{}{} (non fournie par les paquets chargés).
\providecommand{\degree}{^{\circ}}
% \qed (fin de démonstration), utilisé par les modèles sans amsthm
\providecommand{\qed}{\hfill$\square$}
% \barre : double barre des valeurs interdites dans un tableau de variations (tkz-tab)
\providecommand{\barre}{\ensuremath{\|}}
% environnement proof (amsthm non chargé) utilisé par les modèles
\ifdefined\proof\else\newenvironment{proof}[1][Proof]{\par\noindent\textit{#1.}\ }{\qed\par}\fi
\usepackage{lastpage}
\usepackage{hyperref}
\hypersetup{hidelinks}

% ============================================================
% Palette « Pop » OnBuch+ — mêmes hex que la classe P (plus_ui.dart)
% ============================================================
\definecolor{poporange}{HTML}{F59321}
\definecolor{popdark}{HTML}{E07A0C}
\definecolor{popcream}{HTML}{FFF6E8}
\definecolor{popink}{HTML}{1C1714}
\definecolor{poppurple}{HTML}{7A5AE0}
\definecolor{popgreen}{HTML}{2E9E6A}
\definecolor{poppink}{HTML}{E0457B}
\definecolor{popblue}{HTML}{2F7DE1}
\definecolor{popgold}{HTML}{C98A12}
\definecolor{popmuted}{HTML}{8A7F76}
\definecolor{popline}{HTML}{EADFD2}
\definecolor{popgray}{HTML}{8A7F76}
\colorlet{popgrey}{popgray}
% Alias tolérants (le modèle nomme parfois les couleurs différemment)
\colorlet{popwhite}{white}
\colorlet{popblack}{popink}
\colorlet{popred}{poppink}
\colorlet{popyellow}{popgold}
% Teintes claires pour fonds de tableaux / zones
\colorlet{popblueL}{popblue!10!white}
\colorlet{poporangeL}{poporange!14!white}
\colorlet{popgreenL}{popgreen!12!white}
\colorlet{poppurpleL}{poppurple!12!white}
\colorlet{poppinkL}{poppink!12!white}
\colorlet{popgoldL}{popgold!14!white}

\pagecolor{popcream}

% ============================================================
% Typographie
% ============================================================
\defaultfontfeatures{Ligatures=TeX}
\setmainfont{PlusJakartaSans-Medium.ttf}[Path=../../../fonts/,BoldFont=PlusJakartaSans-Bold.ttf]
% Maths en sans-serif (Fira Math) avec chiffres et lettres droites de Plus
% Jakarta, pour que les formules se fondent dans le texte.
\usepackage[math-style=ISO,bold-style=ISO]{unicode-math}
\setmathfont{FiraMath-Regular.otf}[Path=../../../fonts/]
\setmathfont{PlusJakartaSans-Medium.ttf}[Path=../../../fonts/,range={up/num,up/{Latin,latin},\mathup}]
% (icomma retiré : décimales avec point en anglais)
% Caractères mathématiques Unicode absents des polices de texte (Plus Jakarta,
% Archivo Black) : rendus par leur équivalent mathématique, en texte comme en
% formule — sinon ils s'affichent en carré vide (ex. « Arithmétique dans ℤ »).
\usepackage{newunicodechar}
\newunicodechar{ℝ}{\ensuremath{\mathbb{R}}}
\newunicodechar{ℤ}{\ensuremath{\mathbb{Z}}}
\newunicodechar{ℂ}{\ensuremath{\mathbb{C}}}
\newunicodechar{ℕ}{\ensuremath{\mathbb{N}}}
\newunicodechar{ℚ}{\ensuremath{\mathbb{Q}}}
\newunicodechar{𝔻}{\ensuremath{\mathbb{D}}}
\newunicodechar{α}{\ensuremath{\alpha}}
\newunicodechar{β}{\ensuremath{\beta}}
\newunicodechar{γ}{\ensuremath{\gamma}}
\newunicodechar{δ}{\ensuremath{\delta}}
\newunicodechar{ε}{\ensuremath{\varepsilon}}
\newunicodechar{θ}{\ensuremath{\theta}}
\newunicodechar{λ}{\ensuremath{\lambda}}
\newunicodechar{μ}{\ensuremath{\mu}}
\newunicodechar{σ}{\ensuremath{\sigma}}
\newunicodechar{τ}{\ensuremath{\tau}}
\newunicodechar{φ}{\ensuremath{\varphi}}
\newunicodechar{ψ}{\ensuremath{\psi}}
\newunicodechar{ω}{\ensuremath{\omega}}
\newunicodechar{Σ}{\ensuremath{\Sigma}}
% majuscules produites par \MakeUppercase dans les titres (α-aminés -> Α)
\newunicodechar{Α}{\ensuremath{\alpha}}
\newunicodechar{Β}{\ensuremath{\beta}}
\newunicodechar{Γ}{\ensuremath{\Gamma}}
\newunicodechar{Δ}{\ensuremath{\Delta}}
\newunicodechar{Ε}{\ensuremath{\varepsilon}}
\newunicodechar{Θ}{\ensuremath{\Theta}}
\newunicodechar{Λ}{\ensuremath{\Lambda}}
\newunicodechar{Μ}{\ensuremath{\mu}}
\newunicodechar{Τ}{\ensuremath{\tau}}
\newunicodechar{Φ}{\ensuremath{\Phi}}
\newunicodechar{Ψ}{\ensuremath{\Psi}}
\newunicodechar{Ω}{\ensuremath{\Omega}}
\newunicodechar{↦}{\ensuremath{\mapsto}}
\newunicodechar{∈}{\ensuremath{\in}}
\newunicodechar{∉}{\ensuremath{\notin}}
\newunicodechar{∧}{\ensuremath{\wedge}}
\newunicodechar{⇔}{\ensuremath{\Leftrightarrow}}
\newunicodechar{⟺}{\ensuremath{\Longleftrightarrow}}
\newunicodechar{⇒}{\ensuremath{\Rightarrow}}
\newunicodechar{⟹}{\ensuremath{\Longrightarrow}}
\newunicodechar{∘}{\ensuremath{\circ}}
\newunicodechar{∪}{\ensuremath{\cup}}
\newunicodechar{∩}{\ensuremath{\cap}}
\newunicodechar{≡}{\ensuremath{\equiv}}
\newunicodechar{⊂}{\ensuremath{\subset}}
\newunicodechar{⊥}{\ensuremath{\perp}}
\newunicodechar{∃}{\ensuremath{\exists}}
\newunicodechar{∀}{\ensuremath{\forall}}
\newunicodechar{∛}{\ensuremath{\sqrt[3]{\,}}}
\newunicodechar{∜}{\ensuremath{\sqrt[4]{\,}}}
\newunicodechar{⃗}{\ensuremath{{}^{\to}}}
\newunicodechar{ⁿ}{\ifmmode^{n}\else\textsuperscript{\ensuremath{n}}\fi}
\newunicodechar{ˣ}{\ifmmode^{x}\else\textsuperscript{\ensuremath{x}}\fi}
\newunicodechar{ᵇ}{\ifmmode^{b}\else\textsuperscript{\ensuremath{b}}\fi}
\newunicodechar{ʳ}{\ifmmode^{r}\else\textsuperscript{\ensuremath{r}}\fi}
\newunicodechar{ᵘ}{\ifmmode^{u}\else\textsuperscript{\ensuremath{u}}\fi}
\newunicodechar{ʸ}{\ifmmode^{y}\else\textsuperscript{\ensuremath{y}}\fi}
\newunicodechar{ᵃ}{\ifmmode^{a}\else\textsuperscript{\ensuremath{a}}\fi}
\newunicodechar{ᵉ}{\ifmmode^{e}\else\textsuperscript{\ensuremath{e}}\fi}
\newunicodechar{ᵏ}{\ifmmode^{k}\else\textsuperscript{\ensuremath{k}}\fi}
\newunicodechar{ᵖ}{\ifmmode^{p}\else\textsuperscript{\ensuremath{p}}\fi}
\newunicodechar{ᴺ}{\ifmmode^{N}\else\textsuperscript{\ensuremath{N}}\fi}
\newunicodechar{ᶜ}{\ifmmode^{c}\else\textsuperscript{\ensuremath{c}}\fi}
\newunicodechar{ʰ}{\ifmmode^{h}\else\textsuperscript{\ensuremath{h}}\fi}
\newunicodechar{ᵠ}{\ifmmode^{q}\else\textsuperscript{\ensuremath{q}}\fi}
\newunicodechar{ₙ}{\ifmmode_{n}\else\textsubscript{\ensuremath{n}}\fi}
\newunicodechar{ₐ}{\ifmmode_{a}\else\textsubscript{\ensuremath{a}}\fi}
\newunicodechar{ᵢ}{\ifmmode_{i}\else\textsubscript{\ensuremath{i}}\fi}
\newunicodechar{ᵣ}{\ifmmode_{r}\else\textsubscript{\ensuremath{r}}\fi}
\newunicodechar{ₖ}{\ifmmode_{k}\else\textsubscript{\ensuremath{k}}\fi}
\newunicodechar{ᵦ}{\ifmmode_{\beta}\else\textsubscript{\ensuremath{\beta}}\fi}
\newunicodechar{ₓ}{\ifmmode_{x}\else\textsubscript{\ensuremath{x}}\fi}
\newfontfamily\popdisplay{ArchivoBlack-Regular.ttf}[Path=../../../fonts/]
\newfontfamily\poplogo{SpaceGrotesk-Bold.ttf}[Path=../../../fonts/]
\newfontfamily\popsemi{PlusJakartaSans-SemiBold.ttf}[Path=../../../fonts/]
\newfontfamily\popxbold{PlusJakartaSans-ExtraBold.ttf}[Path=../../../fonts/]
\newfontfamily\popxboldls{PlusJakartaSans-ExtraBold.ttf}[Path=../../../fonts/,LetterSpace=3.0]

\setlist[enumerate]{leftmargin=1.7em,itemsep=5pt,topsep=4pt}
\setlist[itemize]{leftmargin=1.7em,itemsep=4pt,topsep=4pt,label={\color{poporange}\textbullet}}
\setlength{\parindent}{0pt}
\setlength{\parskip}{5pt}
\renewcommand{\arraystretch}{1.25}

% pgfplots : style Pop par défaut
\pgfplotsset{
  every axis/.append style={axis line style={popink,line width=1pt},tick style={popink},
    grid style={popline,line width=0.5pt},label style={font=\small},tick label style={font=\footnotesize}},
  popcurve/.style={poporange,line width=1.8pt,no markers,smooth},
}
% Même style disponible dans un \draw TikZ simple (hors pgfplots)
\tikzset{popcurve/.style={poporange,line width=1.8pt}}
\tikzset{popgrid/.style={popline,very thin}}
\colorlet{popcurve}{poporange} % le nom du style est parfois utilisé comme couleur

% ============================================================
% Pill / squiggle
% ============================================================
\newcommand{\poppill}[3]{%
  \tikz[baseline=-0.55ex]\node[draw=popink,line width=1.2pt,rounded corners=2.6pt,%
    fill=#2,inner xsep=6.5pt,inner ysep=3pt]{%
    \popxboldls\fontsize{7}{7}\selectfont\color{#3}\MakeUppercase{#1}};%
}
\newcommand{\popsquiggle}[1]{%
  \tikz[baseline=-0.4ex]{
    \draw[#1,line width=2.6pt,line cap=round]
      plot [smooth] coordinates {(0,0.09) (0.34,0.01) (0.68,0.21) (1.02,0.01) (1.36,0.10)};
  }%
}
% Mot-clé surligné (marqueur orange sous le texte)
\newcommand{\cle}[1]{{\popxbold\color{popdark}#1}}

% ============================================================
% En-tête / pied
% ============================================================
\pagestyle{fancy}
\fancyhf{}
\renewcommand{\headrulewidth}{0pt}
\renewcommand{\footrulewidth}{0pt}
\lhead{\raisebox{-0.15ex}{\poplogo\fontsize{13}{13}\selectfont\color{popink}OnBuch\color{poporange}+}}
\rhead{\poppill{\DOCMATIERE\ \textbullet\ \DOCNIVEAU\ \textbullet\ Chapter \DOCLECON}{white}{popink}}
\lfoot{\popxbold\fontsize{7.6}{7.6}\selectfont\color{popmuted}\MakeUppercase{Signed course OnBuch+}\ \textbullet\ {\popsemi\DOCTITRE}}
\rfoot{\tikz[baseline=-0.6ex]\node[rounded corners=8.5pt,fill=popink,minimum height=17pt,minimum width=17pt,inner xsep=5pt,inner sep=0]{\popxbold\fontsize{8}{8}\selectfont\color{popcream}\thepage\,/\,\pageref*{LastPage}};}

% ============================================================
% Titres
% ============================================================
\newcommand{\chaptitre}[2]{%
  \begin{tcolorbox}[enhanced,colback=poporange,colframe=popink,boxrule=2.6pt,arc=11pt,
      drop shadow=popink,left=15pt,right=15pt,top=12pt,bottom=11pt,before skip=3pt,after skip=18pt]
    \poppill{#2}{popink}{popcream}\par\vspace{9pt}
    {\raggedright\hyphenpenalty=10000\exhyphenpenalty=10000\popdisplay\fontsize{22}{24}\selectfont\color{popcream}\MakeUppercase{#1}\par}
  \end{tcolorbox}
}
% Section numérotée : pastille orange avec le numéro + titre Archivo
\newcounter{popsec}
\newcommand{\coursec}[1]{%
  \stepcounter{popsec}\setcounter{popsub}{0}%
  \par\vspace{16pt}\noindent
  \begin{minipage}{\linewidth}
  \tikz[baseline=-0.7ex]\node[draw=popink,line width=1.6pt,fill=poporange,rounded corners=4pt,minimum size=22pt,inner sep=0]{\popdisplay\fontsize{12}{12}\selectfont\color{popcream}\thepopsec};%
  \hspace{8pt}{\popdisplay\fontsize{15}{17}\selectfont\color{popink}\MakeUppercase{#1}}\par
  \vspace{4pt}\noindent{\color{poporange}\rule{36pt}{3.6pt}}
  \end{minipage}\par\nopagebreak\vspace{8pt}
}
\newcounter{popsub}
\newcommand{\courssub}[1]{%
  \stepcounter{popsub}%
  \par\vspace{9pt}\noindent{\popxbold\fontsize{11.5}{13}\selectfont\color{popink}{\color{poporange}\thepopsec.\thepopsub}\hspace{6pt}#1}\par\nopagebreak\vspace{4pt}
}

% ============================================================
% Boîtes pédagogiques — même recette : fond blanc, bordure épaisse colorée,
% ombre dure encre, badge plein. Titre optionnel [#1].
% ============================================================
\tcbset{popbase/.style={breakable,enhanced,colback=white,boxrule=2.3pt,arc=9pt,
  shadow={3pt}{-3pt}{0pt}{popink},left=14pt,right=14pt,top=11pt,bottom=11pt,before skip=11pt,after skip=15pt,
  fontupper=\popsemi\fontsize{10.6}{14.5}\selectfont\color{popink},
  fontlower=\popsemi\fontsize{10.6}{14.5}\selectfont\color{popink}}}
% Coche dessinée (le glyphe ✓ n'existe pas dans les polices du document)
\newcommand{\popcheck}[1]{\tikz[baseline=-0.5ex]\draw[#1,line width=1.6pt,line cap=round,line join=round](-0.13,0.0)--(-0.04,-0.09)--(0.13,0.1);}
\providecommand{\coche}{\popcheck{popgreen}}
\providecommand{\resultat}[1]{\par\smallskip\noindent\fcolorbox{popgreen}{popgreenL}{\parbox{\dimexpr\linewidth-2\fboxsep-2\fboxrule\relax}{\textbf{Result.} #1}}\par\smallskip}
\newcommand{\popboxhead}[3]{\poppill{#1}{#2}{white}\ifx\relax#3\relax\else\hspace{6pt}{\popxbold\fontsize{10.6}{12}\selectfont\color{popink}#3}\fi\par\vspace{7pt}}

\newtcolorbox{aretenir}[1][]{popbase,colframe=popgreen,before upper={\popboxhead{Key points}{popgreen}{#1}}}
\newtcolorbox{exemplebox}[1][]{popbase,colframe=poporange,before upper={\popboxhead{Exemple}{poporange}{#1}}}
\newtcolorbox{definition}[1][]{popbase,colframe=popblue,before upper={\popboxhead{Definition}{popblue}{#1}}}
\newtcolorbox{propriete}[1][]{popbase,colframe=poppurple,before upper={\popboxhead{Property}{poppurple}{#1}}}
\newtcolorbox{methode}[1][]{popbase,colframe=popgold,before upper={\popboxhead{Method}{popgold}{#1}}}
\newtcolorbox{attention}[1][]{popbase,colframe=poppink,before upper={\popboxhead{Warning}{poppink}{#1}}}
\newtcolorbox{experience}[1][]{popbase,colframe=popblue,colback=popblueL,before upper={\popboxhead{Experiment}{popblue}{#1}}}
% « Le savais-tu ? » : carte violette pleine (contexte, culture, Cameroun)
\newtcolorbox{savaistu}[1][]{popbase,colback=poppurple,colframe=popink,
  fontupper=\popsemi\fontsize{10.4}{14.2}\selectfont\color{white},
  before upper={\poppill{Did you know?}{white}{poppurple}\ifx\relax#1\relax\else\hspace{6pt}{\popxbold\color{white}#1}\fi\par\vspace{7pt}}}
% Exercice résolu : énoncé en haut, \tcblower puis la solution sur fond crème
\newcounter{popexo}\newcounter{popexr}
\newtcolorbox{exoresolu}[1][]{popbase,colframe=poporange,colbacklower=popcream,
  segmentation style={popink,line width=1.2pt,dashed},
  before upper={\stepcounter{popexr}\popboxhead{Worked example \thepopexr}{poporange}{#1}},
  before lower={\poppill{Solution}{popink}{popcream}\par\vspace{6pt}}}
% Exercice (non résolu en place) — la correction est dans la partie Corrigés
\newtcolorbox{exercice}[1][]{popbase,colframe=popink,
  before upper={\stepcounter{popexo}\popboxhead{Exercise \thepopexo}{popink}{#1}}}
% Corrigé
\newtcolorbox{corrige}[1][]{popbase,colframe=popgreen,colback=popgreenL,
  before upper={\popboxhead{Solution}{popgreen}{#1}}}
% Objectifs / prérequis / bilan
\newtcolorbox{objectifs}{popbase,colframe=popink,colback=poporangeL,
  before upper={\popboxhead{Chapter objectives}{popink}{}}}
\newtcolorbox{prerequis}{popbase,colframe=popink,colback=popblueL,
  before upper={\popboxhead{Prerequisites}{popblue}{}}}
\newtcolorbox{bilan}{popbase,colframe=popink,colback=white,boxrule=2.8pt,
  before upper={\popboxhead{Fiche bilan}{poporange}{L'essentiel à connaître pour l'examen}}}
% Figure encadrée avec légende
\newtcolorbox{popfigure}[1][]{enhanced,breakable=false,colback=white,colframe=popline,boxrule=1.4pt,arc=8pt,
  left=10pt,right=10pt,top=10pt,bottom=8pt,before skip=10pt,after skip=12pt,halign=center,
  lower separated=false,halign lower=center,
  fontlower=\popsemi\fontsize{9}{11.5}\selectfont\color{popmuted},
  #1}
\newcommand{\legende}[1]{\par\vspace{6pt}{\popsemi\fontsize{9}{11.5}\selectfont\color{popmuted}#1}}

% ============================================================
% Signature OnBuch+
% ============================================================
\newcommand{\poplogoinline}[1]{{\poplogo\fontsize{#1}{#1}\selectfont\color{popink}OnBuch\color{poporange}+}}
\newcommand{\signatureonbuch}{%
  \begin{tcolorbox}[enhanced,colback=popink,colframe=popink,boxrule=2.2pt,arc=10pt,
      drop shadow=poporange,left=14pt,right=14pt,top=10pt,bottom=10pt,before skip=0pt,after skip=0pt]
    \tikz[baseline=-0.7ex]\node[circle,fill=popgreen,draw=popcream,line width=1.3pt,minimum size=22pt,inner sep=0]{\popcheck{white}};%
    \hspace{9pt}\begin{minipage}[c]{0.62\linewidth}
      {\popxbold\fontsize{12}{13}\selectfont\color{popcream}Signed\ \textbullet\ {\poplogo OnBuch\color{poporange}+}}\par\vspace{2pt}
      {\raggedright\popsemi\fontsize{8}{10.5}\selectfont\color{popline}\MakeUppercase{OnBuch+ teaching team}\\\MakeUppercase{Follows the official MINESEC syllabus}\par}
    \end{minipage}\hfill
    {\poplogo\fontsize{22}{22}\selectfont\color{popcream}OnBuch\color{poporange}+}
  \end{tcolorbox}%
}

% ============================================================
% Page de garde
% #1 titre · #2 matière · #3 niveau · #4 module · #5 n° leçon · #6 durée ·
% #7 description / objectifs (texte ou itemize)
% ============================================================
\newcommand{\pagegarde}[7]{%
  \thispagestyle{empty}
  \newgeometry{a4paper,margin=1.7cm,top=1.6cm,bottom=1.4cm}%
  \begin{tikzpicture}[remember picture,overlay]
    \fill[poporange!18] ([xshift=-3.2cm,yshift=-3.6cm]current page.north east) circle (4.6cm);
    \fill[poppurple!14] ([xshift=2.4cm,yshift=7.6cm]current page.south west) circle (3.8cm);
  \end{tikzpicture}%
  {\poplogo\fontsize{30}{30}\selectfont\color{popink}OnBuch\color{poporange}+}\hfill
  \raisebox{0.6ex}{\poppill{Signed course \textbullet\ Premium}{popink}{popcream}}\par
  \vspace{1pt}\popsquiggle{poppurple}\par
  \vspace{8pt}
  \begin{tcolorbox}[enhanced,colback=poporange,colframe=popink,boxrule=2.8pt,arc=13pt,
      drop shadow=popink,left=18pt,right=18pt,top=12pt,bottom=13pt,after skip=10pt]
    \poppill{#2 \textbullet\ #3}{popink}{popcream}\hspace{5pt}\poppill{#4}{white}{popink}\par\vspace{8pt}
    {\popdisplay\fontsize{14}{14}\selectfont\color{popink}CHAPTER #5}\par\vspace{4pt}
    {\raggedright\hyphenpenalty=10000\exhyphenpenalty=10000\popdisplay\fontsize{25}{27}\selectfont\color{popcream}\MakeUppercase{#1}\par}
    \vspace{8pt}
    {\popxbold\fontsize{9.5}{9.5}\selectfont\color{popink}\MakeUppercase{Suggested time: #6}}
  \end{tcolorbox}
  \begin{tcolorbox}[enhanced,colback=white,colframe=popink,boxrule=2.2pt,arc=10pt,
      drop shadow=popink,left=16pt,right=16pt,top=10pt,bottom=10pt,after skip=8pt]
    \poppill{In this chapter}{poppurple}{white}\par\vspace{5pt}
    \popsemi\fontsize{9.3}{12.6}\selectfont\color{popink}#7
  \end{tcolorbox}
  \noindent\begin{minipage}[t]{0.487\linewidth}
    \begin{tcolorbox}[enhanced,colback=popgreen,colframe=popink,boxrule=2.2pt,arc=10pt,
        drop shadow=popink,left=11pt,right=11pt,top=10pt,bottom=10pt,height=3.1cm,valign=center]
      \tcbox[colback=white,colframe=popink,boxrule=1.3pt,arc=5pt,boxsep=0pt,nobeforeafter,
        left=3pt,right=3pt,top=3pt,bottom=3pt,valign=center]{\qrcode[height=1.9cm]{https://play.google.com/store/apps/details?id=com.luvvix.onbuch&hl=en}}%
      \hspace{8pt}\begin{minipage}[c]{0.5\linewidth}
        {\popdisplay\fontsize{11}{12.5}\selectfont\color{white}\MakeUppercase{Download OnBuch}}\par\vspace{4pt}
        {\raggedright\popsemi\fontsize{8.4}{11}\selectfont\color{white}Free \textbullet\ Google Play\\AI tutor, past papers, results\par}
      \end{minipage}
    \end{tcolorbox}
  \end{minipage}\hfill
  \begin{minipage}[t]{0.487\linewidth}
    \begin{tcolorbox}[enhanced,colback=poppurple,colframe=popink,boxrule=2.2pt,arc=10pt,
        drop shadow=popink,left=11pt,right=11pt,top=10pt,bottom=10pt,height=3.1cm,valign=center]
      \tikz[baseline=-0.7ex]\node[circle,fill=white,draw=popink,line width=1.3pt,minimum size=1.6cm,inner sep=0]{\popdisplay\fontsize{24}{24}\selectfont\color{poppurple}+};%
      \hspace{8pt}\begin{minipage}[c]{0.6\linewidth}
        {\popdisplay\fontsize{11}{12.5}\selectfont\color{white}\MakeUppercase{Go OnBuch+}}\par\vspace{4pt}
        {\raggedright\popsemi\fontsize{8.4}{11}\selectfont\color{white}All signed courses, unlimited AI tutor, PDF sheets — OnBuch+ tab in the app\par}
      \end{minipage}
    \end{tcolorbox}
  \end{minipage}
  \vspace{8pt}
  \popcontenu
  \vfill
  \signatureonbuch
  \restoregeometry
  \newpage
}

% Rangée « Ce cours contient » (page de garde)
\newcommand{\poptile}[3]{% #1 couleur #2 gros texte #3 légende
  \begin{tcolorbox}[enhanced,colback=#1,colframe=popink,boxrule=2pt,arc=9pt,shadow={2.5pt}{-2.5pt}{0pt}{popink},
      left=6pt,right=6pt,top=9pt,bottom=9pt,halign=center,valign=center,height=1.75cm,width=\linewidth,nobeforeafter]
    {\popdisplay\fontsize{12.5}{13}\selectfont\color{white}\MakeUppercase{#2}}\par\vspace{3pt}
    {\centering\popsemi\fontsize{7.8}{9.5}\selectfont\color{white}#3\par}
  \end{tcolorbox}}
\newcommand{\popcontenu}{%
  \par\noindent{\popxboldls\fontsize{7.5}{7.5}\selectfont\color{popmuted}THIS SIGNED COURSE CONTAINS}\par\vspace{7pt}
  \noindent\begin{minipage}[t]{0.235\linewidth}\poptile{popblue}{Course}{complete, illustrated, step by step}\end{minipage}\hfill
  \begin{minipage}[t]{0.235\linewidth}\poptile{poporange}{Methods}{and worked examples}\end{minipage}\hfill
  \begin{minipage}[t]{0.235\linewidth}\poptile{poppink}{Exercises}{exam style + solutions}\end{minipage}\hfill
  \begin{minipage}[t]{0.235\linewidth}\poptile{popgreen}{Summary sheet}{the essentials to remember}\end{minipage}\par}

% Sommaire
\newtcolorbox{sommaire}{enhanced,colback=white,colframe=popink,boxrule=2.2pt,arc=10pt,
  drop shadow=popink,left=18pt,right=18pt,top=13pt,bottom=13pt,before skip=6pt,after skip=20pt,
  before upper={\poppill{Contents}{poppurple}{white}\par\vspace{9pt}\popsemi\fontsize{10.6}{14.5}\selectfont\color{popink}}}

% Signature de fin de cours — poussée tout en bas de la dernière page
\newcommand{\findecours}{%
  \par\vfill\nopagebreak
  \begin{center}\popsquiggle{poporange}\end{center}\vspace{4pt}
  \signatureonbuch
}
\AtBeginDocument{\def\llbracket{\mathopen{[\mkern-3.5mu[}}\def\rrbracket{\mathclose{]\mkern-3.5mu]}}} % intervalles d'entiers (glyphe ⟦ absent de la police)
% Glyphes absents de Fira Math : redéfinis à partir de glyphes disponibles
\AtBeginDocument{%
  \def\setminus{\mathbin{\backslash}}\let\smallsetminus\setminus
  \def\vdots{\mathord{\vbox{\baselineskip3.2pt\lineskiplimit0pt\kern4pt\hbox{.}\hbox{.}\hbox{.}}}}%
  \def\ddots{\mathinner{\mkern1mu\raise6.5pt\hbox{.}\mkern2mu\raise3.6pt\hbox{.}\mkern2mu\raise0.7pt\hbox{.}\mkern1mu}}%
  \def\triangle{\mathord{\scalebox{0.9}{$\symup{\Delta}$}}}%
  \def\nabla{\mathord{\raisebox{\depth}{\scalebox{1}[-1]{$\symup{\Delta}$}}}}%
  \def\checkmark{\popcheck{popdark}}%
  \def\star{\mathbin{\ast}}\def\bigstar{\mathord{\ast}}%
  \def\diamond{\mathbin{\scalebox{0.6}{$\Diamond$}}}%
  \def\bigcup{\mathop{\mathchoice{\raisebox{-0.3em}{\scalebox{1.7}{$\cup$}}}{\scalebox{1.25}{$\cup$}}{\cup}{\cup}}\displaylimits}%
  \def\bigcap{\mathop{\mathchoice{\raisebox{-0.3em}{\scalebox{1.7}{$\cap$}}}{\scalebox{1.25}{$\cap$}}{\cap}{\cap}}\displaylimits}%
  \def\lesseqqgtr{\mathrel{\lessgtr}}\def\bigoplus{\mathop{\oplus}\displaylimits}%
}
\providecommand{\ssi}{if and only if} % abréviation fréquente
\AtBeginDocument{\providecommand{\wideparen}{\overparen}} % arc de cercle

%% FIGFIX-BEGIN v3 — correctifs globaux des figures (docs/onbuch-generation/tools/figfix)
\usepackage{adjustbox}\usepackage{etoolbox}
% toute figure TikZ/pgfplots plus large que la ligne est réduite à la largeur disponible
\BeforeBeginEnvironment{tikzpicture}{\begin{adjustbox}{max width=\linewidth}}
\AfterEndEnvironment{tikzpicture}{\end{adjustbox}}
% légendes pgfplots lisibles par-dessus les courbes
\pgfplotsset{every axis/.append style={legend style={fill=white,fill opacity=0.92,text opacity=1,draw=black!15,inner sep=3pt}}}
% étiquettes d'arêtes (syntaxe edge["..."]) avec fond blanc
\tikzset{every edge quotes/.append style={fill=white,inner sep=1.2pt}}
% bulles de cartes mentales : texte à la ligne dans la bulle, bulle un peu plus grande
\tikzset{every concept/.append style={text width=2.0cm,align=center,minimum size=2.3cm,inner sep=2pt}}
%% FIGFIX-END
\begin{document}

\pagegarde{Non-infectious diseases}{Biology}{Upper Sixth}{Upper Sixth — Biology}{6}{4 periods}{This chapter explores the major non-infectious diseases affecting the human reproductive, nervous, endocrine, renal, cardiovascular and musculoskeletal systems. For each system, learners link normal structure and function to the causes, symptoms, prevention and management of its characteristic disorders, with emphasis on conditions of public-health importance in Cameroon.
\vspace{4pt}
\begin{multicols}{2}\small
\begin{itemize}
\item By the end of this chapter I can define non-infectious disease and distinguish its main categories (degenerative, hormonal, genetic, nutritional, neoplastic, autoimmune).
\item By the end of this chapter I can describe the causes, symptoms and prevention of common disorders of the male and female reproductive systems.
\item By the end of this chapter I can explain how disorders of the nervous system (e.g. epilepsy, stroke, meningitis complications, Parkinson's disease) arise from failures of neurones, synapses or brain regions.
\item By the end of this chapter I can relate endocrine diseases (diabetes mellitus, goitre, gigantism, cretinism) to hypo- or hyper-secretion of specific hormones.
\item By the end of this chapter I can describe kidney diseases (nephritis, kidney stones, renal failure) and the principle of dialysis.
\item By the end of this chapter I can explain cardiovascular diseases (hypertension, atherosclerosis, coronary heart disease, stroke) in terms of risk factors and vascular changes.
\end{itemize}\end{multicols}}

\begin{sommaire}
\begin{multicols}{2}
\begin{enumerate}
\item Section 1: Non-infectious diseases — overview and disorders of the reproductive systems
\item Section 2: Disorders of the nervous and endocrine systems
\item Section 3: Diseases of the renal and cardiovascular systems
\item Section 4: Disorders of the musculoskeletal system and integrated prevention of NCDs
\item Integration activity
\item Methods and worked examples
\item Exercises
\item Solutions to the exercises
\item Summary sheet
\end{enumerate}
\end{multicols}
\end{sommaire}

\begin{objectifs}
\begin{itemize}
\item By the end of this chapter I can define non-infectious disease and distinguish its main categories (degenerative, hormonal, genetic, nutritional, neoplastic, autoimmune).
\item By the end of this chapter I can describe the causes, symptoms and prevention of common disorders of the male and female reproductive systems.
\item By the end of this chapter I can explain how disorders of the nervous system (e.g. epilepsy, stroke, meningitis complications, Parkinson's disease) arise from failures of neurones, synapses or brain regions.
\item By the end of this chapter I can relate endocrine diseases (diabetes mellitus, goitre, gigantism, cretinism) to hypo- or hyper-secretion of specific hormones.
\item By the end of this chapter I can describe kidney diseases (nephritis, kidney stones, renal failure) and the principle of dialysis.
\item By the end of this chapter I can explain cardiovascular diseases (hypertension, atherosclerosis, coronary heart disease, stroke) in terms of risk factors and vascular changes.
\item By the end of this chapter I can describe musculoskeletal disorders (arthritis, osteoporosis, rickets, muscular dystrophy) and their causes.
\item By the end of this chapter I can interpret clinical data (blood pressure readings, blood/urine test results, X-ray descriptions) to suggest a likely diagnosis.
\item By the end of this chapter I can propose lifestyle and public-health measures that reduce the burden of non-infectious diseases in a Cameroonian community.
\end{itemize}
\end{objectifs}

\begin{prerequis}
\begin{itemize}
\item Structure and functions of the male and female reproductive systems (Lower Sixth)
\item Structure of the neurone, reflex arc and organisation of the nervous system
\item Endocrine glands, hormones and negative feedback control (e.g. insulin and blood glucose)
\item Structure of the kidney nephron, ultrafiltration and selective reabsorption
\item Structure of the heart, blood vessels and the double circulation; bones, joints and muscles
\end{itemize}
\end{prerequis}

\newpage

\coursec{Section 1: Non-infectious diseases — overview and disorders of the reproductive systems}

At the Regional Hospital of Bamenda, a 45‑year‑old trader is told her persistent fatigue and swollen ankles come from failing kidneys, while her neighbour's infertility after years of marriage is traced to an untreated pelvic infection. Neither illness was caused by a germ that could be \emph{caught} from someone else — both are \cle{non‑infectious diseases}, disorders arising from within the body's own systems. In Cameroon today, such conditions (hypertension, diabetes, kidney failure, arthritis) are silently overtaking malaria and cholera as causes of death and disability. Why do organs that worked perfectly for decades suddenly malfunction, and how can understanding the biology of each system help us prevent, detect and manage these diseases? This section equips you to answer exactly that.

\courssub{What are non‑infectious (non‑communicable) diseases?}

\begin{definition}[Non‑infectious disease]
A \cle{non‑infectious disease} (also called a \cle{non‑communicable disease, NCD}) is a medical condition that is \textbf{not caused by a transmissible pathogen} (bacterium, virus, fungus, parasite). It arises from genetic, physiological, environmental or lifestyle factors and cannot be spread directly from person to person.
\end{definition}

\begin{definition}[Neoplasm]
A \cle{neoplasm} is an abnormal mass of tissue whose growth exceeds and is uncoordinated with that of the normal tissues. A \cle{benign tumour} remains localised, does not invade surrounding tissue and does not metastasise. A \cle{malignant tumour} (cancer) invades locally and can spread to distant sites (metastasis).
\end{definition}

NCDs are commonly grouped into seven major categories. The flowchart below summarises each category with a typical example.

\begin{popfigure}
\begin{tikzpicture}[
  node distance=1.2cm and 2.5cm,
  cat/.style={rectangle, draw=popblue, thick, fill=popblueL, rounded corners, text width=3.2cm, align=center, font=\small},
  ex/.style={rectangle, draw=poporange, thick, fill=poporange!20, rounded corners, text width=3.2cm, align=center, font=\small},
  arrow/.style={->, thick, popdark}
]
\node[cat] (deg) {Degenerative};
\node[ex, right=of deg] (degEx) {Osteoarthritis};
\node[cat, below=of deg] (horm) {Hormonal / Metabolic};
\node[ex, right=of horm] (hormEx) {Diabetes mellitus};
\node[cat, below=of horm] (gen) {Genetic / Congenital};
\node[ex, right=of gen] (genEx) {Sickle cell disease};
\node[cat, below=of gen] (nut) {Nutritional‑deficiency};
\node[ex, right=of nut] (nutEx) {Iron‑deficiency anaemia};
\node[cat, below=of nut] (neo) {Neoplastic};
\node[ex, right=of neo] (neoEx) {Breast cancer};
\node[cat, below=of neo] (auto) {Autoimmune};
\node[ex, right=of auto] (autoEx) {Rheumatoid arthritis};
\node[cat, below=of auto] (life) {Lifestyle‑related};
\node[ex, right=of life] (lifeEx) {Hypertension};

\draw[arrow] (deg) -- (degEx);
\draw[arrow] (horm) -- (hormEx);
\draw[arrow] (gen) -- (genEx);
\draw[arrow] (nut) -- (nutEx);
\draw[arrow] (neo) -- (neoEx);
\draw[arrow] (auto) -- (autoEx);
\draw[arrow] (life) -- (lifeEx);
\end{tikzpicture}
\legende{Classification of non‑infectious diseases with one representative example per category.}
\end{popfigure}

\courssub{General risk factors and the growing burden in Cameroon}

The main modifiable and non‑modifiable risk factors for NCDs are:

\begin{itemize}
  \item \textbf{Age} — physiological reserve declines with time.
  \item \textbf{Heredity} — family history of hypertension, diabetes, certain cancers.
  \item \textbf{Diet} — high salt, saturated fats, refined sugars; low fruit/vegetable intake.
  \item \textbf{Tobacco} — smoking damages vasculature and is carcinogenic.
  \item \textbf{Alcohol} — excessive intake raises blood pressure and cancer risk.
  \item \textbf{Physical inactivity} — contributes to obesity, insulin resistance.
  \item \textbf{Environmental pollutants} — indoor smoke from biomass fuels, outdoor air pollution.
  \item \textbf{Urbanisation} — rapid lifestyle changes, reduced physical activity, processed diets.
\end{itemize}

Cameroon, like many sub‑Saharan countries, faces a \cle{double burden of disease}: infectious diseases remain prevalent while NCDs are rising sharply. The illustrative bar chart below shows the estimated proportion of total deaths attributed to NCDs versus infectious diseases over three recent decades (data are indicative, not official statistics).

\begin{popfigure}
\begin{tikzpicture}
\begin{axis}[
  ybar,
  bar width=1.2cm,
  width=\linewidth,
  height=7cm,
  symbolic x coords={1990,2000,2010,2020},
  xtick=data,
  ylabel={Percentage of total deaths (\%)},
  ymin=0,ymax=100,
  legend style={at={(0.5,-0.2)},anchor=north,legend columns=-1},
  nodes near coords,
  nodes near coords align={vertical},
  enlarge x limits=0.2,
  ymajorgrids=true,
  grid style=dashed,
]
\addplot[fill=popblue] coordinates {(1990,30) (2000,38) (2010,48) (2020,58)};
\addplot[fill=poporange] coordinates {(1990,70) (2000,62) (2010,52) (2020,42)};
\legend{NCDs (illustrative), Infectious diseases (illustrative)}
\end{axis}
\end{tikzpicture}
\legende{Illustrative trend: rising share of deaths due to NCDs in Cameroon (1990–2020).}
\end{popfigure}

\begin{aretenir}[Key points on NCD burden]
\begin{itemize}
  \item NCDs now account for more than half of all deaths in many urban Cameroonian settings.
  \item Early detection and lifestyle modification can prevent up to 80\% of premature NCD deaths (WHO estimate).
  \item Health systems must integrate NCD screening into primary care (e.g. blood pressure, blood glucose, cervical screening).
\end{itemize}
\end{aretenir}

\courssub{Male reproductive system disorders}

\begin{popfigure}
\begin{tikzpicture}[
  organ/.style={ellipse, draw=popdark, thick, fill=popblueL, text width=2.2cm, align=center, font=\small},
  label/.style={rectangle, draw=poporange, thick, fill=poporange!20, rounded corners, text width=2.5cm, align=center, font=\small, font=\bfseries},
  arrow/.style={->, thick, popdark}
]
\node[organ, align=center] (testisL) at (0,3){Testis\\(spermatogenesis)};
\node[organ, align=center] (epidL) at (3,3){Epididymis\\(maturation)};
\node[organ, align=center] (vasL) at (6,3){Vas deferens\\(transport)};
\node[organ, align=center] (pros) at (9,1.5){Prostate\\(seminal fluid)};
\node[organ] (semv) at (9,-0.5) {Seminal vesicles};
\node[organ] (ureth) at (9,-2.5) {Urethra};
\node[organ] (penis) at (11,0) {Penis};

% disorder labels
\node[label] (bp) at (9,3.5) {Benign prostatic hyperplasia (BPH)};
\node[label] (pc) at (9,4.5) {Prostate cancer};
\node[label] (tc) at (0,5) {Testicular cancer};
\node[label] (var) at (-2,3) {Varicocele};
\node[label] (hyd) at (-2,1.5) {Hydrocele};
\node[label] (cry) at (-2,0) {Cryptorchidism};

\draw[arrow] (testisL) -- (epidL);
\draw[arrow] (epidL) -- (vasL);
\draw[arrow] (vasL) -- (pros);
\draw[arrow] (pros) -- (ureth);
\draw[arrow] (semv) -- (pros);
\draw[arrow] (ureth) -- (penis);

% dashed lines to disorder labels
\draw[dashed, poporange] (pros) -- (bp);
\draw[dashed, poporange] (pros) -- (pc);
\draw[dashed, poporange] (testisL) -- (tc);
\draw[dashed, poporange] (testisL) -- (var);
\draw[dashed, poporange] (testisL) -- (hyd);
\draw[dashed, poporange] (testisL) -- (cry);
\end{tikzpicture}
\legende{Male reproductive system with common disorder sites highlighted.}
\end{popfigure}

\courssub{Prostate disorders}

\begin{propriete}[Benign prostatic hyperplasia (BPH)]
Non‑cancerous enlargement of the prostate gland, common after age 50. It compresses the urethra causing \cle{lower urinary tract symptoms (LUTS)}: hesitancy, weak stream, frequency, nocturia, incomplete emptying.
\end{propriete}

\begin{propriete}[Prostate cancer]
Malignant tumour of the prostate epithelium. Risk factors: age $>50$, family history, African ancestry, high‑fat diet. Early stages often asymptomatic; later may mimic BPH plus haematuria, bone pain (metastasis). Screening: serum \cle{prostate‑specific antigen (PSA)} test and digital rectal examination (DRE). Definitive diagnosis by trans‑rectal ultrasound‑guided biopsy.
\end{propriete}

\begin{methode}[PSA screening interpretation (simplified for examination)]
\begin{enumerate}
  \item Obtain serum PSA (ng/mL).
  \item Age‑adjusted reference (simplified): $<4.0$ ng/mL generally normal; $4.0$–$10.0$ ng/mL borderline; $>10.0$ ng/mL high suspicion. \emph{Note: modern guidelines use age‑specific ranges (e.g. 40–49 yr: 0–2.5; 50–59 yr: 0–3.5; 60–69 yr: 0–4.5; 70–79 yr: 0–6.5 ng/mL).}
  \item If PSA elevated, repeat after 4–6 weeks (rule out prostatitis).
  \item Persistent elevation $\rightarrow$ refer for DRE and possible biopsy.
\end{enumerate}
\end{methode}

\courssub{Testicular disorders}

\begin{itemize}
  \item \textbf{Testicular cancer} — most common solid tumour in males 15–35 years. Presents as painless lump; high cure rate if detected early (orchiectomy + chemotherapy).
  \item \textbf{Varicocele} — dilation of pampiniform plexus veins; raises scrotal temperature, impairs spermatogenesis $\rightarrow$ infertility.
  \item \textbf{Hydrocele} — fluid accumulation in tunica vaginalis; usually benign, may cause discomfort.
  \item \textbf{Cryptorchidism} — failure of testis to descend into scrotum; if untreated, increases risk of malignancy and infertility.
\end{itemize}

\courssub{Male infertility and erectile dysfunction}

\begin{propriete}[Causes of male infertility]
\begin{itemize}
  \item \textbf{Pre‑testicular}: hormonal (hypogonadotropic hypogonadism, hyperprolactinaemia).
  \item \textbf{Testicular}: varicocele, cryptorchidism, infection (mumps orchitis), genetic (Klinefelter syndrome), idiopathic.
  \item \textbf{Post‑testicular}: obstruction of vas deferens (congenital absence, post‑vasectomy, infection).
\end{itemize}
\end{propriete}

\begin{propriete}[Erectile dysfunction (ED)]
Inability to achieve/maintain erection sufficient for intercourse. \textbf{Vascular} (atherosclerosis, hypertension), \textbf{neurological} (diabetic neuropathy), \textbf{psychological} (stress, depression), \textbf{hormonal} (low testosterone), \textbf{iatrogenic} (antihypertensives, antidepressants).
\end{propriete}

\begin{exoresolu}[Worked example: interpreting a semen analysis (WHO 2021, 5th edition)]
\textbf{Statement:} A 32‑year‑old man presents with 2 years of infertility. Semen analysis shows: volume 1.8 mL, concentration 8 million/mL, progressive motility 25\%, normal morphology 3\%. WHO 2021 lower reference limits (5th centiles): volume $\ge$1.5 mL, concentration $\ge$15 million/mL, progressive motility $\ge$30\%, normal morphology $\ge$4\%. Classify the abnormalities.

\tcblower
\textbf{Solution:}
\begin{itemize}
  \item Volume 1.8 mL $\ge$ 1.5 mL $\rightarrow$ normal.
  \item Concentration 8 million/mL $<$ 15 million/mL $\rightarrow$ \cle{oligozoospermia}.
  \item Progressive motility 25\% $<$ 30\% $\rightarrow$ \cle{asthenozoospermia}.
  \item Normal morphology 3\% $<$ 4\% $\rightarrow$ \cle{teratozoospermia}.
\end{itemize}
Combined pattern = \cle{oligoasthenoteratozoospermia (OAT)}. Further hormonal profile (FSH, LH, testosterone) and scrotal ultrasound are indicated to identify cause.
\end{exoresolu}

\courssub{Female reproductive system disorders}

\begin{popfigure}
\begin{tikzpicture}[
  organ/.style={ellipse, draw=popdark, thick, fill=poppink!30, text width=2.2cm, align=center, font=\small},
  label/.style={rectangle, draw=poppurple, thick, fill=poppurple!20, rounded corners, text width=2.5cm, align=center, font=\small, font=\bfseries},
  arrow/.style={->, thick, popdark}
]
\node[organ, align=center] (ovL) at (-3,2){Ovary\\(oogenesis)};
\node[organ, align=center] (tubeL) at (0,2){Fallopian tube\\(fertilisation)};
\node[organ, align=center] (uterus) at (3,0){Uterus\\(implantation)};
\node[organ] (cerv) at (3,-2) {Cervix};
\node[organ] (vag) at (3,-4) {Vagina};
\node[organ] (breast) at (6,2) {Breast};

% disorder labels
\node[label] (pcos) at (-3,4) {PCOS};
\node[label] (cyst) at (-5,2) {Ovarian cyst};
\node[label] (fib) at (3,2) {Fibroids (leiomyoma)};
\node[label] (endo) at (5,0) {Endometriosis};
\node[label] (cervCa) at (6,-2) {Cervical cancer};
\node[label] (ovCa) at (-3,0) {Ovarian cancer};
\node[label] (brCa) at (8,2) {Breast cancer};

\draw[arrow] (ovL) -- (tubeL);
\draw[arrow] (tubeL) -- (uterus);
\draw[arrow] (uterus) -- (cerv);
\draw[arrow] (cerv) -- (vag);

\draw[dashed, poppurple] (ovL) -- (pcos);
\draw[dashed, poppurple] (ovL) -- (cyst);
\draw[dashed, poppurple] (ovL) -- (ovCa);
\draw[dashed, poppurple] (uterus) -- (fib);
\draw[dashed, poppurple] (uterus) -- (endo);
\draw[dashed, poppurple] (cerv) -- (cervCa);
\draw[dashed, poppurple] (breast) -- (brCa);
\end{tikzpicture}
\legende{Female reproductive system with common disorder sites highlighted.}
\end{popfigure}

\courssub{Menstrual disorders (hormonal basis)}

\begin{itemize}
  \item \textbf{Amenorrhoea} — absence of menses. \emph{Primary}: no menarche by age 16; \emph{Secondary}: cessation for $\ge$3 cycles. Causes: hypothalamic‑pituitary failure (low GnRH/FSH/LH), ovarian failure (premature ovarian insufficiency), outflow tract obstruction, pregnancy.
  \item \textbf{Dysmenorrhoea} — painful menstruation. \emph{Primary}: prostaglandin‑mediated uterine contractions (no pathology). \emph{Secondary}: endometriosis, adenomyosis, fibroids.
  \item \textbf{Menorrhagia} — excessively heavy/prolonged bleeding ($>80$ mL/cycle). Causes: anovulatory cycles (unopposed oestrogen), fibroids, coagulopathies, copper IUD.
\end{itemize}

\begin{propriete}[Hormonal control of the menstrual cycle]
\begin{itemize}
  \item \textbf{Follicular phase}: rising \cle{oestradiol} from developing follicle $\rightarrow$ proliferative endometrium.
  \item \textbf{Ovulation}: LH surge triggered by positive feedback of oestradiol.
  \item \textbf{Luteal phase}: \cle{progesterone} from corpus luteum $\rightarrow$ secretory endometrium.
  \item \textbf{Menses}: withdrawal of progesterone+oestradiol $\rightarrow$ spiral artery vasoconstriction $\rightarrow$ endometrial shedding.
\end{itemize}
\end{propriete}

\courssub{Ovarian and uterine disorders}

\begin{itemize}
  \item \textbf{Ovarian cysts} — fluid‑filled sacs; functional (follicular, corpus luteum) usually resolve; pathological (dermoid, cystadenoma) may require surgery.
  \item \textbf{Polycystic ovary syndrome (PCOS)} — hyperandrogenism, oligo‑/anovulation, polycystic ovaries on ultrasound (Rotterdam criteria: 2 of 3). Insulin resistance central; long‑term risks: type 2 diabetes, endometrial hyperplasia, cardiovascular disease.
  \item \textbf{Uterine fibroids (leiomyomas)} — benign smooth‑muscle tumours; oestrogen‑dependent; cause menorrhagia, pressure symptoms, infertility.
  \item \textbf{Adenomyosis} — endometrial glands within myometrium; causes dysmenorrhoea, menorrhagia, enlarged uterus.
  \item \textbf{Endometriosis} — ectopic endometrial tissue (ovaries, peritoneum); causes chronic pelvic pain, dysmenorrhoea, infertility via adhesions and inflammatory milieu.
\end{itemize}

\courssub{Gynaecological cancers}

\begin{center}
\begin{tabularx}{\linewidth}{|l|X|X|X|}
\hline
\rowcolor{popblueL}
\textbf{Cancer} & \textbf{Main risk factors} & \textbf{Early signs} & \textbf{Screening / early detection} \\
\hline
Cervix & HPV (types 16,18), early sexual debut, multiple partners, smoking, immunosuppression & Post‑coital bleeding, intermenstrual bleeding, foul discharge & \cle{Pap smear} (cytology) every 3 yr (25–65 yr); HPV DNA testing; visual inspection with acetic acid (VIA) in low‑resource settings \\
\hline
Ovary & Nulliparity, late menopause, BRCA1/2 mutations, family history & Vague: bloating, pelvic discomfort, early satiety, urinary frequency & No population screening; high‑risk: transvaginal ultrasound + CA‑125 annually \\
\hline
Breast & Female sex, age, family history, BRCA1/2, early menarche/late menopause, nulliparity, hormone therapy, alcohol, obesity & Painless lump, skin dimpling, nipple retraction, discharge & \cle{Breast self‑examination (BSE)} monthly; clinical breast exam annually; mammography 40–50 yr (or earlier if high risk) \\
\hline
\end{tabularx}
\end{center}

\courssub{Pelvic inflammatory disease (PID) and ectopic pregnancy}

\begin{propriete}[Pelvic inflammatory disease]
Ascending infection (often \emph{Chlamydia trachomatis}, \emph{Neisseria gonorrhoeae}) from cervix to uterus, tubes, peritoneum. Acute PID: lower abdominal pain, fever, cervical motion tenderness. Sequelae: tubal scarring $\rightarrow$ \cle{blocked Fallopian tubes} $\rightarrow$ infertility; increased risk of \cle{ectopic pregnancy}.
\end{propriete}

\begin{attention}[Ectopic pregnancy — medical emergency]
Implantation outside uterine cavity (95\% in Fallopian tube). Presents with amenorrhoea, unilateral pelvic pain, vaginal bleeding, shoulder tip pain (diaphragmatic irritation). Rupture $\rightarrow$ haemorrhagic shock. Diagnosis: $\beta$-hCG + transvaginal ultrasound (empty uterus, adnexal mass). Management: methotrexate (unruptured, stable) or laparoscopic salpingectomy/salpingostomy.
\end{attention}

\courssub{Prevention and management of reproductive NCDs}

\begin{methode}[Breast self‑examination (BSE) — 5 steps]
\begin{enumerate}
  \item \textbf{Visual inspection} — stand before mirror, arms at sides, then raised; look for asymmetry, skin changes, nipple inversion.
  \item \textbf{Palpation lying down} — place pillow under shoulder, use flat fingers in circular motion covering whole breast and axilla.
  \item \textbf{Palpation standing} — in shower with soapy fingers, same systematic pattern.
  \item \textbf{Nipple check} — gently squeeze each nipple; note any discharge (bloody, clear).
  \item \textbf{Record and report} — note date, findings; any new lump or change $\rightarrow$ consult health worker within 2 weeks.
\end{enumerate}
\end{methode}

\begin{methode}[Rationale for cervical screening (Pap smear)]
\begin{enumerate}
  \item HPV infection $\rightarrow$ cervical intraepithelial neoplasia (CIN) grades 1–3 over 10–20 years.
  \item Pap smear detects abnormal squamous cells (ASC‑US, LSIL, HSIL) before invasion.
  \item Treatment of CIN (cryotherapy, LEEP) prevents progression to invasive carcinoma.
  \item Screening interval 3 years balances benefit/harm; HPV vaccination (9–14 yr) reduces future burden.
\end{enumerate}
\end{methode}

\begin{attention}[Common mistake: confusing STIs with non‑infectious reproductive disorders]
\begin{itemize}
  \item \textbf{STIs} (e.g. gonorrhoea, syphilis, HIV) are \emph{infectious} — caused by transmissible pathogens.
  \item \textbf{Fibroids, BPH, PCOS, endometriosis, prostate cancer} are \emph{non‑infectious} — arise from hormonal, genetic, lifestyle factors.
  \item In exam tables, always separate \textbf{cause} (pathogen vs hormonal/genetic), \textbf{symptom}, and \textbf{prevention} (condom/vaccine vs lifestyle/screening).
\end{itemize}
\end{attention}

\begin{propriete}[Hormonal imbalance and reproductive disorders]
\begin{itemize}
  \item \textbf{Oestrogen excess / progesterone deficiency} $\rightarrow$ endometrial hyperplasia, menorrhagia, fibroid growth, increased breast/endometrial cancer risk.
  \item \textbf{Androgen excess} (e.g. PCOS) $\rightarrow$ anovulation, hirsutism, acne, insulin resistance.
  \item \textbf{Testosterone deficiency} (male hypogonadism) $\rightarrow$ low libido, ED, reduced spermatogenesis, osteoporosis.
  \item \textbf{Prolactin excess} (prolactinoma) $\rightarrow$ amenorrhoea/galactorrhoea (female), hypogonadism (male).
\end{itemize}
\end{propriete}

\begin{savaistu}[HPV vaccination in Cameroon]
Cameroon introduced the quadrivalent HPV vaccine into the routine immunisation programme for girls aged 9–13 years in 2020 (GAVI support). Coverage remains below 50\% in many regions due to vaccine hesitancy, cold‑chain challenges and school‑based delivery gaps. Catch‑up campaigns and community sensitisation are ongoing. Vaccination of boys is under consideration for future expansion.
\end{savaistu}

\begin{aretenir}[Exam tip: three‑column cause–symptom–prevention table]
GCE questions frequently ask: \emph{“Complete the table below for the following conditions: cervical cancer, BPH, PCOS.”} Practise drawing a table with columns \textbf{Condition}, \textbf{Cause / Risk factors}, \textbf{Key symptom(s)}, \textbf{Prevention / Early detection}. Fill each row concisely; use bullet points inside cells if needed. This structure earns full marks for organisation and completeness.
\end{aretenir}

\begin{exercice}[Practice: cause–symptom–prevention table]
Complete the table for the four conditions listed. Write your answer in the three‑column format described above.

\begin{center}
\begin{tabularx}{\linewidth}{|l|X|X|X|}
\hline
\rowcolor{popblueL}
\textbf{Condition} & \textbf{Cause / Risk factors} & \textbf{Key symptom(s)} & \textbf{Prevention / Early detection} \\
\hline
Cervical cancer & & & \\
\hline
Benign prostatic hyperplasia & & & \\
\hline
Polycystic ovary syndrome & & & \\
\hline
Testicular cancer & & & \\
\hline
\end{tabularx}
\end{center}
\end{exercice}

\begin{corrige}[Exercise 1]
\begin{center}
\begin{tabularx}{\linewidth}{|l|X|X|X|}
\hline
\rowcolor{popblueL}
\textbf{Condition} & \textbf{Cause / Risk factors} & \textbf{Key symptom(s)} & \textbf{Prevention / Early detection} \\
\hline
Cervical cancer & HPV 16/18 infection; early sexual debut; multiple partners; smoking; immunosuppression & Post‑coital bleeding; intermenstrual bleeding; foul discharge & HPV vaccination (9–14 yr); Pap smear q3 yr (25–65 yr); VIA in low‑resource settings \\
\hline
Benign prostatic hyperplasia & Age $>50$; hormonal changes (DHT); family history & LUTS: hesitancy, weak stream, frequency, nocturia, incomplete emptying & No proven primary prevention; early DRE + PSA for early detection; lifestyle (weight control, exercise) may reduce severity \\
\hline
Polycystic ovary syndrome & Insulin resistance; hyperandrogenism; genetic predisposition; obesity & Oligo/amenorrhoea; hirsutism; acne; polycystic ovaries on ultrasound & Weight loss (5–10\%); combined oral contraceptives for cycle control; metformin if insulin resistant; regular endometrial surveillance \\
\hline
Testicular cancer & Cryptorchidism; family history; age 15–35; Caucasian ethnicity & Painless testicular lump; heaviness; sometimes gynaecomastia (hCG) & Monthly testicular self‑exam (15–35 yr); early orchiectomy if mass detected; no population screening \\
\hline
\end{tabularx}
\end{center}
\end{corrige}

\coursec{Section 2: Disorders of the nervous and endocrine systems}

In Section 1 we saw that non-infectious diseases (NCDs) are not caused by pathogens but by malfunctions of the body's own structures, and we applied this idea to the reproductive systems. We now turn to two systems that act as the body's \cle{control networks}: the \cle{nervous system}, which sends rapid electrical messages, and the \cle{endocrine system}, which sends slower chemical messages (hormones) through the blood. Because these two systems coordinate every other organ, their disorders often have widespread effects — and several of them, such as stroke and type 2 diabetes, are among the fastest-growing health problems in Cameroonian towns today.

\courssub{Locating nervous-system disorders: a quick review}

The nervous system is divided into the \cle{central nervous system} (CNS: brain and spinal cord) and the \cle{peripheral nervous system} (PNS: cranial and spinal nerves). The functional cell is the \cle{neurone}: impulses travel along the axon (often insulated by a \cle{myelin sheath}) and cross junctions called \cle{synapses} using chemical \cle{neurotransmitters}. This organisation gives us a logical way to classify disorders:

\begin{itemize}
\item \textbf{Electrical disorders} — abnormal firing of neurones (epilepsy).
\item \textbf{Vascular disorders} — failure of blood supply to the brain (stroke).
\item \textbf{Degenerative disorders} — progressive death of neurones (Parkinson's, Alzheimer's).
\item \textbf{Autoimmune/structural disorders} — damage to myelin (multiple sclerosis), or irritation of pain pathways (neuralgia, migraine).
\end{itemize}

\begin{popfigure}
\begin{tikzpicture}[scale=1.0, every node/.style={font=\small}]
% cell body
\draw[fill=popblueL, draw=popblue, thick] (0,0) circle (0.55);
\draw[fill=popblue, draw=popblue] (0,0) circle (0.15);
\node at (0,-0.95) {cell body};
% dendrites
\draw[thick, popblue] (-0.4,0.35) -- (-1.1,0.9);
\draw[thick, popblue] (-0.5,0.1) -- (-1.3,0.25);
\draw[thick, popblue] (-0.4,-0.35) -- (-1.1,-0.85);
\node at (-1.7,0.9) {dendrites};
% axon with myelin
\draw[thick, popink] (0.55,0) -- (5.6,0);
\foreach \x in {1.0,2.2,3.4,4.6}{
  \draw[fill=popgreenL, draw=popgreen, thick] (\x,-0.18) rectangle (\x+0.8,0.18);
}
\node at (2.8,0.55) {myelin sheath};
\node at (5.9,0.35) {axon};
% terminals
\draw[thick, popink] (5.6,0) -- (6.4,0.5);
\draw[thick, popink] (5.6,0) -- (6.4,-0.5);
\draw[fill=poporange, draw=poporange] (6.55,0.55) circle (0.12);
\draw[fill=poporange, draw=poporange] (6.55,-0.55) circle (0.12);
\node[align=center] at (7.7,0) {synaptic\\terminals};
% annotations of disorders
\node[align=center, text=poppurple] at (2.8,-1.1) {demyelination\\(multiple sclerosis)};
\draw[->, poppurple, thick] (2.8,-0.75) -- (2.6,-0.22);
\node[align=center, text=poppurple] at (7.7,1.3) {transmitter deficiency\\(Parkinson's)};
\draw[->, poppurple, thick] (7.2,0.9) -- (6.6,0.62);
\node[align=center, text=poppurple] at (0,1.5) {abnormal discharge\\(epilepsy)};
\draw[->, poppurple, thick] (0,1.15) -- (0,0.6);
\end{tikzpicture}
\legende{A neurone with the main sites at which nervous-system disorders arise.}
\end{popfigure}

\courssub{Epilepsy}

\begin{definition}[Epilepsy]
\cle{Epilepsy} is a chronic non-infectious disorder of the brain in which groups of neurones suddenly produce \cle{abnormal, excessive electrical discharges}, causing recurrent \cle{seizures} (fits).
\end{definition}

Seizures vary in form. In \textbf{generalised tonic-clonic seizures} (grand mal), the person loses consciousness, the body stiffens (tonic phase) then jerks rhythmically (clonic phase). In \textbf{absence seizures} (petit mal), common in children, the person briefly stares blankly for a few seconds without falling. Triggers include fever, flashing lights, lack of sleep, stress and missed medication. Epilepsy is controlled — though usually not cured — by regular \cle{anticonvulsant} medicines; interrupting treatment is a common cause of relapse.

\begin{methode}[First aid for an epileptic seizure]
\begin{enumerate}
\item Stay calm and note the time; most seizures stop within a few minutes.
\item Protect the person from injury: move away hard or sharp objects, cushion the head.
\item Do \textbf{not} restrain the person's movements and do \textbf{not} put anything (spoon, stick, fingers) into the mouth.
\item Loosen tight clothing around the neck.
\item When the jerking stops, place the person in the \cle{recovery position} (on the side) so saliva drains and breathing is free.
\item Stay until full consciousness returns; call for medical help if the seizure lasts more than about five minutes or repeats.
\end{enumerate}
\end{methode}

\begin{attention}[Stigma is not science]
In some Cameroonian communities epilepsy is wrongly attributed to witchcraft or thought to be contagious through saliva. Epilepsy is \textbf{neither contagious nor supernatural}: it is an electrical disorder of the brain, treatable with anticonvulsant medicines. Forcing objects into a patient's mouth causes broken teeth and choking — a harmful practice based on a myth.
\end{attention}

\courssub{Stroke (cerebrovascular accident)}

\begin{definition}[Stroke]
A \cle{stroke} is the sudden death of brain tissue caused by an interruption of its blood supply. It is \textbf{ischaemic} when a clot blocks a cerebral artery, and \textbf{haemorrhagic} when a weakened vessel bursts and bleeds into the brain.
\end{definition}

The strongest risk factor is \cle{hypertension} (high blood pressure), together with smoking, diabetes and high cholesterol. Because neurones deprived of oxygen die within minutes, a stroke is a medical emergency. Typical warning signs can be recalled with the word \textbf{FAST}: \textbf{F}ace drooping on one side, \textbf{A}rm weakness (one arm cannot be raised), \textbf{S}peech slurred or lost, \textbf{T}ime to rush to hospital. Sudden severe headache and one-sided numbness are also common. Rapid hospital treatment (clot-dissolving drugs for ischaemic stroke, given early) greatly improves recovery — "time is brain". Prevention rests on controlling blood pressure, not smoking, exercising and limiting salt and alcohol.

\courssub{Degenerative, autoimmune and other nervous disorders}

\textbf{Parkinson's disease} results from the progressive loss of \cle{dopamine-producing neurones} in a region of the midbrain (the substantia nigra). Dopamine is needed for smooth, coordinated movement, so patients show \textbf{tremor at rest}, muscle rigidity and slow movement (bradykinesia). Treatment replaces dopamine activity (e.g. L-dopa).

\textbf{Alzheimer's disease} is a degenerative disorder causing \textbf{progressive memory loss}, confusion and personality change, mainly in the elderly; care is mostly supportive.

\textbf{Multiple sclerosis (MS)} is an \cle{autoimmune} disease in which the immune system attacks the \cle{myelin sheath} of CNS neurones. Demyelination slows or blocks impulse conduction, producing weakness, numbness and vision problems that come and go.

\textbf{Migraine} is recurrent severe headache, often one-sided with nausea and light sensitivity; \textbf{neuralgia} is sharp pain along a nerve (e.g. sciatica). Finally, remember that \textbf{alcohol, drugs and head injuries} damage the nervous system: chronic alcohol destroys neurones and the liver's detoxification capacity, while head injuries from motorcycle accidents — sadly common in Yaoundé and Douala — cause concussions and lasting brain damage. Wearing a helmet is disease prevention.

\courssub{The endocrine system and negative feedback}

The \cle{endocrine glands} (pituitary, thyroid, parathyroids, adrenals, pancreas, ovaries, testes) secrete \cle{hormones} directly into the blood. Hormone release is normally controlled by \cle{negative feedback}: when the blood level of a hormone (or of the substance it controls) rises above the set point, further secretion is switched off; when it falls, secretion resumes. This keeps internal conditions stable (homeostasis).

\begin{propriete}[Negative feedback and endocrine disease]
In a healthy feedback loop, a rise in the controlled variable inhibits the gland, and a fall stimulates it. If a gland becomes \textbf{over-active} (hyper-secretion), the controlled variable stays abnormally high because the gland no longer responds to the "off" signal; if it becomes \textbf{under-active} (hypo-secretion), the variable stays abnormally low. Hence endocrine diseases are classified as \cle{hypo-secretion} or \cle{hyper-secretion} disorders.
\end{propriete}

\begin{popfigure}
\begin{tikzpicture}[scale=0.95, every node/.style={font=\small}]
% glands (left column) and diseases (right column)
\node[fill=popblueL, draw=popblue, rounded corners, minimum width=3.4cm] (pit) at (0,4.4) {Pituitary};
\node[fill=popblueL, draw=popblue, rounded corners, minimum width=3.4cm] (thy) at (0,2.9) {Thyroid};
\node[fill=popblueL, draw=popblue, rounded corners, minimum width=3.4cm] (adr) at (0,1.4) {Adrenals};
\node[fill=popblueL, draw=popblue, rounded corners, minimum width=3.4cm] (pan) at (0,-0.1) {Pancreas (islets)};
\node[fill=popblueL, draw=popblue, rounded corners, minimum width=3.4cm] (gon) at (0,-1.6) {Ovaries / Testes};
% disease columns
\node[align=left, anchor=west, text=popdark] at (4.8,4.4) {\textbf{Hypo:} dwarfism; diabetes insipidus (ADH)\\ \textbf{Hyper:} gigantism (child), acromegaly (adult)};
\node[align=left, anchor=west, text=popdark] at (4.8,2.9) {\textbf{Hypo:} cretinism (child), myxoedema, goitre\\ \textbf{Hyper:} hyperthyroidism (exophthalmic goitre)};
\node[align=left, anchor=west, text=popdark] at (4.8,1.4) {\textbf{Hypo:} Addison's disease\\ \textbf{Hyper:} Cushing's syndrome};
\node[align=left, anchor=west, text=popdark] at (4.8,-0.1) {\textbf{Hypo (insulin):} diabetes mellitus type 1\\ \textbf{Resistance:} diabetes mellitus type 2};
\node[align=left, anchor=west, text=popdark] at (4.8,-1.6) {\textbf{Hypo:} infertility, delayed puberty\\ \textbf{Hyper:} precocious puberty};
\foreach \g in {pit,thy,adr,pan,gon}{
  \draw[->, popmuted, thick] (\g.east) -- ++(1.0,0);
}
\end{tikzpicture}
\legende{Major endocrine glands with the diseases caused by hypo- and hyper-secretion of each.}
\end{popfigure}

\courssub{Diabetes mellitus}

\begin{definition}[Diabetes mellitus]
\cle{Diabetes mellitus} is a chronic disorder in which blood \cle{glucose} remains abnormally high (hyperglycaemia) because the hormone \cle{insulin} is absent, insufficient or ineffective.
\end{definition}

\textbf{Type 1 diabetes} is an \textbf{autoimmune} destruction of the insulin-secreting $\beta$-cells of the islets of Langerhans, usually appearing in childhood or adolescence: insulin is \emph{deficient}, so daily insulin injections are essential for survival. \textbf{Type 2 diabetes}, the commonest form in adults, results from \cle{insulin resistance}: the pancreas still secretes insulin, but body cells respond poorly to it. Risk factors include obesity, physical inactivity, age and family history — it is increasingly common in urban Cameroon with changing diets and sedentary lifestyles.

Symptoms of both types follow directly from hyperglycaemia: when blood glucose exceeds the \cle{renal threshold} (about $10\ \mathrm{mmol\,L^{-1}}$), the kidneys cannot reabsorb it all, so glucose spills into the urine, dragging water with it (\cle{polyuria}), causing thirst (\cle{polydipsia}); because cells cannot use glucose properly, the body breaks down fat and muscle (\textbf{weight loss}). Diagnosis is by \textbf{blood glucose tests}: fasting blood glucose, or an \cle{oral glucose tolerance test} (OGTT) in which glucose is drunk and blood glucose is followed over two hours.

\begin{popfigure}
\begin{tikzpicture}
\begin{axis}[
  width=11cm, height=6.5cm,
  xlabel={Time after meal (min)}, ylabel={Blood glucose (mmol/L)},
  xmin=0, xmax=150, ymin=0, ymax=14,
  xtick={0,30,60,90,120,150}, ytick={0,2,4,6,8,10,12,14},
  grid=both, grid style={popline!40},
  legend style={at={(0.03,0.97)}, anchor=north west, font=\small},
]
\addplot[very thick, popblue, smooth] coordinates {(0,4.5) (30,8.0) (60,6.5) (90,5.2) (120,4.6) (150,4.5)};
\addlegendentry{Healthy person}
\addplot[very thick, poporange, smooth] coordinates {(0,7.5) (30,12.5) (60,11.5) (90,10.0) (120,9.0) (150,8.4)};
\addlegendentry{Diabetic patient}
\addplot[dashed, popmuted] coordinates {(0,10.0) (150,10.0)};
\node[font=\footnotesize, text=popmuted] at (axis cs:118,10.4) {renal threshold ($\approx 10$ mmol/L)};
\end{axis}
\end{tikzpicture}
\legende{Blood glucose after a meal: the diabetic curve starts higher, rises further (above the renal threshold, so glucose appears in urine) and returns to normal much more slowly.}
\end{popfigure}

\begin{exoresolu}[Interpreting a glucose-tolerance curve]
Two patients each drink a glucose solution. Patient A's blood glucose rises from 4.6 to 8.1 mmol/L and returns below 6 mmol/L within 2 hours. Patient B's rises from 7.4 to 12.6 mmol/L and is still 9.0 mmol/L after 2 hours. Which patient is diabetic? Justify.
\tcblower
\textbf{Solution.} Patient B is diabetic. Justification: (1) her \emph{fasting} glucose (7.4 mmol/L) is already above the normal range; (2) the peak (12.6 mmol/L) exceeds the renal threshold (about 10 mmol/L), so glucose appears in the urine; (3) glucose remains elevated after 2 hours, showing that insulin action is insufficient or ineffective. Patient A shows a normal response: a moderate peak followed by a rapid return to the set point under insulin action.
\end{exoresolu}

Long-term \textbf{complications} of poorly controlled diabetes include \cle{retinopathy} (eye damage, blindness), \cle{neuropathy} (nerve damage, foot ulcers that may lead to amputation), and \cle{nephropathy} (kidney damage). Management combines insulin (type 1, and sometimes type 2), a balanced diet low in refined sugars, regular exercise and weight control.

\begin{attention}[Mellitus $\neq$ insipidus]
Do not confuse \textbf{diabetes mellitus} (insulin/glucose problem, sugar in urine) with \textbf{diabetes insipidus} (deficiency of \cle{ADH} released from the posterior pituitary, or kidney unresponsiveness to it). Both cause polyuria and thirst, but in diabetes insipidus the urine is dilute and contains \emph{no glucose} — it is a water-balance disorder, not a sugar disorder.
\end{attention}

\begin{attention}[Type 2 diabetes is not just "eating sugar"]
Attributing type 2 diabetes solely to sugar intake is a common mistake. The core defect is \cle{insulin resistance}, driven by several interacting risk factors: genetic predisposition, obesity (especially abdominal fat), lack of exercise and age. Sugar is only one contributor among many.
\end{attention}

\begin{popfigure}
\begin{tikzpicture}[scale=1.0, every node/.style={font=\small, align=center}]
\node[fill=popblueL, draw=popblue, rounded corners, minimum width=3.2cm, align=center] (high) at (0,2.2){blood glucose HIGH\\(after meal)};
\node[fill=popgreenL, draw=popgreen, rounded corners, minimum width=3.2cm, align=center] (beta) at (4.6,2.2){$\beta$-cells of pancreas\\secrete \textbf{insulin}};
\node[fill=popgreenL, draw=popgreen, rounded corners, minimum width=3.4cm, align=center] (cells) at (4.6,-0.6){liver \& muscle cells\\take up / store glucose};
\node[fill=popblueL, draw=popblue, rounded corners, minimum width=3.2cm, align=center] (norm) at (0,-0.6){blood glucose returns\\to set point};
\draw[->, thick, popink] (high) -- (beta);
\draw[->, thick, popink] (beta) -- (cells);
\draw[->, thick, popink] (cells) -- (norm);
\draw[->, thick, poppurple] (norm.west) .. controls (-2.4,-0.6) and (-2.4,2.2) .. (high.west) node[midway, left, text=poppurple, align=center]{negative\\feedback};
\node[text=poporange] at (4.6,3.5) {\textbf{Type 1:} $\beta$-cells destroyed $\rightarrow$ no insulin};
\node[text=poporange] at (4.6,-1.9) {\textbf{Type 2:} target cells resist insulin};
\end{tikzpicture}
\legende{Negative-feedback control of blood glucose, showing where type 1 and type 2 diabetes break the loop.}
\end{popfigure}

\courssub{Thyroid, growth-hormone and adrenal disorders}

\textbf{Thyroid disorders.} The thyroid needs \cle{iodine} to make thyroxine. In iodine deficiency the gland enlarges as it works harder to capture iodine, producing \cle{goitre}; endemic goitre has historically occurred in some Cameroonian highland areas where soils and diets are poor in iodine — the reason table salt is iodised. \textbf{Hypothyroidism} slows metabolism (fatigue, weight gain, cold intolerance); in children, severe deficiency causes \cle{cretinism} (congenital hypothyroidism: stunted growth and mental impairment). \textbf{Hyperthyroidism} accelerates metabolism (weight loss, rapid heartbeat, bulging eyes — exophthalmic goitre).

\textbf{Growth-hormone disorders} (pituitary): hyper-secretion in childhood causes \cle{gigantism}, in adulthood \cle{acromegaly} (enlarged hands, feet, jaw); hypo-secretion in childhood causes pituitary \cle{dwarfism}.

\textbf{Adrenal disorders} (briefly): \cle{Addison's disease} is hypo-secretion of adrenal cortex hormones (weakness, low blood pressure, darkened skin); \cle{Cushing's syndrome} is hyper-secretion of cortisol (weight gain, "moon face", high blood pressure).

\begin{savaistu}[Iodised salt and goitre in Cameroon]
Iodine deficiency disorders were once a serious public-health problem in parts of the Adamawa, North-West and West highlands of Cameroon, where iodine-poor soils produce iodine-poor food crops. The national policy of \cle{salt iodisation} has greatly reduced goitre and, above all, protects unborn babies and children from cretinism. Checking that household salt is iodised is a simple, concrete act of NCD prevention.
\end{savaistu}

\courssub{Linking the two control systems}

The nervous and endocrine systems meet at the \cle{hypothalamus–pituitary axis}: the hypothalamus (nervous tissue) monitors the blood and directs the pituitary (the "master gland"), which in turn controls the thyroid, adrenals and gonads. Failure at any level of this axis produces disease — for example, lack of ADH (made in the hypothalamus, released by the posterior pituitary) causes diabetes insipidus. This is why the two systems are studied together.

\begin{aretenir}[Key points]
\begin{itemize}
\item Nervous disorders are located by mechanism: electrical (epilepsy), vascular (stroke), degenerative (Parkinson's, Alzheimer's), autoimmune (multiple sclerosis).
\item Stroke is an emergency linked to hypertension; think \textbf{FAST}. Epilepsy is managed by protecting the patient, never restraining, never putting objects in the mouth.
\item Endocrine diseases = hypo- or hyper-secretion breaking a negative-feedback loop.
\item Diabetes mellitus: type 1 = insulin deficiency (autoimmune); type 2 = insulin resistance (lifestyle-linked). Complications: retinopathy, neuropathy, nephropathy.
\item Goitre (iodine deficiency), cretinism, gigantism/acromegaly/dwarfism, Addison's and Cushing's complete the picture.
\end{itemize}
\end{aretenir}

\begin{exercice}[Glucose curves and feedback]
The table gives blood glucose (mmol/L) after a glucose drink:
\begin{center}
\begin{tabular}{|l|c|c|c|c|}
\hline
\rowcolor{popblueL} Time (min) & 0 & 30 & 60 & 120 \\
\hline
Person X & 4.8 & 8.2 & 6.4 & 5.0 \\
\hline
Person Y & 7.9 & 13.0 & 11.8 & 9.6 \\
\hline
\end{tabular}
\end{center}
\begin{enumerate}
\item Which person is diabetic? Give two pieces of evidence.
\item Explain, using negative feedback, why person X's glucose falls after 30 minutes.
\item State one long-term complication person Y risks if untreated.
\end{enumerate}
\end{exercice}

\begin{corrige}[Exercise 1]
\begin{enumerate}
\item Person Y is diabetic. Evidence: fasting glucose is already high (7.9 mmol/L); the peak (13.0 mmol/L) is far above normal and exceeds the renal threshold (about 10 mmol/L); glucose remains elevated (9.6 mmol/L) after 120 minutes.
\item In person X, rising blood glucose stimulates the $\beta$-cells of the islets of Langerhans to secrete insulin; insulin promotes glucose uptake by liver and muscle cells and its storage as glycogen, so blood glucose falls back to the set point — a negative-feedback response.
\item Any one of: retinopathy (blindness), neuropathy (nerve damage/foot ulcers), nephropathy (kidney damage), or increased risk of cardiovascular disease.
\end{enumerate}
\end{corrige}

\coursec{Section 3: Diseases of the renal and cardiovascular systems}

In Section 2 we saw how diabetes mellitus, an endocrine disorder, damages small blood vessels throughout the body. Two systems suffer most from this silent vascular damage: the \cle{renal system}, which filters the blood, and the \cle{cardiovascular system}, which pumps and carries it. The two are so closely linked that disease in one almost always provokes disease in the other. This section studies the main non-infectious diseases of the renal and cardiovascular systems, which together form the leading causes of adult death and disability in Cameroon today.

\courssub{The nephron: structural basis of renal disease}

Each kidney contains about one million microscopic filtering units called \cle{nephrons}. Understanding where a nephron can fail is the key to understanding every renal disease.

\begin{definition}[The nephron]
The \cle{nephron} is the functional unit of the kidney. It consists of:
\begin{itemize}
\item the \cle{glomerulus}, a ball of capillaries where blood is filtered under pressure;
\item \cle{Bowman's capsule}, a cup surrounding the glomerulus that collects the filtrate;
\item the \cle{proximal convoluted tubule}, the \cle{loop of Henl\'e} and the \cle{distal convoluted tubule}, where useful substances (all glucose, most water, salts) are reabsorbed into the blood;
\item the \cle{collecting duct}, where final water reabsorption (under ADH control) concentrates the urine.
\end{itemize}
\end{definition}

The logic is simple: the glomerulus is a \emph{filter}, the tubules are \emph{reclaimers}. If the filter is damaged, large molecules (proteins) and even cells (red blood cells) leak into the urine. If the tubules fail, useful substances such as glucose are lost and water balance collapses. Keep this filter--reclaimer picture in mind: it explains every urine test result you will meet in the examination.

\begin{popfigure}
\begin{center}
\begin{tikzpicture}[scale=0.95, every node/.style={font=\small}]
% Bowman's capsule and glomerulus
\draw[thick, popblue] (0,0) circle (1.1);
\draw[thick, poppink, fill=poppink!15] (0,0.15) circle (0.55);
\node[popdark] at (0,0.15) {glomerulus};
\node[popblue, align=center] at (0,-1.55) {Bowman's\\capsule};
% afferent / efferent arterioles
\draw[thick, poppink] (-2.2,0.6) -- (-0.6,0.35) node[midway, above, popdark]{afferent arteriole};
\draw[thick, poppink] (0.6,0.45) -- (2.2,0.8) node[midway, above, popdark]{efferent arteriole};
% proximal tubule
\draw[thick, popgreen] (1.0,-0.5) .. controls (2.2,-0.9) and (1.2,-1.6) .. (2.4,-1.8) .. controls (3.4,-2.0) and (2.6,-2.6) .. (3.2,-2.9);
\node[popgreen!60!black, align=center] at (4.6,-2.2) {proximal\\tubule};
% loop of Henle
\draw[thick, poporange] (3.2,-2.9) -- (3.2,-5.2) .. controls (3.2,-5.9) and (4.2,-5.9) .. (4.2,-5.2) -- (4.2,-2.9);
\node[poporange!80!black, align=center] at (3.7,-6.35) {loop of\\Henl\'e};
% distal tubule
\draw[thick, poppurple] (4.2,-2.9) .. controls (5.0,-2.5) and (4.4,-1.8) .. (5.2,-1.6);
\node[poppurple, align=center] at (6.3,-2.3) {distal\\tubule};
% collecting duct
\draw[thick, popblue] (5.2,-1.6) -- (5.2,-5.4);
\node[popblue, align=center] at (5.2,-5.85) {collecting duct};
% disorder labels
\draw[thick, popdark, ->] (-0.9,1.0) -- (-2.6,1.9);
\node[popdark, align=center, fill=poppink!20, rounded corners] at (-3.4,2.2) {\textbf{1} glomerulonephritis\\(inflamed filter)};
\draw[thick, popdark, ->] (5.2,-4.4) -- (7.0,-4.6);
\node[popdark, align=center, fill=popgold!25, rounded corners] at (8.3,-4.6) {\textbf{2} stone blockage\\(renal calculus)};
\draw[thick, popdark, ->] (2.9,-2.4) -- (1.4,-3.6);
\node[popdark, align=center, fill=popgreenL, rounded corners] at (0.6,-4.0) {\textbf{3} tubular failure\\(poor reabsorption)};
\end{tikzpicture}
\end{center}
\legende{The nephron and the main sites of non-infectious renal disorders: (1) glomerulus, (2) urinary tract downstream, (3) tubules.}
\end{popfigure}

\courssub{Diseases of the kidneys and urinary tract}

\textbf{Nephritis (glomerulonephritis)}\par
\cle{Glomerulonephritis} is an inflammation of the glomeruli. It often appears one to three weeks after a streptococcal infection (for example a poorly treated sore throat): antibodies produced against the bacterium form immune complexes that lodge in the glomerular membrane and inflame it (\emph{post-streptococcal} glomerulonephritis). It may also be \emph{autoimmune}, the immune system attacking the kidney's own tissue. The inflamed filter becomes leaky, so the characteristic signs are \cle{proteinuria} (protein in urine, making it foamy), \cle{haematuria} (blood in urine) and \cle{oedema} (swelling of the face and ankles, because lost plasma protein lowers the osmotic pressure of the blood, so water is retained less effectively in the capillaries and accumulates in the tissues).

\textbf{Kidney stones (renal calculi)}\par
When urine becomes too concentrated, dissolved mineral salts (calcium oxalate, uric acid) crystallise into hard stones. Risk factors include \emph{chronic dehydration} — very relevant in the hot Far North of Cameroon, where water intake may be low — and a diet rich in salt and animal protein. A stone stuck in the ureter causes intense \cle{flank pain} (renal colic), often with blood in the urine. Small stones are flushed out by drinking abundantly; large ones may be broken by shock waves (lithotripsy) or removed surgically. Prevention is simple and cheap: \emph{drink enough water to keep the urine dilute}.

\textbf{Renal failure}\par

\begin{definition}[Renal failure]
\cle{Renal failure} is the inability of the kidneys to filter the blood adequately. \emph{Acute} renal failure appears suddenly (severe infection, poisoning, shock) and is often reversible. \emph{Chronic} renal failure develops slowly over years, most commonly as a consequence of \emph{hypertension} and \emph{diabetes mellitus}, and is irreversible. Waste products such as urea accumulate in the blood (\cle{uraemia}), with dangerous fluid and electrolyte (especially potassium) imbalance.
\end{definition}

When both kidneys fail, two options exist. \cle{Haemodialysis} replaces filtration artificially: the patient's blood is pumped through a machine where it flows along one side of a \emph{semi-permeable membrane}, while a clean fluid, the \cle{dialysate}, flows on the other side. Urea and excess salts diffuse from blood (high concentration) into the dialysate (low concentration) down their concentration gradients; blood cells and proteins, too large, are retained. A kidney \cle{transplant} is the only true cure, but requires a compatible donor and lifelong immunosuppressive drugs. In Cameroon, dialysis sessions are expensive and dialysis centres are concentrated in big cities such as Yaound\'e and Douala, so many patients cannot attend the two to three weekly sessions needed — a major public-health challenge.

\begin{popfigure}
\begin{center}
\begin{tikzpicture}[scale=0.9, every node/.style={font=\small}]
% membrane
\draw[very thick, popdark, dashed] (0,-3) -- (0,3) node[above, align=center]{semi-permeable\\membrane};
% blood side
\fill[popink!8] (-4.5,-3) rectangle (0,3);
\draw[very thick, poppink, ->] (-4.2,2.2) -- (-0.4,2.2) node[midway, above, popdark]{blood from patient (rich in urea)};
\draw[very thick, poppink, ->] (-0.4,-2.2) -- (-4.2,-2.2) node[midway, below, popdark]{purified blood returned};
\node[popdark] at (-2.3,0.9) {blood side};
\node[poppink] at (-2.3,0.3) {urea $\bullet$ \; salts $\bullet$};
\node[popmuted] at (-2.3,-0.5) {cells, proteins: retained};
% dialysate side
\fill[popblue!8] (0,-3) rectangle (4.5,3);
\draw[very thick, popblue, ->] (4.2,-1.4) -- (0.4,-1.4) node[midway, below, popdark]{used dialysate out};
\draw[very thick, popblue, ->] (0.4,1.4) -- (4.2,1.4) node[midway, above, popdark]{fresh dialysate in};
\node[popdark] at (2.3,0.5) {dialysate side};
% diffusion arrows
\draw[thick, poporange, ->] (-0.9,0.1) -- (0.9,0.1) node[midway, above, poporange!80!black]{urea diffuses};
\draw[thick, poporange, ->] (-0.9,-0.9) -- (0.9,-0.9) node[midway, below, poporange!80!black]{excess salts};
\end{tikzpicture}
\end{center}
\legende{Principle of haemodialysis: waste diffuses from blood to dialysate across a semi-permeable membrane; cells and proteins cannot cross.}
\end{popfigure}

\begin{attention}[Dialysis is not a cure]
A very common mistake is to write that dialysis \emph{cures} kidney failure. It does not: it only \emph{replaces} the filtration function temporarily and must be repeated two to three times per week for life (or until a transplant). It also does not replace the kidney's hormonal roles (e.g. erythropoietin production, which stimulates red blood cell formation — hence the anaemia of dialysis patients). Only a transplant restores full kidney function.
\end{attention}

\textbf{Urinary tract disorders (non-infectious aspects)}\par
\cle{Cystitis}, inflammation of the bladder, is usually infectious, but it can also result from irritation by stones or chemicals; it causes frequent, painful urination. \cle{Urinary incontinence} — involuntary leakage of urine — is non-infectious: it results from weakened pelvic floor muscles (after childbirth, in old age) or from nerve damage that prevents proper control of the bladder sphincter.

\begin{center}
\begin{tabularx}{\linewidth}{|l|X|X|}
\hline
\rowcolor{popblueL} \textbf{Disorder} & \textbf{Site / mechanism} & \textbf{Key signs} \\
\hline
Glomerulonephritis & Inflamed, leaky glomerular filter & Proteinuria, haematuria, oedema \\
\hline
Kidney stones & Crystallised salts blocking urinary tract & Renal colic, haematuria \\
\hline
Chronic renal failure & Progressive nephron loss (hypertension, diabetes) & Uraemia, fluid and electrolyte imbalance \\
\hline
Urinary incontinence & Weak sphincter or nerve damage & Involuntary urine leakage \\
\hline
\end{tabularx}
\end{center}

\courssub{Hypertension: the silent killer}

Blood pressure is the force exerted by blood on artery walls. It is expressed by two numbers: the \cle{systolic} pressure (when the ventricles contract) and the \cle{diastolic} pressure (when the heart relaxes), measured in millimetres of mercury (mmHg).

\begin{definition}[Hypertension]
\cle{Hypertension} is a persistently raised arterial blood pressure. The traditional diagnostic threshold used in most clinical and examination contexts is $\geq 140/90\ \mathrm{mmHg}$ on repeated measurements; more recent international guidelines already classify $130$--$139/80$--$89\ \mathrm{mmHg}$ as stage 1 hypertension. It is called the ``silent killer'' because it usually causes no symptoms for years while silently damaging the heart, brain, kidneys and eyes. It is extremely common among Cameroonian adults and is frequently discovered only after a stroke or heart failure.
\end{definition}

Risk factors include a \emph{salt-rich diet} (think of heavily salted stock cubes and smoked fish), obesity, physical inactivity, chronic stress, excessive alcohol, tobacco and heredity. Consequences include heart failure, stroke, kidney damage and retinopathy (damage to the retina that can lead to blindness).

\begin{methode}[Reading and classifying a blood-pressure value]
A value written $140/90\ \mathrm{mmHg}$ means: systolic $= 140$, diastolic $= 90$.
\begin{enumerate}
\item Identify the systolic (first, higher number) and the diastolic (second number).
\item Compare each with the thresholds: normal $< 120/80$; elevated $120$--$129/<80$; hypertension stage 1: $130$--$139$ or $80$--$89$; stage 2: $\geq 140$ or $\geq 90$.
\item Classify by the \emph{higher} of the two categories. Example: $145/85$ is stage 2 (because of the systolic value).
\item A single high reading is not enough: diagnosis requires repeated measurements on different occasions.
\end{enumerate}
\end{methode}

\begin{popfigure}
\begin{center}
\begin{tikzpicture}
\begin{axis}[
ybar, bar width=0.5cm,
width=12cm, height=6.5cm,
symbolic x coords={Normal, Elevated, Stage 1, Stage 2},
xtick=data,
ymin=0, ymax=160,
ylabel={Pressure (mmHg)},
legend style={at={(0.02,0.97)}, anchor=north west, font=\small},
nodes near coords, every node near coord/.append style={font=\scriptsize},
]
\addplot[fill=popblue] coordinates {(Normal,120) (Elevated,129) (Stage 1,139) (Stage 2,140)};
\addplot[fill=poporange] coordinates {(Normal,80) (Elevated,80) (Stage 1,89) (Stage 2,90)};
\legend{systolic threshold, diastolic threshold}
\end{axis}
\end{tikzpicture}
\end{center}
\legende{Upper limits of systolic and diastolic pressure for each blood-pressure category. Stage 2 values shown are the minimum thresholds ($\geq 140/90$).}
\end{popfigure}

\courssub{Atherosclerosis, coronary heart disease and heart failure}

\begin{definition}[Atherosclerosis and arteriosclerosis]
\cle{Atherosclerosis} is the deposition of fatty \cle{plaques} (mainly cholesterol) in the inner wall of arteries, narrowing the lumen. \cle{Arteriosclerosis} is the hardening and loss of elasticity of artery walls, often accompanying plaque formation. Both are promoted by a diet rich in saturated fats, smoking, diabetes and hypertension.
\end{definition}

A narrowed artery delivers less blood; a plaque can also rupture and trigger a clot (\emph{thrombus}) that blocks the vessel completely. The consequences depend on the organ supplied: heart (angina, infarction), brain (stroke), kidneys (renal failure).

\begin{popfigure}
\begin{center}
\begin{tikzpicture}[scale=1.0, every node/.style={font=\small}]
% healthy artery
\draw[thick, poppink, fill=poppink!15] (-4,0) circle (1.5);
\draw[thick, poppink, fill=white] (-4,0) circle (1.05);
\node[popdark] at (-4,0) {wide lumen};
\node[popdark, align=center] at (-4,-2.1) {\textbf{Healthy artery}\\elastic wall, free blood flow};
% atherosclerotic artery
\draw[thick, poppink, fill=poppink!15] (3,0) circle (1.5);
\draw[thick, popgold, fill=popgold!50] (3,0) ++(140:1.05) arc (140:-40:1.05) -- ++(-40:0.55) arc (-40:140:0.5) -- cycle;
\draw[thick, poppink, fill=white] (3,0) circle (0.5);
\node[popdark] at (3,0) {narrow};
\node[popdark, align=center] at (3,-2.1) {\textbf{Atherosclerotic artery}\\cholesterol plaque, reduced flow};
\draw[thick, popdark, ->] (4.3,0.9) -- (5.3,1.4);
\node[popdark, align=center] at (5.6,1.7) {plaque\\(cholesterol)};
\end{tikzpicture}
\end{center}
\legende{Cross-sections of a healthy artery (left) and an atherosclerotic artery (right): the plaque narrows the lumen and stiffens the wall.}
\end{popfigure}

\textbf{Coronary heart disease}\par
The heart muscle itself is fed by the \emph{coronary arteries}. Partial narrowing causes \cle{angina pectoris}: chest pain on effort, when oxygen demand exceeds supply. Complete blockage (usually by a clot on a plaque) causes a \cle{myocardial infarction} (heart attack): the deprived heart muscle dies. Symptoms include crushing central chest pain radiating to the left arm and jaw, sweating and breathlessness. The emergency response is to stop all activity, call for medical help immediately and, if available, give aspirin to limit clotting. Hospital treatment restores blood flow by a \cle{stent} (a small mesh tube keeping the artery open) or a coronary \cle{bypass} (a grafted vessel detouring around the blockage).

\begin{attention}[Heart attack $\neq$ cardiac arrest]
Do not confuse the two. A \emph{heart attack} (myocardial infarction) is a \textbf{circulation} problem: a blocked coronary artery; the heart usually keeps beating and the victim is conscious. A \emph{cardiac arrest} is an \textbf{electrical} problem: the heart stops pumping, the victim collapses and needs immediate cardiopulmonary resuscitation (CPR) and defibrillation. A heart attack can \emph{lead to} cardiac arrest, but the words are not interchangeable.
\end{attention}

\textbf{Heart failure, valves and congenital defects}\par
\cle{Heart failure} is the inability of the heart to pump enough blood for the body's needs — often the end result of hypertension or infarction — causing breathlessness and oedema of the legs. \emph{Valvular disorders} (narrowed or leaking valves) force the heart to work harder. \emph{Congenital defects} are present at birth, for example a ``hole in the heart'' (septal defect) allowing blood to shunt abnormally between chambers; many can be corrected surgically.

\textbf{Stroke, aneurysms and varicose veins}\par
A \cle{stroke} is the meeting point of cardiovascular and nervous disease (Section 2): a cerebral artery blocked by a clot or ruptured by hypertension deprives brain tissue of oxygen, causing paralysis or loss of speech. An \cle{aneurysm} is a balloon-like swelling of a weakened artery wall that may burst. \cle{Varicose veins} are dilated, tortuous leg veins caused by faulty valves, common in people who stand for long hours.

\begin{propriete}[The heart--kidney vicious circle]
Hypertension damages the small arteries of the kidney, reducing filtration; diseased kidneys retain salt and water and release hormones (renin) that raise blood pressure further. Thus \textbf{hypertension causes kidney disease, and kidney disease worsens hypertension} — a self-reinforcing vicious circle that explains why the two diseases so often occur together and why controlling blood pressure protects the kidneys (and vice versa). (Understanding, not memorisation, is examinable.)
\end{propriete}

\begin{exoresolu}[Interpreting a urine test and a blood-pressure reading]
A 52-year-old trader in Bamenda has a blood pressure of $155/95\ \mathrm{mmHg}$. His urine contains protein and traces of blood, but no glucose. (a) Classify his blood pressure. (b) Which region of the nephron is most likely damaged? Justify. (c) Explain why his ankles are swollen.
\tcblower
\textbf{(a)} Systolic $155 \geq 140$ and diastolic $95 \geq 90$: \textbf{hypertension stage 2}.\\
\textbf{(b)} Protein and blood cells are normally \emph{too large to be filtered}; their presence in urine means the \emph{filter} is leaky, so the damage is in the \textbf{glomerulus / Bowman's capsule} (glomerulonephritis), not the tubules. Absence of glucose confirms that tubular reabsorption still works.\\
\textbf{(c)} Loss of plasma protein in urine lowers the osmotic (oncotic) pressure of the blood, so water is retained less effectively in capillaries and accumulates in the tissues: \textbf{oedema} of the ankles. The chronic hypertension has probably damaged his glomeruli — an illustration of the heart--kidney vicious circle.
\end{exoresolu}

\begin{attention}[Exam tip: link urine composition to the nephron region]
In data-response questions, read the urine test logically: \emph{protein or blood} in urine $\rightarrow$ glomerulus (leaky filter); \emph{glucose} in urine $\rightarrow$ proximal tubule (failed reabsorption) or diabetes (blood glucose above the reabsorption threshold); \emph{very dilute, abundant urine} $\rightarrow$ collecting duct / ADH problem. Never write ``the kidney is damaged'' without naming the precise region.
\end{attention}

\courssub{Prevention of cardiovascular and renal disease}

Most of these diseases are preventable, and prevention is far cheaper than dialysis or bypass surgery. The key measures are: a diet \emph{low in salt and saturated fat} and rich in vegetables and fruits; \emph{regular physical exercise} (at least 30 minutes of brisk walking most days); \emph{no tobacco}; \emph{little or no alcohol}; maintaining a healthy body weight; and \emph{regular screening} — checking blood pressure and blood sugar even when feeling perfectly well, because hypertension and diabetes are silent. In Cameroon, community screening campaigns and reducing salt in cooked food are among the most cost-effective weapons against stroke, heart failure and kidney failure.

\begin{aretenir}[Key points of Section 3]
\begin{itemize}
\item The nephron = glomerular filter + tubular reclaimer; the site of damage determines the urine abnormalities.
\item Glomerulonephritis: leaky glomeruli $\rightarrow$ protein, blood in urine, oedema. Kidney stones: crystallised salts, prevented by hydration.
\item Chronic renal failure (mainly from hypertension and diabetes) is treated by dialysis or transplant; dialysis replaces but does not cure.
\item Hypertension (traditionally $\geq 140/90\ \mathrm{mmHg}$) is a silent killer; atherosclerosis narrows arteries with cholesterol plaques.
\item Angina and myocardial infarction result from coronary blockage; heart attack (circulation) differs from cardiac arrest (electrical).
\item Heart and kidney form a vicious circle; prevention rests on diet, exercise, no tobacco and regular screening.
\end{itemize}
\end{aretenir}

\begin{exercice}[Check your understanding]
\begin{enumerate}
\item A patient's urine contains glucose but no protein. Identify the likely nephron region affected and give one possible underlying disease.
\item Explain why a person living in a hot, dry region who drinks little water is at risk of kidney stones.
\item A man collapses with chest pain but remains conscious. Is this more likely a heart attack or a cardiac arrest? Justify and state the first-aid response.
\item Using the vicious-circle property, explain why treating hypertension slows the progression of chronic kidney disease.
\end{enumerate}
\end{exercice}

\begin{corrige}[Exercise 1]
\begin{enumerate}
\item Glucose is normally reabsorbed in the \textbf{proximal convoluted tubule}; its presence in urine (without protein) indicates tubular reabsorption failure, or more commonly \textbf{diabetes mellitus}, where blood glucose exceeds the tubules' reabsorption capacity. The glomerular filter is intact (no protein).
\item Low water intake produces \emph{concentrated urine}, in which dissolved mineral salts easily reach saturation and \emph{crystallise} into stones; heat increases water loss by sweating, worsening concentration.
\item A \textbf{heart attack} (myocardial infarction): the victim is conscious because the heart is still beating — the problem is a blocked coronary artery, not an electrical stop. Response: stop all activity, sit the person down, call for emergency medical help immediately, give aspirin if available and not contraindicated.
\item Lowering blood pressure reduces mechanical damage to the glomerular arterioles, slowing kidney damage; healthier kidneys in turn regulate salt, water and renin better, keeping blood pressure lower. Breaking the circle at any point benefits both organs.
\end{enumerate}
\end{corrige}

\coursec{Section 4: Disorders of the musculoskeletal system and integrated prevention of NCDs}

In Section 3 we saw how the kidneys and the cardiovascular system can fail without any germ being involved: hypertension, atherosclerosis, kidney failure are diseases of \emph{structure and regulation}, not of infection. We now complete our tour of the body with the system that gives us shape and movement: the \cle{musculoskeletal system}. A builder in Bamenda who can no longer straighten his knees, a grandmother in Buea who breaks her hip after a simple fall, a schoolboy in Yaoundé whose legs curve outwards — none of these people caught a microbe. Their bones, joints or muscles have failed. Understanding \emph{why} requires us first to recall how these structures are built and maintained.

\courssub{Review: bones, joints and skeletal muscle}

\textbf{Bones.} Bone is a living, vascularised tissue, not dead scaffolding. It consists of a matrix of \cle{collagen} fibres (giving tensile strength and flexibility) hardened by \cle{calcium phosphate} salts, mainly hydroxyapatite (giving compressive rigidity). Two cell types constantly remodel it: \cle{osteoblasts} deposit new bone matrix and mineralise it, while \cle{osteoclasts} resorb old bone by secreting acid and enzymes. In a healthy adult the two activities are coupled and balanced, so total bone mass is stable. Bone also acts as the body's major mineral \emph{store}: when blood calcium falls, parathyroid hormone stimulates osteoclasts to release calcium from bone into the blood. This dual role — structural support and mineral reservoir — explains most bone diseases.

\textbf{Joints.} Freely movable joints such as the knee, hip and shoulder are \cle{synovial joints}. The ends of the two bones are covered with smooth \cle{articular (hyaline) cartilage}, a slippery, shock-absorbing tissue with no blood supply of its own (it receives nutrients by diffusion from synovial fluid). The joint is enclosed in a fibrous capsule lined by the \cle{synovial membrane}, which secretes \cle{synovial fluid} that lubricates the joint and nourishes the cartilage. Ligaments (dense regular connective tissue) hold the bones together and stabilise the joint.

\textbf{Skeletal muscle.} Muscles are bundles of multinucleated fibres that contract when stimulated by motor neurones at the \cle{neuromuscular junction}, where the neurotransmitter \cle{acetylcholine} is released onto nicotinic receptors on the motor end plate. Tendons (dense regular connective tissue) attach muscle to bone. Movement therefore requires four healthy partners: bone, joint, muscle, and nerve — and a disorder can strike any of them.

\begin{popfigure}
\begin{tikzpicture}[scale=0.9, font=\small]
% Healthy joint
\begin{scope}[xshift=0cm]
\fill[popblueL] (0,1.6) -- (2.4,1.6) -- (2.4,0.9) .. controls (1.2,0.5) .. (0,0.9) -- cycle;
\fill[popgreenL] (0,0.55) .. controls (1.2,0.15) .. (2.4,0.55) -- (2.4,-0.6) -- (0,-0.6) -- cycle;
\draw[popdark, thick] (0,1.6) -- (2.4,1.6) -- (2.4,0.9) .. controls (1.2,0.5) .. (0,0.9) -- cycle;
\draw[popdark, thick] (0,0.55) .. controls (1.2,0.15) .. (2.4,0.55) -- (2.4,-0.6) -- (0,-0.6) -- cycle;
\draw[poporange, very thick] (0,0.9) .. controls (1.2,0.5) .. (2.4,0.9);
\draw[poporange, very thick] (0,0.55) .. controls (1.2,0.15) .. (2.4,0.55);
\node[align=center] at (1.2,2.0) {\textbf{Healthy}};
\node[align=center, text width=2.6cm] at (1.2,-1.1) {thick smooth cartilage};
\end{scope}
% Osteoarthritis
\begin{scope}[xshift=4.2cm]
\fill[popblueL] (0,1.6) -- (2.4,1.6) -- (2.4,0.9) .. controls (1.2,0.5) .. (0,0.9) -- cycle;
\fill[popgreenL] (0,0.55) .. controls (1.2,0.15) .. (2.4,0.55) -- (2.4,-0.6) -- (0,-0.6) -- cycle;
\draw[popdark, thick] (0,1.6) -- (2.4,1.6) -- (2.4,0.9) .. controls (1.2,0.5) .. (0,0.9) -- cycle;
\draw[popdark, thick] (0,0.55) .. controls (1.2,0.15) .. (2.4,0.55) -- (2.4,-0.6) -- (0,-0.6) -- cycle;
\draw[poporange, very thick, dashed] (0,0.82) .. controls (1.2,0.45) .. (2.4,0.82);
\draw[poporange, very thick, dashed] (0,0.6) .. controls (1.2,0.25) .. (2.4,0.6);
\node[align=center] at (1.2,2.0) {\textbf{Osteoarthritis}};
\node[align=center, text width=2.6cm] at (1.2,-1.1) {cartilage worn thin, bones rub};
\end{scope}
% Rheumatoid arthritis
\begin{scope}[xshift=8.4cm]
\fill[popblueL] (0,1.6) -- (2.4,1.6) -- (2.4,0.9) .. controls (1.2,0.5) .. (0,0.9) -- cycle;
\fill[popgreenL] (0,0.55) .. controls (1.2,0.15) .. (2.4,0.55) -- (2.4,-0.6) -- (0,-0.6) -- cycle;
\draw[popdark, thick] (0,1.6) -- (2.4,1.6) -- (2.4,0.9) .. controls (1.2,0.5) .. (0,0.9) -- cycle;
\draw[popdark, thick] (0,0.55) .. controls (1.2,0.15) .. (2.4,0.55) -- (2.4,-0.6) -- (0,-0.6) -- cycle;
\draw[poppink, very thick, decorate, decoration={coil, amplitude=1.5pt, segment length=4pt}] (0,0.95) .. controls (1.2,0.55) .. (2.4,0.95);
\draw[poppink, very thick, decorate, decoration={coil, amplitude=1.5pt, segment length=4pt}] (0,0.5) .. controls (1.2,0.1) .. (2.4,0.5);
\node[align=center] at (1.2,2.0) {\textbf{Rheumatoid}};
\node[align=center, text width=2.6cm] at (1.2,-1.1) {inflamed swollen membrane};
\end{scope}
\end{tikzpicture}
\legende{A synovial joint in three states. Healthy: thick articular cartilage (orange) separates the bone ends. Osteoarthritis: cartilage is worn away so bones rub directly. Rheumatoid arthritis: the synovial membrane (pink) is inflamed and swollen (pannus), invading the joint space and eroding cartilage and bone.}
\end{popfigure}

\courssub{Arthritis: osteoarthritis versus rheumatoid arthritis}

\begin{definition}[Arthritis]
\cle{Arthritis} is inflammation or degeneration of one or more joints, producing the cardinal signs of \textbf{pain}, \textbf{stiffness}, \textbf{swelling} and, in advanced cases, \textbf{deformity} and loss of function. The two commonest forms are \cle{osteoarthritis} and \cle{rheumatoid arthritis}.
\end{definition}

\textbf{Osteoarthritis (OA)} is the ``wear-and-tear'' arthritis. With age, repetitive heavy use, or previous injury, the articular cartilage — which has no blood supply and limited repair capacity — becomes thin, rough and fissured. Subchondral bone thickens (sclerosis) and osteophytes (bony spurs) form at joint margins. Bone ends rub directly against each other, causing pain that is characteristically \emph{worse with movement and weight-bearing, and relieved by rest}. It typically affects weight-bearing joints (knees, hips, lumbar and cervical spine) and the distal interphalangeal joints (Heberden's nodes) and base of thumb in older people. Risk factors include age, obesity (extra mechanical load and adipokine-mediated inflammation), previous joint injury, heavy manual labour and genetic predisposition.

\textbf{Rheumatoid arthritis (RA)} is a chronic systemic \cle{autoimmune} disease: the immune system mistakenly attacks the synovial membrane, which becomes infiltrated with lymphocytes and plasma cells, thickened (pannus formation) and invasive. The pannus erodes articular cartilage and underlying bone. It often begins in young or middle-aged adults (30--50 years), affects \emph{symmetrical} small joints (metacarpophalangeal, proximal interphalangeal, wrists, MTP joints), and causes prolonged \emph{morning stiffness lasting more than one hour} that eases with activity. Systemic features include fatigue, low-grade fever, weight loss, and extra-articular manifestations (rheumatoid nodules, vasculitis, lung fibrosis). Deformities of the fingers (ulnar deviation, swan-neck, boutonnière) are characteristic late signs.

\begin{propriete}[Degenerative versus autoimmune joint disease]
Osteoarthritis is a \textbf{degenerative} disease: the primary lesion is mechanical and biochemical breakdown of \emph{articular cartilage}; pain worsens with \emph{activity} and improves with rest; it is commoner in the elderly and in weight-bearing joints; systemic inflammation is absent. Rheumatoid arthritis is an \textbf{autoimmune inflammatory} disease: the primary lesion is immune-mediated inflammation of the \emph{synovial membrane} (synovitis); stiffness is worst after \emph{rest} (prolonged morning stiffness); it can occur at any adult age and affects small joints symmetrically; systemic inflammation is present (raised ESR, CRP, rheumatoid factor, anti-CCP antibodies). (The justification of this distinction is examinable as a comparison table.)
\end{propriete}

\textbf{Management.} Neither form can be cured, but both can be controlled to preserve function and quality of life:
\begin{itemize}
\item \textbf{OA:} Weight reduction (crucial for lower limbs), physiotherapy (quadriceps strengthening for knee OA), walking aids, analgesics (paracetamol), NSAIDs (topical or oral), intra-articular corticosteroid injections; in severe refractory cases, surgical joint replacement (total hip or knee arthroplasty).
\item \textbf{RA:} Early referral to a rheumatologist. Disease-modifying antirheumatic drugs (DMARDs) — methotrexate is the anchor drug — started as early as possible to induce remission; biologic agents (anti-TNF, anti-IL6, etc.) for inadequate responders; short-term glucocorticoids for flares; NSAIDs for symptom control; physiotherapy and occupational therapy; surgery for established deformities.
\end{itemize}

\courssub{Gout: crystals in the joint}

\cle{Gout} is a painful crystal-induced arthritis caused by deposition of monosodium \cle{urate} crystals in a joint, classically the first metatarsophalangeal joint (podagra). Urate is the end product of purine catabolism. Normally the \textbf{kidneys excrete} the majority of urate (the rest via gut). Gout arises when serum urate exceeds its solubility limit (~6.8 mg/dL or 400 µmol/L), leading to crystal precipitation. This occurs due to \emph{overproduction} (high purine diet: red meat, offal, certain fish, yeast extracts; high fructose intake; tumour lysis) or, more commonly, \emph{underexcretion} (chronic kidney disease, diuretics, alcohol — especially beer — which increases urate production and decreases excretion). The crystals trigger intense inflammation via the NLRP3 inflammasome: a red, hot, exquisitely tender joint, often waking the patient at night. Notice the link with Section 3: gout sits at the crossroads of \emph{diet}, \emph{renal function} and \emph{cardiovascular risk} (metabolic syndrome).

\textbf{Management:} Acute attack — NSAIDs (indomethacin, naproxen), colchicine (low dose), or intra-articular/systemic glucocorticoids. Long-term — urate-lowering therapy (allopurinol or febuxostat) once the acute attack settles, targeting serum urate < 360 µmol/L (6 mg/dL); lifestyle: hydration, limit alcohol and high-purine foods, weight loss, treat comorbidities (hypertension, CKD).

\begin{savaistu}[Gout in the Cameroonian context]
Gout is increasingly diagnosed in urban Cameroon with the nutrition transition towards more red meat, organ meats (``kpem'', tripe), sugary drinks and beer consumption. Traditional diets rich in plantains, vegetables and fish were protective. Screening for hyperuricaemia in hypertensive or diabetic patients in primary care can prevent gout and may slow CKD progression.
\end{savaistu}

\courssub{Osteoporosis: the silent thinning of bone}

\begin{definition}[Osteoporosis]
\cle{Osteoporosis} is a systemic skeletal disease characterised by low \textbf{bone mineral density (BMD)} and microarchitectural deterioration of bone tissue, with a consequent increase in bone fragility and susceptibility to \emph{fragility fractures} (fractures occurring from a fall from standing height or less). Typical sites are the \textbf{hip} (neck of femur), the \textbf{spine} (vertebral crush fractures causing loss of height, kyphosis, chronic back pain) and the \textbf{distal radius} (Colles' fracture).
\end{definition}

Recall that bone is constantly remodelled by osteoblasts and osteoclasts in basic multicellular units (BMUs). Osteoporosis results when \emph{resorption exceeds deposition} over many years, leading to negative bone balance. Risk factors are:

\begin{itemize}
\item \textbf{Age}: bone mass peaks around age 25--30 (peak bone mass) and declines thereafter (~0.5--1\% per year).
\item \textbf{Sex and menopause}: oestrogen normally restrains osteoclast activity and promotes osteoblast function; after menopause women lose bone rapidly (2--3\% per year for 5--10 years) — this is why osteoporosis is commoner in elderly women.
\item \textbf{Calcium and vitamin-D deficiency}: too little raw material and poor intestinal calcium absorption (vitamin D deficiency is common even in sunny climates due to indoor lifestyles, veiling, sunscreen, dark skin).
\item \textbf{Physical inactivity}: bone strengthens in response to mechanical load (Wolff's law); a sedentary life weakens it.
\item \textbf{Secondary causes}: glucocorticoid excess (Cushing's, therapy), hyperthyroidism, hyperparathyroidism, hypogonadism, malabsorption, chronic liver/kidney disease, rheumatoid arthritis.
\item Smoking and excess alcohol, which also accelerate bone loss.
\end{itemize}

\begin{popfigure}
\begin{tikzpicture}[scale=1.0, font=\small]
% Normal bone
\begin{scope}
\draw[popdark, very thick] (0,0) rectangle (4,2.4);
\foreach \x/\y in {0.5/0.5, 1.4/0.5, 2.3/0.5, 3.2/0.5, 0.95/1.2, 1.85/1.2, 2.75/1.2, 0.5/1.9, 1.4/1.9, 2.3/1.9, 3.2/1.9}{
  \draw[popmuted, fill=white] (\x,\y) circle (0.16);}
\node at (2,-0.4) {\textbf{Normal bone}};
\node[align=center, text width=3.6cm] at (2,2.8) {dense matrix, small spaces};
\end{scope}
% Osteoporotic bone
\begin{scope}[xshift=6cm]
\draw[popdark, very thick] (0,0) rectangle (4,2.4);
\draw[popmuted, fill=white] (0.8,0.6) ellipse (0.5 and 0.35);
\draw[popmuted, fill=white] (2.4,0.5) ellipse (0.6 and 0.3);
\draw[popmuted, fill=white] (3.4,1.5) ellipse (0.45 and 0.5);
\draw[popmuted, fill=white] (1.6,1.7) ellipse (0.55 and 0.45);
\draw[popmuted, fill=white] (0.6,1.9) ellipse (0.3 and 0.25);
\node at (2,-0.4) {\textbf{Osteoporotic bone}};
\node[align=center, text width=3.6cm] at (2,2.8) {thinned matrix, large spaces};
\end{scope}
\end{tikzpicture}
\legende{Normal versus osteoporotic bone (schematic cross-section of trabecular bone). The same outer cortical size hides a much weaker interior: fewer and thinner bony trabeculae with larger marrow spaces, so the bone fractures under loads it once tolerated.}
\end{popfigure}

\textbf{Diagnosis:} Dual-energy X-ray absorptiometry (DXA) scan at hip and lumbar spine. T-score $\le$ -2.5 SD = osteoporosis; -1 to -2.5 = osteopenia. \textbf{Prevention} must begin in youth: a diet rich in calcium (milk, yogurt, cheese, small fish eaten with bones like ``canda'', green leafy vegetables) and vitamin D, plus regular \cle{weight-bearing exercise} (walking, running, farm work, dancing, sport) and resistance training, which stimulate osteoblasts. Fall prevention in the elderly (home safety, vision correction, balance training). \textbf{Treatment:} Calcium and vitamin-D supplements; bisphosphonates (alendronate, risedronate) as first-line to inhibit osteoclasts; denosumab (anti-RANKL); teriparatide (recombinant PTH) for severe cases; hormone replacement therapy (HRT) for early postmenopausal women with indications.

\courssub{Rickets and osteomalacia: soft bones}

Where osteoporosis is a loss of \emph{amount} of normally mineralised bone, \cle{rickets} (in children) and \cle{osteomalacia} (in adults) are defects of \emph{mineralisation}: the collagen matrix (osteoid) is laid down but not hardened with calcium salts, because of \textbf{vitamin-D or calcium deficiency}. Vitamin D (cholecalciferol) is needed for active intestinal absorption of calcium and phosphate; it is obtained from the diet (oily fish, egg yolk, fortified foods) \emph{and} synthesised in the skin from 7-dehydrocholesterol under UV-B sunlight.

In children, whose bones are still growing at the growth plates, soft bones bend under body weight: \textbf{bow-legs (genu varum) or knock-knees (genu valgum)}, \textbf{rachitic rosary} (costochondral junction swelling), \textbf{Harrison's groove} (diaphragmatic pull on soft ribs), \textbf{craniotabes} (soft skull), delayed fontanelle closure, delayed dentition, and growth failure — this is rickets. In adults the same deficiency causes diffuse \textbf{bone pain} (often lower back, pelvis, thighs), \textbf{proximal muscle weakness} (difficulty climbing stairs, rising from a chair), and pseudofractures (Looser's zones) — this is osteomalacia. Prevention is simple and cheap: adequate safe sunlight exposure (15--30 minutes on arms/face several times a week; abundant in Cameroon, though heavily veiled infants, strictly indoor lifestyles, and dark-skinned individuals in high latitudes are at risk), and a diet containing calcium and vitamin D (fish, eggs, fortified flour/milk). Treatment: high-dose vitamin D (ergocalciferol or cholecalciferol) and calcium supplements.

\begin{attention}[Rickets is not osteoporosis]
A classic GCE error is to confuse \textbf{rickets} with \textbf{osteoporosis}. Rickets affects \emph{children}: the bone is \emph{poorly mineralised} (excess osteoid) and becomes \emph{soft and deformed} (bow-legs, rachitic rosary) because growth plates are active. Osteoporosis affects \emph{adults}, especially after menopause: the bone is normally mineralised but \emph{reduced in quantity} (thin cortices, lost trabeculae) and becomes \emph{fragile}, breaking rather than bending. Different cause (mineralisation defect vs. resorption excess), different age group, different radiology and biochemistry.
\end{attention}

\courssub{Muscular disorders}

\begin{definition}[Muscular dystrophy]
\cle{Muscular dystrophy} is a group of \textbf{inherited (genetic)} disorders characterised by progressive degeneration and necrosis of skeletal muscle fibres, replaced by fat and fibrous tissue, causing \textbf{progressive muscle wasting and weakness}. The commonest and most severe form, \cle{Duchenne muscular dystrophy (DMD)}, is caused by mutations in the \emph{DMD} gene on the X chromosome (Xp21) encoding \cle{dystrophin}, a protein linking the cytoskeleton to the extracellular matrix. It affects almost exclusively boys, with onset at 2--5 years (delayed walking, Gowers' sign, calf pseudohypertrophy), wheelchair dependence by early teens, and death from cardiomyopathy or respiratory failure in the 20s. Becker muscular dystrophy is a milder allelic variant.
\end{definition}

\begin{attention}[Muscular dystrophy is genetic, not contagious]
Never present muscular dystrophy as contagious, or as the result of laziness, poor feeding or lack of exercise. It is caused by a \textbf{faulty gene} present from conception; no one can ``catch'' it, and no amount of effort by the child could prevent it. Confusing a genetic disease with an infectious or lifestyle disease is a serious conceptual error in the examination. Genetic counselling for carrier mothers is essential.
\end{attention}

\textbf{Myasthenia gravis} is an \cle{autoimmune} disorder of the \emph{neuromuscular junction}: autoantibodies (mainly against the nicotinic acetylcholine receptor, AChR) block, cross-link and accelerate degradation of receptors, so nerve impulses no longer trigger contraction effectively. The hallmark is \emph{fatigable weakness} — muscles (often ocular: ptosis, diplopia; then bulbar: dysphagia, dysarthria; then limb proximal) weaken with repeated use and recover after rest. Diagnosis: positive anti-AChR or anti-MuSK antibodies, decremental response on repetitive nerve stimulation, improvement with edrophonium (Tensilon test). Treatment: acetylcholinesterase inhibitors (pyridostigmine), immunosuppression (corticosteroids, azathioprine, mycophenolate), thymectomy (if thymoma or generalised early-onset), plasma exchange/IVIG for crises.

\textbf{Cramps and sprains} are everyday muscular problems. A \cle{cramp} is a sudden, painful, involuntary tonic contraction, often linked to dehydration, electrolyte loss in sweat (think of athletes in the dry-season heat), muscle fatigue, or nerve compression. A \cle{sprain} is the stretching or tearing of ligaments around a joint, typically the lateral ankle (inversion injury).

\courssub{Back pain, herniated disc, posture and sports injuries}

The spine is a stack of vertebrae separated by \cle{intervertebral discs} — a tough outer \cle{annulus fibrosus} with a soft, gelatinous \cle{nucleus pulposus} that act as shock absorbers. Lifting a heavy load with a flexed spine (instead of bending the knees and keeping the back straight) generates huge compressive and shear forces that can tear the annulus so the nucleus protrudes: a \cle{slipped (herniated/prolapsed) disc}. The protrusion may compress a spinal nerve root (commonly L5 or S1), causing \cle{sciatica} — radiating pain, numbness or weakness down the leg. Chronic \textbf{non-specific low back pain} is also promoted by poor posture, weak core (abdominal and paraspinal) muscles, obesity, smoking (disc degeneration), and long hours sitting — increasingly common among students and office workers. Prevention: correct lifting technique (bend knees, keep load close, avoid twisting), good sitting posture (lumbar support, feet flat), regular core-strengthening exercise, weight control, and a medium-firm mattress.

\textbf{Sports injuries} include:
\begin{itemize}
\item \cle{Sprains}: ligament tears (grades I--III).
\item \cle{Strains}: muscle or tendon tears.
\item \cle{Dislocations}: bone forced out of its joint (shoulder, elbow, fingers) — visible deformity, loss of movement.
\item \cle{Fractures}: broken bone (simple/closed, compound/open with skin breach, greenstick in children, comminuted).
\end{itemize}

\begin{methode}[The RICE first-aid protocol for acute soft-tissue injuries (sprains, strains, contusions)]
\begin{enumerate}
\item \textbf{R — Rest}: stop the activity immediately; do not ``walk it off''. Immobilise the injured part (splint/sling if needed).
\item \textbf{I — Ice}: apply ice wrapped in a damp cloth (never directly on skin) for 15--20 minutes every 2--3 hours for the first 48--72 hours to reduce swelling, pain and metabolic rate.
\item \textbf{C — Compression}: apply a firm elastic bandage (not so tight as to cut off circulation — check capillary refill) to limit swelling and provide support.
\item \textbf{E — Elevation}: raise the injured limb above heart level whenever possible so venous and lymphatic fluid drains away from the injury.
\end{enumerate}
Then seek medical care, especially if a fracture or dislocation is suspected (deformity, inability to bear weight, bone visible, neurovascular compromise). \textbf{Never} try to reduce a dislocated joint yourself — this can damage nerves and vessels.
\end{methode}

\courssub{Congenital skeletal disorders}

Some skeletal defects are present at birth. \cle{Congenital talipes equinovarus (club foot)} is a foot fixed in equinus (plantarflexion), varus (inversion) and adduction; it is treated successfully by the \cle{Ponseti method} — gentle weekly manipulation and above-knee plaster casts in infancy, followed by a minor Achilles tenotomy and bracing (boots and bar) for several years. \cle{Spina bifida} (myelomeningocele) is a failure of the posterior vertebral arches to close over the spinal cord during neurulation (weeks 3--4); severe forms cause paralysis of the legs, sensory loss, and neurogenic bladder/bowel. Its frequency is strongly reduced (by up to 70\%) when the mother has adequate \cle{folic acid} (400 µg daily) before conception and during the first trimester — a powerful argument for preconception care, food fortification, and a varied diet (green leafy vegetables, beans, fortified flour) for women of child-bearing age.

\courssub{Synthesis: shared risk factors and the integrated prevention of NCDs}

Step back now and look at everything this chapter has covered — reproductive cancers, diabetes, hypertension, stroke, kidney failure, arthritis, osteoporosis, muscular disorders. Different organs, different symptoms, yet the \emph{same handful of modifiable risk factors} keeps reappearing:

\begin{itemize}
\item \textbf{Tobacco use} — cancers (cervix, prostate, lung), atherosclerosis, stroke, osteoporosis, poor wound healing;
\item \textbf{Harmful alcohol use} — hypertension, liver cirrhosis, gout, cancers, cardiomyopathy, fetal alcohol syndrome;
\item \textbf{Unhealthy diet} (excess salt, sugar, saturated/trans fats, red/processed meat; too little fruit, vegetables, fibre, calcium) — obesity, type 2 diabetes, hypertension, gout, osteoporosis, constipation/diverticular disease;
\item \textbf{Physical inactivity} — obesity, type 2 diabetes, hypertension, osteoporosis, back pain, colon cancer.
\end{itemize}

Because the risk factors are shared, prevention can be \emph{integrated}. The World Health Organization promotes cost-effective ``best buy'' interventions for low- and middle-income countries: raising excise taxes on tobacco and alcohol, reducing salt in processed foods (reformulation), eliminating industrial trans fats, public awareness campaigns (mass media, school curricula), and screening (blood-pressure checks, cervical smears/HPV testing, blood-glucose, cholesterol) in primary health centres. Vaccination against hepatitis B (liver cancer) and HPV (cervical cancer) is also a ``best buy''. Prevention is a shared responsibility:

\begin{itemize}
\item the \textbf{individual}: balanced diet (traditional Cameroonian dishes: more vegetables, less oil/salt/bouillon cubes), regular exercise (walking, farming, sport), no tobacco, moderate or no alcohol, attending screening;
\item the \textbf{family}: healthy cooking habits (steaming, grilling, less salt and palm oil), encouraging active play for children, supporting antenatal care and folic acid for expectant mothers;
\item the \textbf{school}: health education, compulsory physical education, healthy school meals, tobacco-free campuses, safe water and sanitation;
\item the \textbf{government}: taxation and regulation of tobacco/alcohol/unhealthy foods, equipping health centres for NCD screening and essential medicines, urban planning for walkable cities, food fortification (salt iodisation, flour folic acid/iron), vaccination programmes, and sustained public campaigns.
\end{itemize}

\begin{popfigure}
\begin{tikzpicture}[font=\small, node distance=0.4cm]
\node[draw=popdark, fill=popblueL, rounded corners, align=center, text width=4.2cm, thick] (risk) at (0,0) {\textbf{Shared modifiable risk factors}\\ tobacco, alcohol, poor diet, inactivity};
\node[draw=popdark, fill=popgreenL, rounded corners, align=center, text width=3.4cm] (rep) at (-5.2,2.2) {\textbf{Reproductive}\\ cancers, fibroids, BPH};
\node[draw=popdark, fill=popgreenL, rounded corners, align=center, text width=3.4cm] (nerv) at (0,3.4) {\textbf{Nervous / endocrine}\\ stroke, epilepsy, diabetes, goitre};
\node[draw=popdark, fill=popgreenL, rounded corners, align=center, text width=3.4cm] (card) at (5.2,2.2) {\textbf{Renal / cardiovascular}\\ hypertension, CKD, CHD};
\node[draw=popdark, fill=popgreenL, rounded corners, align=center, text width=3.4cm] (musc) at (0,-3.4) {\textbf{Musculoskeletal}\\ arthritis, osteoporosis, gout};
\node[draw=poporange, very thick, fill=white, rounded corners, align=center, text width=4.6cm] (prev) at (8.6,-1.2) {\textbf{Integrated prevention}\\ diet, exercise, screening, taxation, education, vaccination};
\draw[->, thick, popdark] (risk) -- (rep);
\draw[->, thick, popdark] (risk) -- (nerv);
\draw[->, thick, popdark] (risk) -- (card);
\draw[->, thick, popdark] (risk) -- (musc);
\draw[->, very thick, poporange] (risk.east) -- (prev.west);
\end{tikzpicture}
\legende{Concept map: the major non-infectious diseases of the six body systems share the same four modifiable risk factors, so a single integrated prevention strategy acts on all systems at once.}
\end{popfigure}

\begin{exoresolu}[Building a comparison table under exam conditions]
\textbf{Statement.} Construct a table comparing osteoarthritis and rheumatoid arthritis under the headings: cause, joints affected, pattern of pain/stiffness, and management. (8 marks)
\tcblower
\textbf{Solution.} A clear table with four rows earns marks efficiently:

\begin{center}
\begin{tabularx}{\linewidth}{|l|X|X|}
\hline
\rowcolor{popblueL} \textbf{Feature} & \textbf{Osteoarthritis} & \textbf{Rheumatoid arthritis} \\
\hline
Cause & Degeneration (wear and tear) of articular cartilage; age-related & Autoimmune inflammation of the synovial membrane \\
\hline
Joints affected & Weight-bearing joints (knee, hip, spine); often unilateral/asymmetrical & Small joints of hands and wrists; symmetrical \\
\hline
Pain / stiffness & Pain worse with activity, relieved by rest; brief morning stiffness (<30 min) & Morning stiffness >1 hour, eases with movement; pain at rest \\
\hline
Management & Weight loss, physiotherapy, analgesics/NSAIDs, joint replacement & DMARDs (methotrexate), biologics, glucocorticoids, physiotherapy \\
\hline
\end{tabularx}
\end{center}

Each correct, contrasted row earns 2 marks. Note the technique: identical headings down the side, one disease per column, and \emph{parallel} wording so the contrast is visible at a glance.
\end{exoresolu}

\begin{attention}[Exam tip: master the comparison-table format]
Structured GCE questions very often ask you to ``compare disorders X and Y with respect to cause, symptoms, diagnosis, treatment and prevention'' — for example osteoarthritis vs rheumatoid arthritis, rickets vs osteoporosis, or Type 1 vs Type 2 diabetes. Practise drawing the table skeleton (headings first) in under a minute, then filling one row at a time. Never write the two diseases in two separate essays: examiners award marks for \emph{point-by-point contrast}, which only a table (or paired statements) displays clearly.
\end{attention}

\begin{experience}[Practical: Assessing joint range of motion and RICE application]
\textbf{Objective:} Practise measuring active and passive range of motion (ROM) of the knee and ankle using a goniometer, and demonstrate the RICE protocol on a simulated ankle sprain.
\textbf{Materials:} Goniometer, ice packs, elastic bandages, pillows, consenting partner.
\textbf{Procedure:}
\begin{enumerate}
\item Measure active ROM: ask partner to flex/extend knee and dorsiflex/plantarflex ankle; record angles.
\item Measure passive ROM: you move the joint through its range; note any pain, crepitus, or end-feel (firm, soft, hard).
\item Simulate RICE: Rest (partner sits), Ice (wrapped pack on lateral ankle 15 min), Compression (figure-of-eight bandage), Elevation (leg on pillows > heart level).
\item Re-assess swelling and pain after 20 min.
\end{enumerate}
\textbf{Report:} Record initial and post-RICE measurements. Discuss why ice is wrapped (prevent ice burn) and why compression must not impair circulation (check capillary refill < 3 s).
\end{experience}

\begin{exercice}[Applying your knowledge]
\begin{enumerate}
\item A 65-year-old woman breaks her hip after slipping gently on a wet kitchen floor. Name the most likely underlying bone disease and give three risk factors she may have.
\item Explain why a person with poorly functioning kidneys (chronic kidney disease stage 3) has an increased risk of gout.
\item A footballer twists his ankle during a match in Douala. Outline the first aid you would give on the pitch.
\item State one difference and one similarity between rickets and osteomalacia.
\item A 10-year-old boy presents with progressive difficulty climbing stairs, frequent falls, and enlarged calf muscles. Name the most likely diagnosis, its inheritance pattern, and the deficient protein.
\end{enumerate}
\end{exercice}

\begin{corrige}[Exercise 1]
\begin{enumerate}
\item \textbf{Osteoporosis}: a fragility fracture (hip fracture from standing height fall) is its hallmark. Risk factors: age > 60, post-menopausal (loss of oestrogen), possible calcium/vitamin-D deficiency, physical inactivity, smoking, excess alcohol, low body weight, family history.
\item Uric acid is excreted primarily by the kidneys (glomerular filtration, tubular secretion/reabsorption). In CKD, reduced GFR and tubular dysfunction decrease urate excretion, so serum urate rises above saturation, leading to monosodium urate crystal deposition in joints and soft tissues, causing gout.
\item Apply \textbf{RICE} immediately: \textbf{R}est — stop play, do not weight-bear; \textbf{I}ce — apply ice wrapped in cloth for 15--20 min; \textbf{C}ompression — firm elastic bandage (figure-of-eight); \textbf{E}levation — raise leg on bag/bench above hip level. Refer to health facility for X-ray to exclude fracture (Ottawa ankle rules).
\item \emph{Difference}: rickets occurs in \emph{children} with open growth plates, causing \emph{bone deformities} (bow-legs, rachitic rosary); osteomalacia occurs in \emph{adults} after growth plate fusion, causing \emph{bone pain and proximal muscle weakness} without deformity. \emph{Similarity}: both are caused by vitamin-D and/or calcium deficiency leading to defective mineralisation of osteoid (excess unmineralised bone matrix).
\item \textbf{Duchenne muscular dystrophy (DMD)}; \textbf{X-linked recessive} inheritance (affects boys, carrier mothers); deficient protein: \textbf{dystrophin}.
\end{enumerate}
\end{corrige}

\begin{aretenir}[Key points of Section 4]
\begin{itemize}
\item Bone is living tissue remodelled by osteoblasts (formation) and osteoclasts (resorption); synovial joints rely on articular cartilage and synovial fluid; muscles contract via the neuromuscular junction (acetylcholine).
\item Osteoarthritis = degenerative cartilage loss (age, mechanical load) affecting weight-bearing joints; pain on movement. Rheumatoid arthritis = autoimmune synovitis affecting small joints symmetrically; prolonged morning stiffness, systemic features.
\item Gout = monosodium urate crystal deposition due to hyperuricaemia (overproduction/underexcretion); linked to diet, alcohol, CKD; acute podagra, chronic tophi.
\item Osteoporosis = low bone mass + microarchitectural deterioration $\rightarrow$ fragility fractures (hip, spine, wrist); risk: age, menopause, low Ca/Vit D, inactivity, glucocorticoids. Prevention: peak bone mass in youth, weight-bearing exercise, Ca/Vit D, fall prevention.
\item Rickets (children) / osteomalacia (adults) = defective mineralisation from Vit D/Ca deficiency $\rightarrow$ soft bones (deformity in children, pain/weakness in adults).
\item Muscular dystrophy = genetic (e.g., DMD: X-linked, dystrophin deficiency); myasthenia gravis = autoimmune (anti-AChR antibodies) $\rightarrow$ fatigable weakness.
\item Acute soft-tissue injury: RICE (Rest, Ice, Compression, Elevation). Never reduce a dislocation yourself.
\item Congenital: club foot (Ponseti method), spina bifida (folic acid prevention pre-conception).
\item All major NCDs share four modifiable risk factors (tobacco, alcohol, unhealthy diet, physical inactivity) $\rightarrow$ integrated prevention (WHO best buys) involving individual, family, school, government.
\end{itemize}
\end{aretenir}

\newpage

\coursec{Integration activity}

\courssub{Aims of the activity}

By the end of this integration activity, you should be able to:

\begin{itemize}
\item \textbf{Interpret clinical data} (blood pressure readings, urine tests, blood glucose values, physical examination findings) and link each abnormal result to the body system concerned;
\item \textbf{Propose and justify a diagnosis} of a non-infectious disease affecting the reproductive, nervous, endocrine, renal, cardiovascular or musculoskeletal system;
\item \textbf{Explain the underlying biological mechanism} of each disorder using correct Upper Sixth terminology;
\item \textbf{Suggest confirmatory tests} and realistic \textbf{treatment and lifestyle measures};
\item \textbf{Work as a team} to design a community prevention campaign targeting the shared, modifiable risk factors of non-communicable diseases (NCDs): excess salt, tobacco, alcohol, physical inactivity and unhealthy diet;
\item \textbf{Communicate scientific ideas} clearly in a short oral presentation, as expected in the GCE Advanced Level practical and problem-solving questions.
\end{itemize}

\begin{aretenir}[Why this matters]
Non-communicable diseases such as hypertension, diabetes, kidney failure, stroke, cancers and osteoporosis are now leading causes of illness and death among Cameroonian adults. Unlike malaria or cholera, they cannot be cured with a short course of drugs: they are \cle{prevented} and \cle{managed}. Understanding them makes you a genuine resource person in your family and quarter.
\end{aretenir}

\courssub{Situation}

\textbf{Health-screening day at the Lyc\'ee de Nkolbisson.}\par

As part of the third-term health education programme, your school has organised a fictional \emph{free community screening day} in the school hall. Volunteer nurses have measured blood pressure, tested urine with dipsticks, measured fasting blood glucose with a glucometer, and recorded simple examination findings for volunteer adults from the Nkolbisson neighbourhood. The results have been anonymised and compiled into four \textbf{case files}. Your class is divided into \textbf{community health teams} of four to six students. Each team acts as the medical analysis unit for the day: you must study the case files, propose diagnoses, and then design a prevention poster for the quarter.

\textbf{Case file 1 — Mrs A., 52 years old, market trader.}\par
Blood pressure measured twice, 15 minutes apart: \SI{165/105}{mmHg} both times. She reports headaches and swollen ankles in the evening. Urine dipstick: \textbf{protein +++} (marked proteinuria). She adds a stock cube (\emph{Maggi}) and extra salt to most dishes and does little physical exercise.

\textbf{Case file 2 — Mr B., 30 years old, moto-taxi driver.}\par
For three months he has complained of \textbf{excessive thirst} (polydipsia), \textbf{frequent and abundant urination} (polyuria) and unusual tiredness. Fasting blood glucose measured this morning: \SI{9.5}{mmol/L} (normal fasting value: \SIrange{3.9}{5.5}{mmol/L}; diabetes is suspected above \SI{7.0}{mmol/L} on two occasions). His father was diagnosed with diabetes at 45.

\textbf{Case file 3 — Mr C., 60 years old, retired teacher.}\par
He reports \textbf{difficulty starting urination}, a weak urine stream, and the need to urinate several times at night (nocturia). Digital rectal examination by the volunteer doctor: \textbf{enlarged prostate}, smooth and firm. No pain, no fever, urine dipstick negative for blood and nitrites.

\textbf{Case file 4 — Mrs D., 68 years old, farmer.}\par
She slipped on a wet floor at home — a \emph{minor} fall — and \textbf{fractured her hip}. X-ray and a bone density scan (DEXA) show \textbf{low bone mineral density} (T-score below $-2.5$). She went through menopause at 50, drinks little milk, and has become much less active with age.

\textbf{Reference data provided to all teams:}

\begin{center}
\begin{tabularx}{\linewidth}{|l|X|}
\hline
\rowcolor{popblueL} \textbf{Measurement} & \textbf{Normal / reference value} \\
\hline
Resting blood pressure & below $140/90\ \mathrm{mmHg}$ (ideal about $120/80$) \\
\hline
Fasting blood glucose & $3.9$--$5.5\ \mathrm{mmol\,L^{-1}}$; diabetes suspected $\geq 7.0\ \mathrm{mmol\,L^{-1}}$ twice \\
\hline
Urine protein (dipstick) & negative or trace \\
\hline
Prostate (rectal examination) & small, walnut-sized, smooth \\
\hline
Bone density (DEXA T-score) & above $-1.0$ normal; $-1.0$ to $-2.5$ osteopenia; below $-2.5$ osteoporosis \\
\hline
\end{tabularx}
\end{center}

\courssub{Tasks}

\textbf{Part A — The detective work (for each of the four case files):}

\begin{enumerate}
\item Identify the \textbf{body system} affected and the \textbf{organ(s)} involved.
\item Suggest the \textbf{most likely non-infectious disease}, justifying your answer with \emph{at least two pieces of evidence} taken from the case file and the reference table.
\item Explain the \textbf{underlying biological mechanism} of the disorder (for example: how the structure or regulation of the organ is disturbed, which hormone or tissue is involved, and how this produces the observed symptoms).
\item Propose \textbf{one or two further tests} that would confirm the diagnosis.
\item Recommend appropriate \textbf{treatment and lifestyle changes}, adapted to the patient's daily life in Nkolbisson.
\end{enumerate}

\textbf{Part B — Connecting the cases:}

\begin{enumerate}
\setcounter{enumi}{5}
\item For Case 1, explain how untreated hypertension can \textbf{damage the kidneys}, and how kidney damage in turn \textbf{raises blood pressure} (a vicious circle). Which hormone system, studied in the chapter on the endocrine and renal systems, is involved?
\item For Case 2, explain why high blood glucose causes \textbf{thirst and polyuria} (think of osmosis in the kidney tubules), and name \textbf{two long-term complications} of uncontrolled diabetes affecting other systems studied in this chapter.
\item Identify the \textbf{shared, modifiable risk factors} appearing across the four files (diet, salt, inactivity, age, heredity, hormonal changes). Classify them as \emph{modifiable} or \emph{non-modifiable}.
\end{enumerate}

\textbf{Part C — The community campaign:}

\begin{enumerate}
\setcounter{enumi}{8}
\item As a team, design a \textbf{one-page poster} entitled \emph{``Stop NCDs in our quarter''}. It must contain: (i) one clear message per shared risk factor (salt, tobacco, alcohol, inactivity); (ii) one simple diagram or pictogram; (iii) one slogan in language the neighbourhood understands.
\item Present your four diagnoses and your poster to the class in a \textbf{5-minute plenary}, and be ready to defend your reasoning against questions from the other teams.
\end{enumerate}

\begin{attention}[Common traps]
\begin{itemize}
\item Do not confuse \textbf{benign prostatic hyperplasia} with \textbf{prostate cancer}: a smooth, firm, painless enlargement in a 60-year-old suggests BPH, but only further tests (PSA blood test, biopsy if indicated) can exclude cancer.
\item A \emph{single} high blood glucose or blood pressure reading is not enough for a diagnosis: values must be \textbf{confirmed on a second occasion}.
\item Osteoporosis is \emph{not} an infectious or inflammatory joint disease: it is a loss of \textbf{bone mass} that weakens the skeleton silently until a fracture occurs.
\item Proteinuria is a sign of \textbf{kidney damage}, not a disease in itself.
\end{itemize}
\end{attention}

\courssub{Model answer}

\textbf{Case 1 — Mrs A. (hypertension with kidney damage).}\par
\emph{System:} cardiovascular system, with renal involvement. \emph{Diagnosis:} chronic \cle{hypertension} (BP $165/105\ \mathrm{mmHg}$, confirmed on two readings, well above $140/90$) with signs of \cle{chronic kidney disease} (proteinuria +++, ankle oedema). \emph{Mechanism:} persistently high pressure damages the walls of arterioles, including the glomerular capillaries; the glomerular filter becomes leaky, so plasma proteins pass into the urine (proteinuria). Damaged kidneys retain salt and water, increasing blood volume and activating the \textbf{renin--angiotensin--aldosterone system}, which raises blood pressure further — a vicious circle. Fluid retention explains the swollen ankles. \emph{Confirmatory tests:} serum creatinine and estimated glomerular filtration rate, 24-hour urine protein, ambulatory BP monitoring. \emph{Management:} antihypertensive drugs prescribed by a doctor (e.g.\ a diuretic or an ACE inhibitor, which also protects the kidney), drastic reduction of added salt and stock cubes, regular walking, weight control, and follow-up at the district hospital.

\textbf{Case 2 — Mr B. (diabetes mellitus).}\par
\emph{System:} endocrine system (pancreas / insulin). \emph{Diagnosis:} \cle{diabetes mellitus} — fasting glucose \SI{9.5}{mmol/L} exceeds the \SI{7.0}{mmol/L} threshold, with the classic symptoms of polydipsia and polyuria; family history supports the diagnosis. At age 30 this could be type 1 or early type 2; the distinction needs further tests. \emph{Mechanism:} insulin deficiency (type 1: autoimmune destruction of the $\beta$-cells of the islets of Langerhans) or insulin resistance (type 2) prevents glucose uptake by cells, so blood glucose rises. When the renal threshold is exceeded, glucose appears in the tubular fluid; by \textbf{osmosis} it holds water in the tubule, producing polyuria, dehydration and therefore intense thirst. \emph{Confirmatory tests:} repeat fasting glucose, glycated haemoglobin (HbA1c), urine test for glucose and ketones. \emph{Management:} medical treatment (insulin for type 1; oral antidiabetics for type 2), reduced intake of sugary drinks and refined carbohydrates, regular physical activity, and education on foot care and eye checks, since uncontrolled diabetes damages the \textbf{retina, kidneys, nerves and blood vessels}.

\textbf{Case 3 — Mr C. (benign prostatic hyperplasia).}\par
\emph{System:} male reproductive system (prostate gland), affecting urination. \emph{Diagnosis:} \cle{benign prostatic hyperplasia} (BPH) — age 60, smooth firm enlargement, obstructive urinary symptoms, no signs of infection. \emph{Mechanism:} under the continuing influence of androgens, the prostate tissue around the urethra proliferates with age; the enlarged gland \textbf{compresses the urethra}, so the bladder must contract harder, the stream is weak, emptying is incomplete and nocturia results. \emph{Confirmatory tests:} PSA blood test (to help exclude cancer), ultrasound of the prostate and bladder, urine flow measurement. \emph{Management:} drugs that relax the prostate smooth muscle or shrink the gland, surgery (e.g.\ transurethral resection) if obstruction is severe, reduced evening fluid and caffeine intake.

\textbf{Case 4 — Mrs D. (osteoporosis).}\par
\emph{System:} musculoskeletal system (bone tissue). \emph{Diagnosis:} \cle{osteoporosis} with a fragility fracture — a hip fracture after a minor fall plus a DEXA T-score below $-2.5$ is diagnostic. \emph{Mechanism:} bone is continuously remodelled by \textbf{osteoclasts} (resorption) and \textbf{osteoblasts} (deposition). After menopause the fall in \textbf{oestrogen} removes the brake on osteoclast activity, so resorption exceeds formation; bone mass and density fall, trabeculae thin, and the bone breaks under minimal trauma. Low calcium intake and inactivity accelerate the loss. \emph{Confirmatory tests:} already done (DEXA); blood calcium, phosphate and vitamin D complete the assessment. \emph{Management:} calcium- and vitamin-D-rich diet (milk, small fish eaten with bones, green vegetables), weight-bearing exercise adapted to her age, prescribed anti-resorptive drugs, and fall-proofing the home.

\textbf{Connecting the cases (Part B).}\par
Hypertension and kidney disease reinforce each other through the renin--angiotensin--aldosterone system; diabetes damages kidneys, eyes, nerves and vessels; all four patients share \emph{modifiable} risk factors (excess salt, sugary diet, inactivity, tobacco/alcohol where present) and \emph{non-modifiable} ones (age, sex, heredity, menopause). Prevention must therefore act on the modifiable factors, early, at community level.

\textbf{Poster campaign (Part C) — expected content.}\par
A good poster shows, for example, a heart, a kidney, a pancreas and a bone around the slogan \emph{``Less salt, more steps — our quarter, our health!''}, with four pictograms: a salt cellar crossed out, a cigarette crossed out, a bottle replaced by water, and a walking figure. Messages must be short, positive and actionable: \emph{cook with less Maggi}, \emph{walk 30 minutes a day}, \emph{check your blood pressure once a year}, \emph{screening saves lives}.

\begin{exercice}[Going further]
A fifth volunteer, aged 45, has a fasting blood glucose of \SI{6.2}{mmol/L} and a blood pressure of $138/88\ \mathrm{mmHg}$. Does he have diabetes or hypertension? What advice would your team give him, and why is this ``borderline'' stage actually the most important moment for prevention?
\end{exercice}

\begin{corrige}[Going further]
Neither diagnosis can be made: his glucose ($6.2$) is above the ideal range but below the $7.0\ \mathrm{mmol\,L^{-1}}$ threshold (this is \emph{impaired fasting glucose}, a pre-diabetic state), and his BP ($138/88$) is high-normal but below $140/90$. He should be re-tested to confirm the values, and advised to act \emph{now}: reduce sugar and salt, increase physical activity, control weight, avoid tobacco and excess alcohol, and attend yearly screening. This borderline stage is decisive because lifestyle change at this point can \textbf{prevent or delay} the onset of diabetes and hypertension, whereas once organ damage (kidney, eye, heart) has occurred it is largely irreversible. Prevention is cheaper and more effective than cure — the central message of the whole chapter.
\end{corrige}

\newpage

\coursec{Methods and worked examples}

\courssub{Methods and worked examples}

This part trains the practical skills the GCE examiner expects: reading clinical data (urine, blood, blood pressure), linking symptoms to the correct organ system, reasoning with hormone feedback loops, and building a prevention plan. Each method card is followed by a fully worked GCE-style exercise. Master these patterns and you will be able to transfer them to any unfamiliar case the examiner sets.

\begin{methode}[Identifying the affected organ system from a set of symptoms]
When a question gives a list of signs and symptoms and asks you to name the disorder or the system affected:
\begin{enumerate}
\item \textbf{List} each symptom and translate it into plain physiology (e.g. ``excessive thirst'' = water loss or high blood solute concentration).
\item \textbf{Group} the symptoms: do they point to one organ (kidney, pancreas, heart, joint) or to a whole system?
\item \textbf{Match} the group to a disorder you have studied, checking that \emph{every} symptom is explained by your answer.
\item \textbf{Justify} in your answer by quoting the key diagnostic symptom(s), not by listing everything.
\end{enumerate}
\emph{Reflex:} the examiner awards marks for the \emph{link} between symptom and organ function, not for the name of the disease alone.
\end{methode}

\begin{exoresolu}[Matching symptoms to a system]
A 52-year-old trader in Bamenda complains of persistent lower back pain, swelling of the ankles, reduced urine volume, and fatigue. Her urine is frothy. (a) State the organ system most likely affected. (b) Explain how each symptom supports your answer. (c) Suggest one further test to confirm the diagnosis.
\tcblower
\textbf{(a)} The \cle{renal (urinary) system} is affected; the picture suggests a kidney disorder such as chronic nephritis or chronic kidney disease.

\textbf{(b) Symptom-by-symptom justification:}
\begin{itemize}
\item \emph{Reduced urine volume:} damaged nephrons have a reduced filtration capacity, so less filtrate and therefore less urine is produced.
\item \emph{Frothy urine:} the glomerular filter is leaky and plasma \textbf{proteins} (albumin) pass into the urine (proteinuria); protein makes urine foam, like beaten egg white.
\item \emph{Swollen ankles (oedema):} loss of plasma protein lowers the osmotic (oncotic) pressure of the blood, so tissue fluid is not drawn back into the capillaries and accumulates in the tissues; gravity makes it collect at the ankles.
\item \emph{Fatigue:} failing kidneys produce less \textbf{erythropoietin}, reducing red blood cell formation (anaemia), and toxic wastes such as urea accumulate in the blood (uraemia); both cause tiredness.
\item \emph{Lower back pain:} the kidneys lie against the muscles of the lower back (loin region), so kidney inflammation is felt there.
\end{itemize}
\textbf{(c)} A \textbf{urinalysis} to confirm protein (and possibly blood) in the urine, or a \textbf{blood test} showing raised urea and creatinine levels, would confirm reduced kidney function.

\emph{Examiner's note:} full marks require each symptom to be tied to a mechanism, as above — a bare list scores poorly.
\end{exoresolu}

\begin{methode}[Interpreting urinalysis data]
To analyse a urine test table:
\begin{enumerate}
\item \textbf{Recall the normal values:} urine normally contains \emph{no} glucose, \emph{no} protein, \emph{no} blood cells; it does contain water, urea and salts.
\item \textbf{Compare} each patient's result with the normal values and flag abnormalities.
\item \textbf{Interpret each abnormality:} glucose $\rightarrow$ diabetes mellitus (insulin problem; blood glucose above the kidney's reabsorption threshold); protein or blood $\rightarrow$ damaged glomeruli (nephritis); very large volumes of dilute urine $\rightarrow$ failure of water reabsorption (ADH defect, i.e. diabetes insipidus).
\item \textbf{State the likely disorder and the defective structure} ($\beta$-cells of the islets of Langerhans, glomerulus, tubule/collecting duct, or the hypothalamus--posterior pituitary ADH axis).
\end{enumerate}
\end{methode}

\begin{exoresolu}[Analysing urine samples from three patients]
The table shows the results of urine tests on three patients at a district hospital in Buea.

\begin{center}
\begin{tabularx}{\linewidth}{|l|X|X|X|}\hline
\rowcolor{popblueL} \textbf{Substance} & \textbf{Patient A} & \textbf{Patient B} & \textbf{Patient C} \\ \hline
Glucose & Present & Absent & Absent \\ \hline
Protein & Absent & Present & Absent \\ \hline
Blood cells & Absent & Present & Absent \\ \hline
Volume per day & Very large & Normal & Very large, very dilute \\ \hline
\end{tabularx}
\end{center}

For each patient, state the likely disorder and the defective structure or hormone.
\tcblower
\textbf{Patient A — diabetes mellitus.} Glucose in the urine with normal protein and no blood means the glomerular filter is intact. The blood glucose concentration has exceeded the kidney tubules' maximum reabsorption capacity (the renal threshold), so glucose spills into the urine. The defect is a \textbf{failure of the $\beta$-cells of the islets of Langerhans to secrete enough insulin} (or of body cells to respond to it). The large urine volume follows because unreabsorbed glucose holds water in the tubules by osmosis (osmotic diuresis) — explaining the classic thirst and copious urination.

\textbf{Patient B — nephritis (glomerulonephritis).} Protein \emph{and} blood cells in urine indicate a \textbf{damaged glomerular basement membrane}: normally these molecules and cells are far too large to be filtered. Their presence shows the filter has become leaky, usually through inflammation, often following a bacterial infection.

\textbf{Patient C — diabetes insipidus.} A very large volume of very dilute urine, with no glucose or protein, points to a \textbf{deficiency of antidiuretic hormone (ADH)} from the posterior pituitary (or failure of the collecting ducts to respond to it). Without ADH the collecting duct walls remain impermeable to water, water is not reabsorbed, and huge volumes of dilute urine are excreted, causing intense thirst.

\emph{Key discriminator:} glucose in urine $\rightarrow$ think pancreas/insulin; protein or blood in urine $\rightarrow$ think glomerulus; dilute urine in large volume with no solute abnormality $\rightarrow$ think ADH.
\end{exoresolu}

\begin{methode}[Interpreting blood pressure readings and cardiovascular risk]
\begin{enumerate}
\item \textbf{Know the format:} a reading is written systolic/diastolic in mmHg, e.g. $120/80\ \mathrm{mmHg}$. Systolic = pressure during ventricular contraction; diastolic = pressure during ventricular relaxation.
\item \textbf{Classify:} roughly, normal is about $120/80\ \mathrm{mmHg}$; persistent readings at or above $140/90\ \mathrm{mmHg}$ indicate \cle{hypertension}.
\item \textbf{Explain consequences:} sustained high pressure damages artery walls, promotes atheroma (atherosclerosis), and overloads the heart $\rightarrow$ risk of stroke, heart failure, kidney damage.
\item \textbf{Link to risk factors:} high salt diet, obesity, smoking, stress, lack of exercise, age, family history.
\item \textbf{Recommend:} reduce salt and fat, exercise regularly, stop smoking, manage stress, take prescribed antihypertensive drugs.
\end{enumerate}
\end{methode}

\begin{exoresolu}[Blood pressure data over one month]
A 60-year-old man has his blood pressure measured weekly. Readings (mmHg): 155/98, 148/95, 160/102, 152/96. (a) State what the two numbers in each reading represent. (b) Comment on the data and name the condition. (c) Explain two possible long-term consequences if untreated. (d) Give three lifestyle changes you would advise.
\tcblower
\textbf{(a)} The first (higher) number is the \textbf{systolic pressure} — the pressure in the arteries when the ventricles contract and eject blood. The second (lower) number is the \textbf{diastolic pressure} — the pressure when the ventricles relax and the heart refills.

\textbf{(b)} Every reading has systolic pressure above $140\ \mathrm{mmHg}$ and diastolic pressure above $90\ \mathrm{mmHg}$. Because the elevation is \emph{persistent} (four separate occasions), this is \cle{hypertension} (high blood pressure), not a temporary rise due to stress or activity.

\textbf{(c) Long-term consequences (any two, explained):}
\begin{itemize}
\item \textbf{Atherosclerosis, heart attack and stroke:} high pressure damages the endothelium (inner lining) of arteries; cholesterol is deposited at the damage sites, forming atheroma plaques. A plaque or clot blocking a coronary artery causes a \textbf{heart attack (myocardial infarction)}; blocking or bursting a brain artery causes a \textbf{stroke}.
\item \textbf{Heart failure:} the left ventricle must pump against chronically high pressure; its muscle thickens, then weakens, and can no longer pump effectively.
\item \textbf{Kidney damage:} high pressure damages the delicate glomerular capillaries, progressively destroying nephrons (a cause of chronic kidney disease).
\end{itemize}

\textbf{(d) Lifestyle advice (any three):} reduce salt intake (less salt in cooking, avoid stock cubes in excess); take regular exercise such as brisk walking; stop smoking and reduce alcohol; reduce fatty/fried foods and maintain a healthy body weight; learn to manage stress.

\emph{Examiner's note:} ``persistent'' is the key word — one high reading is not a diagnosis.
\end{exoresolu}

\begin{methode}[Reasoning with hormone feedback loops (endocrine disorders)]
\begin{enumerate}
\item \textbf{Draw the loop} for the hormone concerned: gland $\rightarrow$ hormone $\rightarrow$ target organ $\rightarrow$ effect $\rightarrow$ negative feedback on the gland (or on the pituitary/hypothalamus that controls it).
\item \textbf{Locate the fault} described in the question: is the gland over-secreting, under-secreting, or is the target tissue unresponsive?
\item \textbf{Predict the consequences} by following the loop: if the gland fails, the hormone falls and its effects are lost; if the gland is overactive, effects are exaggerated and feedback cannot switch it off.
\item \textbf{Explain symptoms} as direct results of too much or too little hormone action.
\end{enumerate}
\emph{Key examples to master:} insulin (diabetes mellitus), thyroxine (goitre, hyper-/hypothyroidism), ADH (diabetes insipidus).
\end{methode}

\begin{popfigure}
\begin{center}
\begin{tikzpicture}[font=\small, >=stealth, node distance=1.1cm]
\node[draw=popblue, fill=popblueL, rounded corners, align=center, text width=2.6cm] (pit) at (0,0) {\textbf{Pituitary}\\ secretes TSH};
\node[draw=popgreen, fill=popgreenL, rounded corners, align=center, text width=2.6cm] (thy) at (5.2,0) {\textbf{Thyroid}\\ secretes thyroxine};
\node[draw=poporange, fill=white, rounded corners, align=center, text width=3.2cm] (body) at (5.2,-2.6) {\textbf{Body cells}\\ metabolic rate rises};
\draw[->, thick, popblue] (pit) -- node[above]{TSH stimulates} (thy);
\draw[->, thick, popgreen] (thy) -- node[right, align=center]{thyroxine\\ acts on} (body);
\draw[->, thick, popdark, dashed] (body.west) .. controls (0,-3.4) .. node[below, align=center]{high thyroxine inhibits pituitary\\ (negative feedback)} (pit.south);
\end{tikzpicture}
\end{center}
\legende{Negative feedback control of thyroxine secretion. If iodine is lacking, thyroxine falls, the dashed inhibition is removed, and excess TSH makes the thyroid swell (goitre).}
\end{popfigure}

\begin{exoresolu}[Goitre and the thyroxine feedback loop]
In some highland communities, iodine in the diet is scarce. (a) Name the hormone whose production requires iodine and the gland that produces it. (b) Using the idea of negative feedback, explain why iodine deficiency causes the gland to swell (goitre). (c) State two symptoms you would expect from the resulting hormone deficiency.
\tcblower
\textbf{(a)} The hormone is \textbf{thyroxine}, produced by the \textbf{thyroid gland} in the neck. Iodine is an essential component of the thyroxine molecule.

\textbf{(b) Feedback explanation, step by step:}
\begin{enumerate}
\item With little iodine available, the thyroid cannot make enough thyroxine, so the \textbf{blood thyroxine level falls}.
\item Normally, thyroxine exerts \textbf{negative feedback} on the pituitary, limiting its secretion of \textbf{thyroid-stimulating hormone (TSH)}.
\item With thyroxine low, this feedback is removed, so the pituitary secretes \textbf{excess TSH}.
\item TSH continuously stimulates the thyroid; the gland grows larger as its cells multiply and enlarge in a vain attempt to produce more thyroxine — this swelling of the neck is a \cle{goitre}.
\end{enumerate}

\textbf{(c) Symptoms of thyroxine deficiency (hypothyroidism):} tiredness and a low metabolic rate (feeling cold, weight gain), and slowed physical and mental activity; in children, deficiency impairs growth and mental development (cretinism).

\emph{Prevention link:} the use of \textbf{iodised table salt} — a cheap public-health measure promoted in Cameroon — prevents iodine-deficiency goitre.
\end{exoresolu}

\begin{methode}[Distinguishing disorders of the reproductive systems and advising on management]
\begin{enumerate}
\item \textbf{Classify the disorder} by cause: hormonal (e.g. infertility from low sperm count or failure of ovulation), structural (e.g. blocked Fallopian tubes, enlarged prostate), or cancerous (e.g. cancer of the cervix, breast, prostate).
\item \textbf{Link cause to effect:} blocked tubes prevent fertilisation; low testosterone reduces sperm production; an enlarged prostate compresses the urethra, causing difficulty in urination.
\item \textbf{State the key screening method} where relevant: Pap smear (cervical cancer), breast self-examination and mammography (breast cancer), PSA blood test and rectal examination (prostate).
\item \textbf{Advise} early detection, medical consultation, and a healthy lifestyle; stress that many reproductive cancers are treatable if caught early.
\end{enumerate}
\end{methode}

\begin{exoresolu}[Infertility and cancer screening]
(a) A couple has failed to conceive after three years of regular unprotected intercourse. Give two possible causes in the man and two in the woman, explaining each briefly. (b) Explain why regular screening is emphasised for cancers of the breast, cervix and prostate.
\tcblower
\textbf{(a) Possible causes of infertility:}

\emph{In the man:}
\begin{itemize}
\item \textbf{Low sperm count / abnormal sperm:} poor sperm production in the seminiferous tubules (due to hormonal imbalance, undescended testes, excessive heat around the scrotum, smoking or infection) means too few functional sperm reach the egg.
\item \textbf{Blocked vas deferens:} sperm are produced but cannot pass into the semen, often as a result of infection.
\end{itemize}

\emph{In the woman:}
\begin{itemize}
\item \textbf{Failure of ovulation:} hormonal imbalance (FSH/LH) means no ovum is released, so fertilisation cannot occur.
\item \textbf{Blocked Fallopian tubes:} often a result of pelvic inflammatory disease following an untreated sexually transmitted infection; the ovum and sperm cannot meet, and scarred tubes also raise the risk of ectopic pregnancy.
\end{itemize}

\textbf{(b) Why screening matters:} cancers of the breast, cervix and prostate often develop \textbf{silently}, with few symptoms in the early stages, yet they are among the commonest cancers in adults. Screening (breast self-examination and mammography; Pap smear for the cervix; PSA test and examination for the prostate) detects abnormal cells \emph{before} the cancer spreads. Early-stage cancers can be treated successfully by surgery, radiotherapy or chemotherapy; late-stage cancers that have metastasised (spread to other organs) are far harder to cure. Hence ``early detection saves lives'' is the central public-health message.
\end{exoresolu}

\begin{methode}[Analysing nervous system disorders from described deficits]
\begin{enumerate}
\item \textbf{Identify the level} of the problem: brain (stroke, epilepsy, Alzheimer's disease), spinal cord (paralysis below the injury), or peripheral nerves (neuritis).
\item \textbf{Use the deficit to locate the lesion:} paralysis of one side of the body after a burst or blocked brain artery $\rightarrow$ stroke affecting the opposite cerebral hemisphere; progressive loss of memory and reasoning in an elderly person $\rightarrow$ degenerative brain disease such as Alzheimer's; sudden recurrent seizures $\rightarrow$ epilepsy (abnormal electrical discharge in the brain).
\item \textbf{Explain the mechanism:} neurones cannot regenerate in the CNS, so damage is usually permanent; severity depends on the region and extent damaged.
\item \textbf{Conclude} with management: rehabilitation/physiotherapy, anticonvulsant drugs for epilepsy, control of hypertension to prevent stroke.
\end{enumerate}
\end{methode}

\begin{exoresolu}[A stroke case study]
A 65-year-old hypertensive man suddenly loses movement and sensation on the right side of his body and his speech becomes slurred. (a) Name the likely event and the organ affected. (b) Explain why the \emph{right} side of the body is affected. (c) Explain why recovery is often incomplete. (d) State two measures that could have reduced his risk.
\tcblower
\textbf{(a)} This is a \cle{stroke} (cerebrovascular accident): a region of the \textbf{brain} has been deprived of blood, either because an artery was blocked by a clot or because a weakened artery burst (haemorrhage). His hypertension is the major underlying cause.

\textbf{(b)} The motor and sensory pathways \textbf{cross over} to the opposite side on their way between the brain and the body, so the left cerebral hemisphere controls and receives sensation from the right side of the body. Damage to the left hemisphere therefore paralyses the \emph{right} side; the speech centres are also located in the left hemisphere in most people, explaining the slurred speech.

\textbf{(c)} Neurones of the central nervous system have virtually \textbf{no power of regeneration}. Once brain cells die from lack of oxygen, the lost functions can only be partly recovered as neighbouring areas are retrained through physiotherapy and speech therapy — hence recovery is usually slow and incomplete.

\textbf{(d) Prevention:} control blood pressure with medication and regular checks; reduce salt, stop smoking, exercise, and maintain a healthy weight — all of which reduce damage to cerebral arteries.
\end{exoresolu}

\begin{methode}[Interpreting joint and bone problems (musculoskeletal disorders)]
\begin{enumerate}
\item \textbf{Separate bone from joint problems:} bones $\rightarrow$ osteoporosis (loss of bone mass, fractures), rickets/osteomalacia (poor mineralisation from vitamin D or calcium deficiency); joints $\rightarrow$ arthritis (inflammation: osteoarthritis from wear of cartilage, rheumatoid arthritis from autoimmune attack), sprains and dislocations from injury.
\item \textbf{Match the patient profile:} elderly woman with a hip fracture after a minor fall $\rightarrow$ osteoporosis; elderly person with painful, stiff knees worsening with use $\rightarrow$ osteoarthritis; child with bowed legs and a poor diet $\rightarrow$ rickets.
\item \textbf{Explain the mechanism} at tissue level (cartilage worn away so bones rub; demineralised bone becomes porous and weak).
\item \textbf{Advise:} calcium- and vitamin D-rich diet, weight-bearing exercise, weight control, physiotherapy, and anti-inflammatory treatment as prescribed.
\end{enumerate}
\end{methode}

\begin{exoresolu}[Two elderly patients with joint pain]
Patient X, aged 70, has painful, stiff knee joints that worsen after walking and improve with rest; X-ray shows narrowed joint space. Patient Y, aged 68, broke her hip after a minor fall; bone scans show reduced bone density. (a) Name the most likely disorder in each case. (b) Explain the underlying cause of each. (c) Suggest one preventive measure for each condition.
\tcblower
\textbf{(a)} Patient X: \cle{osteoarthritis}. Patient Y: \cle{osteoporosis}.

\textbf{(b) Underlying causes:}
\begin{itemize}
\item \emph{Patient X:} the \textbf{cartilage} covering the ends of the bones in the knee joint has been worn away by years of use. Without this smooth, cushioning layer, the bone ends rub directly against each other, causing pain, stiffness and the reduced joint space seen on X-ray. Pain worsens with use and eases with rest — the classic pattern.
\item \emph{Patient Y:} bone is living tissue constantly being broken down and rebuilt. In old age — especially in women after the menopause, when protective oestrogen falls — breakdown exceeds rebuilding, so bone loses \textbf{calcium salts and mass}, becoming porous and fragile. Even a minor fall then causes a fracture, typically of the hip, wrist or spine.
\end{itemize}

\textbf{(c) Prevention:}
\begin{itemize}
\item \emph{Osteoarthritis:} maintain a healthy body weight to reduce the load on the knees; avoid repetitive joint strain; gentle regular exercise keeps joints mobile.
\item \emph{Osteoporosis:} a diet rich in \textbf{calcium and vitamin D} (milk, small fish eaten with the bones, green vegetables, sunlight exposure) plus regular weight-bearing exercise throughout life to build and maintain bone mass.
\end{itemize}

\emph{Common trap:} do not confuse osteo\emph{arthritis} (a joint/cartilage disease) with osteo\emph{porosis} (a bone-mass disease) — the GCE examiner tests this distinction directly.
\end{exoresolu}

\begin{methode}[Building an integrated NCD prevention plan]
For an essay or structured question on preventing non-infectious diseases (NCDs):
\begin{enumerate}
\item \textbf{Define} NCDs: diseases not caused by pathogens and not transmissible (cardiovascular diseases, diabetes, cancers, chronic respiratory and renal diseases, musculoskeletal disorders).
\item \textbf{Identify shared risk factors:} tobacco, harmful alcohol use, unhealthy diet (excess salt, sugar, fat), physical inactivity, obesity, plus non-modifiable factors (age, sex, heredity).
\item \textbf{Organise prevention at three levels:} individual (lifestyle choices), community (health education, screening campaigns, smoke-free spaces), and health system (affordable screening, treatment, follow-up).
\item \textbf{Give concrete examples:} iodised salt against goitre, blood pressure checks in churches and markets, vaccination against hepatitis B and HPV to prevent liver and cervical cancers, school sports against obesity.
\item \textbf{Conclude} that most NCDs are largely preventable because their main risk factors are behavioural.
\end{enumerate}
\end{methode}

\begin{exoresolu}[Essay: reducing the burden of NCDs in your community]
``Non-infectious diseases are an increasing cause of death among Cameroonian adults.'' Discuss the statement by (a) naming four NCDs and the organ system each affects, (b) explaining three shared, modifiable risk factors, and (c) proposing a realistic prevention plan for your community.
\tcblower
\textbf{(a) Four NCDs and their systems (examples):}
\begin{itemize}
\item \textbf{Hypertension / stroke} — cardiovascular and nervous systems.
\item \textbf{Diabetes mellitus} — endocrine system (pancreatic insulin defect) with renal and cardiovascular complications.
\item \textbf{Cancer of the cervix or prostate} — reproductive systems.
\item \textbf{Chronic kidney disease} — renal system.
\end{itemize}

\textbf{(b) Three shared modifiable risk factors, explained:}
\begin{enumerate}
\item \emph{Tobacco smoking:} damages artery walls (promoting atherosclerosis, hypertension, stroke) and introduces carcinogens (lung and other cancers).
\item \emph{Unhealthy diet:} excess \textbf{salt} raises blood pressure; excess \textbf{sugar and fat} promote obesity and type 2 diabetes; fatty diets accelerate atheroma formation.
\item \emph{Physical inactivity:} promotes obesity, raises blood pressure and blood glucose, and weakens the heart, bones and muscles — increasing risk across nearly all NCDs.
\end{enumerate}

\textbf{(c) Community prevention plan:}
\begin{itemize}
\item \emph{Health education:} talks in schools, churches, mosques and markets on salt reduction, a balanced diet (local vegetables, fruits, less fried food), and the dangers of tobacco and excessive alcohol.
\item \emph{Screening campaigns:} regular free blood-pressure and blood-sugar measurement days; Pap smear and breast examination drives for women; PSA testing for men over 50 — because early detection greatly improves survival.
\item \emph{Environment and policy:} encourage walking, community sports grounds, smoke-free public places, and promotion of iodised salt.
\item \emph{Health system:} ensure affordable drugs for hypertension and diabetes and proper follow-up of diagnosed patients, since NCDs require lifelong management.
\end{itemize}

\textbf{Conclusion:} NCDs share a small number of behavioural risk factors, so coordinated lifestyle change plus early screening can prevent a large proportion of cases — making prevention far cheaper than treatment.
\end{exoresolu}

\begin{attention}[Frequent errors to avoid in this chapter]
\begin{itemize}
\item Confusing \textbf{diabetes mellitus} (glucose in urine, insulin defect) with \textbf{diabetes insipidus} (dilute urine, ADH defect).
\item Writing that protein in urine indicates diabetes — it indicates \textbf{glomerular damage} (nephritis).
\item Calling any single high blood pressure reading ``hypertension'' — it must be \textbf{persistent}.
\item Mixing up \textbf{osteoarthritis} (worn joint cartilage) with \textbf{osteoporosis} (porous, fragile bone).
\item Describing NCDs as ``caught'' from other people — by definition they are \textbf{not transmissible}.
\item Naming a disease without \textbf{linking each symptom to the failing organ's function} — this is where marks are lost.
\end{itemize}
\end{attention}

\newpage

\coursec{Exercises}

\courssub{I check my knowledge}

\begin{exercice}[Multiple choice questions]
\textbf{1.} Which of the following is a non-infectious disease of the male reproductive system?
\begin{itemize}
\item[A] Prostatitis caused by \emph{Escherichia coli}
\item[B] Benign prostatic hyperplasia (BPH)
\item[C] Gonorrhoea
\item[D] Syphilis
\end{itemize}

\textbf{2.} A patient presents with polyuria, polydipsia and hyperglycaemia but no ketones in urine. The most likely diagnosis is:
\begin{itemize}
\item[A] Diabetes insipidus
\item[B] Type 1 diabetes mellitus
\item[C] Type 2 diabetes mellitus
\item[D] Cushing's syndrome
\end{itemize}

\textbf{3.} The glomerular filtration rate (GFR) is best estimated by the clearance of:
\begin{itemize}
\item[A] Urea
\item[B] Creatinine
\item[C] Glucose
\item[D] Sodium
\end{itemize}

\textbf{4.} Which joint is characteristically spared in rheumatoid arthritis but commonly affected in osteoarthritis?
\begin{itemize}
\item[A] Metacarpophalangeal joints
\item[B] Proximal interphalangeal joints
\item[C] Distal interphalangeal joints
\item[D] Wrist joints
\end{itemize}

\textbf{5.} Atherosclerosis primarily affects which layer of the arterial wall?
\begin{itemize}
\item[A] Tunica adventitia
\item[B] Tunica media
\item[C] Tunica intima
\item[D] External elastic lamina
\end{itemize}
\end{exercice}

\begin{exercice}[True/False with justification]
State whether each statement is true or false and justify your answer briefly.
\begin{enumerate}
\item Hypertension is defined as a sustained systolic blood pressure $\ge 140$ mmHg and/or diastolic $\ge 90$ mmHg.
\item In type 1 diabetes mellitus, insulin resistance is the primary defect.
\item Glomerulonephritis always presents with nephrotic syndrome (massive proteinuria, hypoalbuminaemia, oedema, hyperlipidaemia).
\item Osteoarthritis is an autoimmune disease characterised by symmetrical small-joint involvement.
\item Benign prostatic hyperplasia (BPH) increases the risk of prostate cancer.
\end{enumerate}
\end{exercice}

\begin{exercice}[Definitions and key terms]
Define the following terms precisely:
\begin{enumerate}
\item \cle{Non-communicable disease (NCD)}
\item \cle{Glomerular filtration rate (GFR)}
\item \cle{Insulin resistance}
\item \cle{Atheroma (atherosclerotic plaque)}
\item \cle{Rheumatoid factor}
\item \cle{Benign prostatic hyperplasia (BPH)}
\item \cle{Goitre}
\item \cle{Myocardial infarction}
\end{enumerate}
\end{exercice}

\begin{exercice}[Direct calculations]
\begin{enumerate}
\item A 45-year-old man has a serum creatinine of 120 $\mu$mol/L. Using the simplified MDRD equation (for adults): 
\[ \text{eGFR (mL/min/1.73 m$^2$)} = 186 \times \left(\frac{\text{creatinine in } \mu\text{mol/L}}{88.4}\right)^{-1.154} \times \text{age}^{-0.203} \times (0.742 \text{ if female}) \]
Calculate his estimated GFR. Is this consistent with chronic kidney disease (CKD) stage 3 (GFR 30--59 mL/min/1.73 m$^2$)?
\item A woman's blood pressure readings on three separate occasions are: 148/92, 152/96, 146/94 mmHg. Classify her hypertension according to WHO/ISH guidelines (Grade 1, 2 or 3).
\item Calculate the Body Mass Index (BMI) of a patient weighing 92 kg and measuring 1.68 m. Classify the nutritional status according to WHO categories.
\item In a glucose tolerance test, a patient's 2-hour plasma glucose is 11.5 mmol/L. Fasting glucose is 6.8 mmol/L. Interpret these results using WHO criteria (fasting $\ge$ 7.0 mmol/L or 2-h $\ge$ 11.1 mmol/L for diabetes; fasting 6.1--6.9 mmol/L for impaired fasting glycaemia; 2-h 7.8--11.0 mmol/L for impaired glucose tolerance).
\end{enumerate}
\end{exercice}

\courssub{I practise}

\begin{exercice}[Kidney disease: proteinuria, haematuria and oedema]
A 12-year-old boy presents with puffy eyes, swollen ankles and cola-coloured urine two weeks after a sore throat. Urinalysis shows +++ protein and numerous red blood cells.
\begin{enumerate}
\item Name the specific part of the nephron most likely damaged.
\item Name the probable disease.
\item Explain the pathophysiological mechanism leading to oedema in this condition.
\item List two other clinical features you would expect to find.
\item Outline the principle of management.
\end{enumerate}
\end{exercice}

\begin{exercice}[Glucose tolerance test interpretation]
The graph below shows plasma glucose concentrations (mmol/L) over 2 hours after a 75 g oral glucose load for three individuals A, B and C.

\begin{popfigure}
\begin{tikzpicture}[scale=0.7]
\begin{axis}[
    width=12cm, height=7cm,
    xlabel={Time (min)}, ylabel={Plasma glucose (mmol/L)},
    xmin=0, xmax=120, ymin=0, ymax=22,
    xtick={0,30,60,90,120}, ytick={0,5,7.8,11.1,15,20},
    grid=major, legend pos=north west,
    legend style={font=\small},
]
\addplot[popblue, thick, mark=*] coordinates {(0,5.2) (30,8.5) (60,7.2) (90,6.5) (120,5.8)};
\addlegendentry{Individual A}
\addplot[poporange, thick, mark=square*] coordinates {(0,7.5) (30,14.2) (60,18.5) (90,16.8) (120,15.2)};
\addlegendentry{Individual B}
\addplot[popgreen, thick, mark=triangle*] coordinates {(0,6.0) (30,12.0) (60,14.5) (90,13.0) (120,11.5)};
\addlegendentry{Individual C}
\end{axis}
\end{tikzpicture}
\legende{Oral glucose tolerance test curves for three individuals.}
\end{popfigure}

\begin{enumerate}
\item Identify which curve corresponds to: (i) a healthy individual, (ii) type 2 diabetes mellitus, (iii) type 1 diabetes mellitus. Justify each choice using WHO diagnostic criteria (fasting and 2-hour values).
\item Explain why the curve for type 1 diabetes typically shows a higher peak and slower return to baseline than type 2.
\item State two other laboratory tests that help distinguish type 1 from type 2 diabetes.
\end{enumerate}
\end{exercice}

\begin{exercice}[Comparison: osteoarthritis vs rheumatoid arthritis]
Complete the table below contrasting osteoarthritis (OA) and rheumatoid arthritis (RA).

\begin{center}
\begin{tabularx}{\linewidth}{|l|X|X|}
\hline
\rowcolor{popblueL}
\textbf{Feature} & \textbf{Osteoarthritis (OA)} & \textbf{Rheumatoid arthritis (RA)} \\
\hline
Primary cause / aetiology & & \\
\hline
Typical age of onset & & \\
\hline
Pattern of joint involvement (symmetry, joints affected) & & \\
\hline
Morning stiffness duration & & \\
\hline
Systemic features (e.g. fever, weight loss, nodules) & & \\
\hline
Radiographic findings (key differences) & & \\
\hline
Autoantibodies & & \\
\hline
First-line pharmacological management & & \\
\hline
\end{tabularx}
\end{center}
\end{exercice}

\begin{exercice}[Case study: endemic goitre]
A 35-year-old woman living in a remote highland village in the Northwest Region of Cameroon presents with a diffuse, non-tender swelling in the anterior neck. She reports fatigue, cold intolerance and weight gain. The village diet consists mainly of cassava and maize; iodised salt is unavailable.
\begin{enumerate}
\item Name the condition and its primary cause.
\item Explain the hormonal feedback failure involved, mentioning TRH, TSH, T3 and T4.
\item Why does the thyroid enlarge (goitre) in this condition?
\item Propose a national prevention measure that has proven effective in Cameroon.
\item List two other consequences of iodine deficiency in pregnancy and early childhood.
\end{enumerate}
\end{exercice}

\courssub{I prepare for the GCE}

\begin{exercice}[{Essay: NCDs as a public health threat in Cameroon] [25 marks}]
\textbf{``Non-infectious diseases are now a greater public-health threat in Cameroon than infectious diseases.''} 
Discuss this statement with reference to at least \textbf{four} body systems. In your answer, include:
\begin{itemize}
\item Epidemiological evidence (trends, risk factors).
\item Major NCDs affecting each system, their pathophysiology and complications.
\item Socio-economic impact on individuals, families and the health system.
\item Prevention strategies at primary, secondary and tertiary levels, with Cameroonian examples.
\end{itemize}
\end{exercice}

\begin{exercice}[{Short answers: Hypertension] [15 marks}]
\begin{enumerate}
\item Define hypertension according to current WHO/ISH guidelines. [2 marks]
\item State \textbf{two} consequences of untreated hypertension on the kidney and \textbf{two} on the brain. [4 marks]
\item Explain why hypertension is called a ``silent killer''. [3 marks]
\item Describe the mechanism of action of \textbf{one} class of antihypertensive drugs (e.g. ACE inhibitors, calcium channel blockers, thiazide diuretics, beta-blockers). [3 marks]
\item A 55-year-old trader has BP 160/100 mmHg on two occasions. Outline the initial investigations you would request and lifestyle advice you would give. [3 marks]
\end{enumerate}
\end{exercice}

\begin{exercice}[{Practical: Blood pressure measurement and interpretation] [15 marks}]
\begin{enumerate}
\item Describe, step by step, the correct procedure for measuring a classmate's blood pressure using a manual sphygmomanometer and stethoscope. Include patient preparation, cuff placement, inflation, deflation and reading. [6 marks]
\item A reading of 150/95 mmHg is obtained. Interpret this value (classification, urgency). [3 marks]
\item List \textbf{four} common errors that can falsely elevate or lower the reading. [2 marks]
\item Advise the person on immediate next steps (lifestyle, follow-up, when to seek medical care). [4 marks]
\end{enumerate}
\end{exercice}

\newpage

\coursec{Solutions to the exercises}

\coursec{Corrections and Worked Solutions — Non-infectious Diseases (BIO06)}

\begin{corrige}[Exercise 1 — Multiple choice questions]
\textbf{1. B — Benign prostatic hyperplasia (BPH).} BPH is a non-infectious, age-related proliferation of prostatic stromal and glandular tissue driven by dihydrotestosterone. Options A, C and D are all infectious: prostatitis caused by \emph{E. coli} is bacterial, gonorrhoea is caused by \emph{Neisseria gonorrhoeae} and syphilis by \emph{Treponema pallidum}.

\textbf{2. C — Type 2 diabetes mellitus.} The triad polyuria–polydipsia–hyperglycaemia indicates diabetes mellitus. The \emph{absence of ketones} points to type 2: residual insulin secretion is sufficient to suppress lipolysis and ketogenesis. Type 1 (B) typically presents with ketoacidosis. Diabetes insipidus (A) causes polyuria and polydipsia but \emph{not} hyperglycaemia. Cushing's syndrome (D) may cause hyperglycaemia but is not the most likely diagnosis here.

\textbf{3. B — Creatinine.} Creatinine is freely filtered at the glomerulus and is neither significantly reabsorbed nor secreted (only minimal secretion), so its clearance approximates GFR. Urea is partly reabsorbed (underestimates GFR); glucose is fully reabsorbed in health; sodium is extensively reabsorbed.

\textbf{4. C — Distal interphalangeal joints.} The DIP joints (Heberden's nodes) are characteristically affected in osteoarthritis but spared in rheumatoid arthritis, which targets MCP, PIP and wrist joints symmetrically.

\textbf{5. C — Tunica intima.} Atherosclerosis begins with endothelial dysfunction and accumulation of LDL-derived lipids in the \cle{tunica intima}, forming fatty streaks then fibrous plaques. The media and adventitia are affected only secondarily.
\end{corrige}

\begin{corrige}[Exercise 2 — True/False with justification]
\begin{enumerate}
\item \textbf{True.} WHO/ISH defines hypertension as sustained systolic BP $\ge 140$ mmHg and/or diastolic BP $\ge 90$ mmHg, confirmed on at least two separate occasions.
\item \textbf{False.} In type 1 diabetes the primary defect is \emph{autoimmune destruction of pancreatic $\beta$-cells} causing absolute insulin deficiency. Insulin resistance is the hallmark of \emph{type 2} diabetes.
\item \textbf{False.} Glomerulonephritis has several presentations: nephritic syndrome (haematuria, mild proteinuria, hypertension, oliguria), nephrotic syndrome, or rapidly progressive renal failure. Nephrotic syndrome occurs only in some forms (e.g. membranous nephropathy, minimal change disease).
\item \textbf{False.} Osteoarthritis is a \emph{degenerative} (``wear-and-tear'') disease of articular cartilage, typically asymmetrical and involving weight-bearing joints and DIP joints. Symmetrical small-joint autoimmune disease describes \emph{rheumatoid arthritis}.
\item \textbf{False.} BPH arises from the \emph{transition zone} of the prostate and is benign; it does \emph{not} transform into or increase the risk of prostate cancer (which usually arises in the peripheral zone). The two conditions share age as a risk factor and may coexist, which explains the confusion.
\end{enumerate}
\end{corrige}

\begin{corrige}[Exercise 3 — Definitions and key terms]
\begin{enumerate}
\item \textbf{Non-communicable disease (NCD):} a chronic disease that is \emph{not} caused by an infectious agent and cannot be transmitted from person to person; it results from genetic, physiological, environmental and behavioural factors (e.g. hypertension, diabetes, cancers, chronic respiratory diseases).
\item \textbf{Glomerular filtration rate (GFR):} the volume of plasma filtered through the glomeruli into Bowman's capsule per unit time (normally about $120\ \mathrm{mL/min/1.73\,m^2}$ in young adults); it is the best overall index of kidney function, estimated clinically from creatinine clearance or equations such as MDRD or CKD-EPI.
\item \textbf{Insulin resistance:} a reduced biological response of target tissues (muscle, liver, adipose tissue) to normal circulating insulin concentrations, so that glucose uptake and utilisation are impaired and compensatory hyperinsulinaemia occurs.
\item \textbf{Atheroma:} a focal, raised lesion of the arterial intima composed of a lipid core (cholesterol esters, foam cells) covered by a fibrous cap of smooth-muscle cells and connective tissue; it narrows the lumen and may rupture, triggering thrombosis.
\item \textbf{Rheumatoid factor (RF):} an autoantibody (usually IgM) directed against the Fc portion of IgG, found in about 70--80\% of patients with rheumatoid arthritis (also in some infections and healthy elderly people, so it is not fully specific).
\item \textbf{Benign prostatic hyperplasia:} a non-malignant, age- and androgen-dependent nodular enlargement of the transition zone of the prostate, causing bladder-outlet obstruction with hesitancy, weak stream, nocturia and urinary retention.
\item \textbf{Goitre:} any enlargement of the thyroid gland; in endemic goitre it results from prolonged TSH stimulation when iodine deficiency limits T3/T4 synthesis.
\item \textbf{Myocardial infarction:} irreversible necrosis of cardiac myocytes caused by prolonged ischaemia, usually from occlusion of a coronary artery by thrombus on a ruptured atherosclerotic plaque; diagnosed by typical chest pain, ECG changes and raised cardiac biomarkers (troponin).
\end{enumerate}
\end{corrige}

\begin{corrige}[Exercise 4 — Direct calculations]
\begin{enumerate}
\item \textbf{eGFR (MDRD).} The patient is male, so the 0.742 factor is omitted. Serum creatinine = $120\ \mu\mathrm{mol/L}$ ($= 120/88.4 = 1.357\ \mathrm{mg/dL}$). Age = 45 years.
\begin{align*}
\text{eGFR} &= 186 \times (\text{Scr})^{-1.154} \times (\text{Age})^{-0.203} \\
&= 186 \times 1.357^{-1.154} \times 45^{-0.203} \\
1.357^{-1.154} &= e^{-1.154 \ln 1.357} = e^{-1.154 \times 0.3056} = e^{-0.3527} \approx 0.703 \\
45^{-0.203} &= e^{-0.203 \ln 45} = e^{-0.203 \times 3.807} = e^{-0.7728} \approx 0.462 \\
\text{eGFR} &= 186 \times 0.703 \times 0.462 \approx 60.4\ \mathrm{mL/min/1.73\,m^2}
\end{align*}
The eGFR $\approx 60\ \mathrm{mL/min/1.73\,m^2}$ lies just \emph{above} the 30--59 band, so it is \textbf{not} consistent with CKD stage 3; it corresponds to stage 2 (60--89) if chronic kidney damage markers are present, and is borderline normal otherwise. (Accept values in the 58--62 range depending on rounding; the conclusion about stage 3 is unchanged if the value falls below 60 — the key skill is correct substitution and comparison with the reference band.)

\item \textbf{Blood pressure classification.} Mean values: systolic $\approx (148+152+146)/3 \approx 149$ mmHg; diastolic $\approx (92+96+94)/3 = 94$ mmHg. All readings fall in the ranges SBP 140--159 and/or DBP 90--99 mmHg $\Rightarrow$ \textbf{Grade 1 (mild) hypertension}. The diagnosis is valid because elevated values were recorded on three separate occasions.

\item \textbf{BMI.}
\[ \text{BMI} = \frac{m}{h^2} = \frac{92}{1.68^2} = \frac{92}{2.8224} \approx 32.6\ \mathrm{kg/m^2} \]
WHO classification: 30.0--34.9 $\Rightarrow$ \textbf{obesity class I (moderate obesity)}.

\item \textbf{Glucose tolerance test.} Fasting glucose 6.8 mmol/L lies in the 6.1--6.9 band $\Rightarrow$ impaired fasting glycaemia; but the 2-hour value 11.5 mmol/L $\ge 11.1$ mmol/L $\Rightarrow$ \textbf{diabetes mellitus}. When categories conflict, the more severe category prevails: the patient has diabetes mellitus (confirmed by a repeat test on another day if asymptomatic).
\end{enumerate}
\end{corrige}

\begin{corrige}[Exercise 5 — Kidney disease: proteinuria, haematuria and oedema]
\begin{enumerate}
\item \textbf{Structure damaged:} the \cle{glomerulus}, specifically the glomerular basement membrane and its filtration barrier (capillary endothelium, basement membrane, podocytes).
\item \textbf{Probable disease:} \textbf{acute post-streptococcal glomerulonephritis} (nephritic syndrome), occurring 1--3 weeks after a group A $\beta$-haemolytic streptococcal pharyngitis.
\item \textbf{Mechanism of oedema:} immune complexes deposit in the glomerular basement membrane $\Rightarrow$ inflammation increases permeability and reduces the filtration surface $\Rightarrow$ GFR falls, so Na$^+$ and water are retained (increased tubular reabsorption, activation of the renin–angiotensin–aldosterone system) $\Rightarrow$ plasma volume expands and hydrostatic pressure rises $\Rightarrow$ fluid shifts into the interstitium. Loss of protein in urine additionally lowers plasma oncotic pressure, worsening the oedema (periorbital in the morning, ankles later).
\item \textbf{Other expected features (any two):} hypertension; oliguria; mild fever/malaise; loin discomfort; elevated serum creatinine; low complement C3.
\item \textbf{Principle of management:} mainly supportive — rest, salt and water restriction, treat hypertension and fluid overload (diuretics), eradicate residual streptococcal infection with penicillin, monitor renal function; most children recover fully.
\end{enumerate}
\end{corrige}

\begin{corrige}[Exercise 6 — Glucose tolerance test interpretation]
\begin{enumerate}
\item \textbf{Identification:}
\begin{itemize}
\item \textbf{Individual A — healthy:} fasting 5.2 mmol/L ($<6.1$) and 2-h 5.8 mmol/L ($<7.8$): both normal; glucose returns to baseline within 2 h.
\item \textbf{Individual B — type 1 diabetes:} fasting 7.5 mmol/L ($\ge 7.0$) and 2-h 15.2 mmol/L ($\ge 11.1$): diabetic by both criteria, with a very high peak (18.5) and persistently marked hyperglycaemia — the pattern of absolute insulin deficiency.
\item \textbf{Individual C — type 2 diabetes:} fasting 6.0 mmol/L (normal/IFG borderline) but 2-h 11.5 mmol/L ($\ge 11.1$): diabetic on the 2-h criterion; the moderate peak and partial decline reflect residual but insufficient insulin secretion with insulin resistance.
\end{itemize}
\item \textbf{Why type 1 peaks higher and falls slower:} in type 1 there is \emph{no} endogenous insulin, so ingested glucose cannot be taken up by muscle and adipose tissue and hepatic glucose output is not suppressed — glucose accumulates to a higher peak and is cleared only slowly (mainly by renal excretion once the tubular threshold $\approx 10$ mmol/L is exceeded, causing glycosuria). In type 2, residual insulin secretion, though delayed and inadequate against insulin resistance, still promotes some peripheral uptake, so the peak is lower and the fall faster.
\item \textbf{Two distinguishing tests:} (i) \emph{autoantibodies} — anti-GAD, anti-islet cell or anti-insulin antibodies (present in type 1); (ii) \emph{C-peptide / fasting insulin} — low or absent in type 1, normal or high in type 2. (Also acceptable: presence of ketones/ketoacidosis, age and BMI at presentation.)
\end{enumerate}
\end{corrige}

\begin{corrige}[Exercise 7 — Comparison: osteoarthritis vs rheumatoid arthritis]
\begin{center}
\begin{tabularx}{\linewidth}{|l|X|X|}
\hline
\rowcolor{popblueL}
\textbf{Feature} & \textbf{Osteoarthritis (OA)} & \textbf{Rheumatoid arthritis (RA)} \\
\hline
Primary cause / aetiology & Degenerative ``wear-and-tear'' loss of articular cartilage; age, obesity, joint overload & Autoimmune chronic synovitis (autoantibodies, immune complexes) destroying cartilage and bone \\
\hline
Typical age of onset & Usually $>50$ years & Any age, peak 30--50 years; female predominance \\
\hline
Pattern of joint involvement & Asymmetrical; weight-bearing joints (hips, knees, spine) and DIP joints & Symmetrical small joints: MCP, PIP, wrists; spares DIP \\
\hline
Morning stiffness duration & Short ($<30$ min) & Prolonged ($>1$ h) \\
\hline
Systemic features & Absent & Present: fever, fatigue, weight loss, rheumatoid nodules, anaemia \\
\hline
Radiographic findings & Joint-space narrowing, osteophytes, subchondral sclerosis and cysts & Juxta-articular osteoporosis, marginal erosions, uniform space narrowing, deformities (ulnar deviation) \\
\hline
Autoantibodies & None & Rheumatoid factor, anti-CCP \\
\hline
First-line pharmacological management & Analgesics (paracetamol), NSAIDs, physiotherapy, weight loss & DMARDs (e.g. methotrexate) early; NSAIDs/corticosteroids for flares \\
\hline
\end{tabularx}
\end{center}
\end{corrige}

\begin{corrige}[Exercise 8 — Case study: endemic goitre]
\begin{enumerate}
\item \textbf{Condition:} endemic (iodine-deficiency) goitre with hypothyroidism. \textbf{Primary cause:} chronic dietary iodine deficiency, aggravated by goitrogens in cassava (cyanogenic glucosides that inhibit iodine uptake).
\item \textbf{Feedback failure:} iodine is required to synthesise the thyroid hormones T3 and T4. Iodine deficiency $\Rightarrow$ low T3/T4 $\Rightarrow$ loss of negative feedback on the hypothalamus and pituitary $\Rightarrow$ TRH and especially TSH secretion rise continuously. Chronically elevated TSH is the driver of the disease.
\item \textbf{Why the gland enlarges:} TSH is trophic: it stimulates thyroid follicular cells to grow and divide (hypertrophy and hyperplasia) and increases colloid accumulation in an attempt to trap more iodine and produce more hormone. Sustained stimulation therefore produces diffuse thyroid enlargement — the goitre.
\item \textbf{National prevention measure:} universal salt iodisation — Cameroon has implemented mandatory iodisation of table salt, which has markedly reduced endemic goitre; supplementation with iodised oil in high-risk areas is complementary.
\item \textbf{Consequences in pregnancy/childhood (any two):} miscarriage and stillbirth; congenital hypothyroidism with cretinism (severe, irreversible mental retardation and stunted growth); impaired cognitive development and reduced school performance; deaf-mutism in severe deficiency.
\end{enumerate}
\end{corrige}

\begin{corrige}[Exercise 9 — Essay: NCDs as a public health threat in Cameroon (25 marks)]
\textbf{Indicative mark scheme and examiner expectations.}

\textbf{Introduction (2 marks):} define NCDs; present the epidemiological transition — while malaria, HIV and tuberculosis remain important, urbanisation, dietary change, tobacco/alcohol use, physical inactivity and ageing are driving a rapid rise of NCDs, creating a \emph{double burden} of disease. A nuanced thesis (NCDs are a growing, comparable threat rather than simply ``greater'') earns full credit.

\textbf{Body — at least four systems (4 $\times$ 3 = 12 marks; 1 mark for the disease + pathophysiology, 1 for complications, 1 for Cameroonian context):}
\begin{itemize}
\item \emph{Cardiovascular:} hypertension (very prevalent, often undiagnosed), stroke, myocardial infarction, heart failure; atherosclerosis mechanism; complications (renal failure, sudden death).
\item \emph{Endocrine/metabolic:} type 2 diabetes rising with obesity; complications — retinopathy, nephropathy, neuropathy, diabetic foot amputations.
\item \emph{Renal:} chronic kidney disease from hypertension, diabetes and glomerulonephritis; dialysis is costly and scarce.
\item \emph{Musculoskeletal:} osteoarthritis in an ageing population, road-traffic injuries; disability and loss of productivity.
\item \emph{Reproductive/other:} cancers (cervix, breast, prostate) — late presentation, limited radiotherapy.
\end{itemize}

\textbf{Socio-economic impact (4 marks):} catastrophic out-of-pocket spending, loss of breadwinners, caregiving burden, absenteeism, pressure on an under-resourced health system, treatment costs (insulin, dialysis, antihypertensives) pushing families into poverty.

\textbf{Prevention (5 marks):}
\begin{itemize}
\item \emph{Primary:} health education, salt reduction, tobacco and alcohol control, promotion of physical activity, iodised salt, vaccination against hepatitis B and HPV (cancer prevention).
\item \emph{Secondary:} screening — BP and blood glucose campaigns in churches, markets and workplaces; cervical cancer screening; early detection programmes.
\item \emph{Tertiary:} affordable essential medicines, management of complications, rehabilitation, physiotherapy, palliative care.
\end{itemize}

\textbf{Conclusion (2 marks):} balanced judgement — NCDs now rival infectious diseases and require an integrated national response.

\textbf{Examiner expectations:} accurate terminology, correct pathophysiology, genuine Cameroonian examples (not invented statistics), logical structure, and coverage of all four requested elements. Vague lists without explanation score poorly.
\end{corrige}

\begin{corrige}[Exercise 10 — Short answers: Hypertension (15 marks)]
\begin{enumerate}
\item \textbf{Definition [2]:} a sustained elevation of arterial blood pressure $\ge 140$ mmHg systolic and/or $\ge 90$ mmHg diastolic, confirmed on at least two separate measurements/occasions (WHO/ISH). [1 mark for values, 1 for ``sustained/confirmed''.]
\item \textbf{Consequences [4]:} \emph{Kidney (any 2):} nephrosclerosis (hyaline arteriolosclerosis), chronic kidney disease/proteinuria, end-stage renal failure. \emph{Brain (any 2):} ischaemic or haemorrhagic stroke, hypertensive encephalopathy, vascular dementia, transient ischaemic attacks.
\item \textbf{``Silent killer'' [3]:} hypertension usually causes \emph{no symptoms} for years while progressively damaging arteries, heart, brain, kidneys and retina; the first manifestation may be a fatal event (stroke, myocardial infarction, sudden death); hence many Cameroonians are unaware of their condition until complications occur — only routine screening reveals it.
\item \textbf{Mechanism of one drug class [3] — e.g. ACE inhibitors:} they inhibit angiotensin-converting enzyme $\Rightarrow$ less angiotensin II $\Rightarrow$ reduced vasoconstriction and reduced aldosterone secretion $\Rightarrow$ decreased Na$^+$/water retention and vasodilation $\Rightarrow$ lower blood pressure; they also reduce glomerular pressure, protecting the kidney. (Accept: thiazides — inhibit NaCl reabsorption in the distal convoluted tubule causing diuresis and reduced plasma volume; beta-blockers — block $\beta_1$ receptors reducing heart rate and contractility; calcium channel blockers — block L-type Ca$^{2+}$ channels in vascular smooth muscle causing vasodilation.)
\item \textbf{Investigations and advice [3]:} \emph{Investigations:} confirm BP, urinalysis (proteinuria), serum creatinine/electrolytes, fasting glucose and lipids, ECG, fundoscopy. \emph{Lifestyle:} reduce salt ($<5$ g/day), reduce alcohol, stop tobacco, weight reduction, regular physical activity, fruit/vegetable-rich diet, stress management, adherence to follow-up.
\end{enumerate}
\end{corrige}

\begin{corrige}[Exercise 11 — Practical: Blood pressure measurement (15 marks)]
\begin{enumerate}
\item \textbf{Procedure [6 marks — 1 per key step]:}
\begin{enumerate}
\item \emph{Preparation:} the subject rests seated for 5 minutes, back supported, feet flat, no caffeine, tobacco or exercise in the previous 30 min; arm bare and supported at heart level.
\item \emph{Cuff placement:} choose the correct cuff size (bladder length encircling $\approx 80\%$ of arm circumference); wrap snugly 2--3 cm above the elbow crease, with the artery marker over the brachial artery.
\item \emph{Locate the pulse:} palpate the brachial artery in the antecubital fossa; place the stethoscope diaphragm over it (not under the cuff).
\item \emph{Inflation:} close the valve and inflate rapidly to about 20--30 mmHg above the point where the radial pulse disappears (avoiding auscultatory gap errors).
\item \emph{Deflation:} open the valve and deflate slowly ($\approx 2$--3 mmHg/s). The first clear tapping sound (Korotkoff I) = \textbf{systolic} pressure; the point where sounds disappear (Korotkoff V) = \textbf{diastolic} pressure.
\item \emph{Reading and recording:} read to the nearest 2 mmHg, record arm, position and time; repeat after 1--2 min and average; measure on both arms at first assessment.
\end{enumerate}
\item \textbf{Interpretation of 150/95 mmHg [3]:} both values fall in \textbf{Grade 1 (mild) hypertension} (SBP 140--159, DBP 90--99). It is \emph{not} an emergency; it requires confirmation on at least one or two further occasions (or ambulatory/home monitoring) before diagnosis and treatment decisions.
\item \textbf{Four common errors [2]:} cuff too small or too loose (falsely high); arm below heart level or unsupported (falsely high); deflation too fast (inaccurate readings); talking or recent caffeine/tobacco/exercise (falsely high); cuff over clothing; stethoscope pressed too hard (falsely low diastolic). [Any four, $\frac{1}{2}$ mark each.]
\item \textbf{Next steps [4]:} repeat measurements on separate days to confirm; begin lifestyle measures immediately (salt reduction, weight control, exercise, limit alcohol, stop smoking); check for target-organ damage (urinalysis, creatinine, glucose, ECG); seek medical care promptly for confirmed persistent hypertension, and \emph{urgently} if BP $\ge 180/110$ mmHg or symptoms such as severe headache, chest pain, visual disturbance or neurological signs occur.
\end{enumerate}
\end{corrige}

\newpage

\coursec{Summary sheet}

\begin{popfigure}
\begin{tikzpicture}[grow cyclic, text width=2.6cm, align=flush center, level 1/.style={level distance=3.6cm, sibling angle=51.4}, every node/.style={font=\small}]
\node[circle, draw=popdark, fill=popblueL, text width=2.8cm, align=center, font=\small\bfseries] {Non-infectious diseases (NCDs)}
child[concept color=popblue] { node[draw=popblue, fill=popblueL, rounded corners] {Reproductive system disorders}
  child { node[draw=popblue, fill=white, rounded corners, font=\scriptsize] {Male: prostate enlargement, prostate \& testicular cancer, infertility} }
  child { node[draw=popblue, fill=white, rounded corners, font=\scriptsize] {Female: fibroids, ovarian cysts, cervical \& breast cancer, infertility} } }
child[concept color=popgreen] { node[draw=popgreen, fill=popgreenL, rounded corners] {Nervous system disorders}
  child { node[draw=popgreen, fill=white, rounded corners, font=\scriptsize] {Stroke, epilepsy, Parkinson, Alzheimer} } }
child[concept color=poporange] { node[draw=poporange, fill=white, rounded corners] {Endocrine diseases}
  child { node[draw=poporange, fill=white, rounded corners, font=\scriptsize] {Diabetes mellitus (types 1 \& 2), goitre, hypo-/hyperthyroidism} } }
child[concept color=popgold] { node[draw=popgold, fill=white, rounded corners] {Renal diseases}
  child { node[draw=popgold, fill=white, rounded corners, font=\scriptsize] {Nephritis, kidney stones, renal failure, dialysis} } }
child[concept color=poppurple] { node[draw=poppurple, fill=white, rounded corners] {Cardiovascular diseases}
  child { node[draw=poppurple, fill=white, rounded corners, font=\scriptsize] {Hypertension, atherosclerosis, heart attack, heart failure} } }
child[concept color=poppink] { node[draw=poppink, fill=white, rounded corners] {Musculoskeletal disorders}
  child { node[draw=poppink, fill=white, rounded corners, font=\scriptsize] {Arthritis, osteoporosis, rickets, muscular dystrophy, fractures} } }
child[concept color=popmuted] { node[draw=popmuted, fill=white, rounded corners] {Prevention of NCDs}
  child { node[draw=popmuted, fill=white, rounded corners, font=\scriptsize] {Balanced diet, exercise, no smoking/alcohol, screening} } };
\end{tikzpicture}
\legende{Mind map of chapter BIO06: the main families of non-infectious diseases, their examples and prevention.}
\end{popfigure}

\begin{aretenir}[Key definitions]
\begin{itemize}
\item \cle{Non-infectious disease} (NCD): a disease \textbf{not caused by a pathogen} and \textbf{not transmissible} from person to person; it results from genetic factors, lifestyle, ageing, hormonal imbalance or environmental factors.
\item \cle{Risk factor}: any behaviour or condition (smoking, harmful alcohol use, obesity, high salt intake, physical inactivity, heredity) that increases the probability of developing an NCD.
\item \cle{Benign tumour}: a localised, slow-growing cell mass that remains encapsulated and does not spread; \cle{malignant tumour} (cancer): invades neighbouring tissues and forms \cle{metastases} (secondary tumours) via the blood or lymph.
\item \cle{Infertility}: the inability of a couple to achieve pregnancy after about one year of regular unprotected intercourse.
\item \cle{Hypertension}: persistently raised arterial blood pressure, conventionally $\geq 140/90\ \mathrm{mmHg}$ (systolic/diastolic) on repeated measurements.
\item \cle{Atherosclerosis}: deposition of fatty plaques (rich in cholesterol) in the inner wall of arteries, narrowing the lumen and reducing blood flow.
\item \cle{Renal failure}: inability of the kidneys to filter the blood and excrete nitrogenous wastes (urea); treated by \cle{dialysis} or kidney transplant.
\end{itemize}
\end{aretenir}

\begin{aretenir}[Essential facts by system]
\begin{itemize}
\item \emph{Reproductive:} male — benign prostatic hyperplasia (enlarged prostate, difficulty urinating), prostate and testicular cancers, low sperm count; female — uterine fibroids, ovarian cysts, cervical cancer (associated with human papillomavirus, HPV; preventable by vaccination and screening), breast cancer (detected early by self-examination and mammography).
\item \emph{Nervous:} \cle{stroke} (interrupted blood supply to part of the brain), \cle{epilepsy} (recurrent seizures from abnormal electrical discharges in the brain), \cle{Parkinson disease} (degeneration of dopamine-producing neurones; tremor, rigidity, slow movement), \cle{Alzheimer disease} (progressive memory loss and dementia).
\item \emph{Endocrine:} \cle{diabetes mellitus} — chronic hyperglycaemia; type 1 (autoimmune destruction of the $\beta$ cells of the islets of Langerhans, no insulin, insulin-dependent, usually juvenile onset), type 2 (insulin resistance, associated with obesity and lifestyle, usually adult onset). \cle{Goitre} and hypothyroidism from iodine deficiency; hyperthyroidism (excess thyroxine: weight loss, rapid heartbeat, exophthalmia).
\item \emph{Renal:} nephritis (inflammation of the nephrons), kidney stones (mineral crystals in the kidney or ureter), chronic renal failure; normal urine contains \textbf{no glucose and no protein}.
\item \emph{Cardiovascular:} hypertension, atherosclerosis, coronary thrombosis $\Rightarrow$ \cle{myocardial infarction} (heart attack: death of heart muscle), heart failure, stroke.
\item \emph{Musculoskeletal:} \cle{osteoarthritis} (wear of articular cartilage), \cle{rheumatoid arthritis} (autoimmune inflammation of joints), \cle{osteoporosis} (loss of bone mineral density; calcium/vitamin D deficiency; common in post-menopausal women), \cle{rickets} (vitamin D deficiency in children; bowed legs), muscular dystrophy (inherited progressive muscle wasting), sprains and fractures.
\end{itemize}
\end{aretenir}

\begin{methode}[Methods to master]
\begin{enumerate}
\item Interpret urine/blood test results: glucose in urine $\Rightarrow$ suspect diabetes mellitus; protein in urine $\Rightarrow$ suspect kidney damage; high blood urea $\Rightarrow$ suspect renal failure.
\item Read a blood pressure value: systolic over diastolic (e.g. $120/80\ \mathrm{mmHg}$); classify as normal, elevated or hypertensive.
\item Compare type 1 and type 2 diabetes in a table (cause, age of onset, treatment).
\item Draw a labelled flow diagram: risk factor $\rightarrow$ lesion $\rightarrow$ disease $\rightarrow$ complication $\rightarrow$ prevention.
\end{enumerate}
\end{methode}

\begin{attention}[Common mistakes]
\begin{itemize}
\item Confusing \textbf{infectious} and \textbf{non-infectious} diseases: NCDs are never contagious.
\item Saying diabetes is caused directly by ``eating sugar'': it is a failure of insulin secretion (type 1) or of insulin action (type 2).
\item Confusing \cle{osteoporosis} (porous, fragile bones) with \cle{osteoarthritis} (worn joint cartilage).
\item Confusing \cle{rickets} (children, vitamin D deficiency) with \cle{osteomalacia} (the equivalent condition in adults).
\item Believing hypertension always produces symptoms: it is a ``silent killer'' often detected only by measurement.
\item Writing that all tumours are cancers: benign tumours do not invade tissues and do not metastasise.
\item Listing meningitis as a nervous system NCD: meningitis is an \textbf{infectious} disease; only its after-effects (sequelae) are non-infectious.
\end{itemize}
\end{attention}

\begin{aretenir}[GCE examination tips]
\begin{itemize}
\item Structured questions often ask: ``State two causes, two symptoms and two preventive measures'' — prepare such triplets for each disease.
\item Use the Cameroonian context wisely: cervical and prostate cancer screening campaigns, iodised salt against goitre, hypertension and diabetes as rising NCDs, sickle cell disease as a genetic (non-infectious) disorder.
\item Always link prevention to the \textbf{four shared behavioural risk factors}: tobacco use, harmful alcohol use, unhealthy diet and physical inactivity.
\item In data questions, quote figures from the table or graph to justify your diagnosis (e.g. a fasting blood glucose value, a blood pressure reading).
\end{itemize}
\end{aretenir}

\findecours

\end{document}
